\documentclass[11pt]{article}
\usepackage{ifthen}
\usepackage[utf8]{inputenc}
\usepackage[T1]{fontenc}
\usepackage[english]{babel}
\usepackage{geometry}
\geometry{a4paper,margin=2cm}
\usepackage{longtable,graphicx,multirow,enumitem,tabularx,setspace,ragged2e}
\setlist{nosep,leftmargin=*}
\renewcommand\arraystretch{1.15}
\pagestyle{empty}
\begin{document}
\begin{center}
\begin{tabular}{lccr}
 \multirow{3}{*}{\includegraphics[height=2.7cm]{brasao.png}} &
 \multicolumn{2}{c}{\bfseries FEDERAL UNIVERSITY OF OURO PRETO} &
 \ \ \ \ \multirow{3}{*}{\includegraphics[height=2.7cm]{ufop.png}} \\
 & \multicolumn{2}{c}{\bfseries DEAN'S OFFICE FOR UNDERGRADUATE STUDIES} & \\
 & \multicolumn{2}{c}{} & \\
 & \multicolumn{2}{c}{\Large\bfseries COURSE PROGRAM} & \\
\end{tabular}
\end{center}

\begin{center}
\begin{longtable}{|p{4cm}|p{4cm}|p{4cm}|p{4cm}|}
\hline
\multicolumn{3}{|p{12cm}|}{Course Name in English:} &
\multicolumn{1}{p{4cm}|}{Code:} \\ 
\multicolumn{3}{|p{12cm}|}{\textbf{Dissemination of Scientific Knowledge}} &
\textbf{BCC292}\\ 
\multicolumn{3}{|p{12cm}|}{Course Name in Portuguese:} & \\ 
\multicolumn{3}{|p{12cm}|}{\textbf{Divulgação Científica}} & \\ 
\hline
\multicolumn{3}{|p{12cm}|}{Department Name and Acronym:} & Academic Unit: \\ 
\multicolumn{3}{|p{12cm}|}{Department of Computer Science (DECOM)} & {ICEB} \\ 
\hline
\multicolumn{4}{|p{16cm}|}{Delivery Mode:
[X] in person \hspace{0.6cm}
[ ] blended \hspace{0.6cm}
[ ] distance learning}\\
\hline
\multicolumn{2}{|p{8cm}|}{Semester workload} &
\multicolumn{2}{p{8cm}|}{Weekly workload}\\
\hline
\multicolumn{1}{|p{4cm}|}{Total} &
\multicolumn{1}{p{4cm}|}{Extension} &
\multicolumn{1}{p{4cm}|}{Theoretical} &
\multicolumn{1}{p{4cm}|}{Practical} \\ 
\multicolumn{1}{|p{4cm}|}{120\,hours} &
\multicolumn{1}{p{4cm}|}{60\;hours} &
\multicolumn{1}{p{4cm}|}{0\;class hours} &
\multicolumn{1}{p{4cm}|}{8\;class hours} \\ 
\hline
\multicolumn{4}{|p{16cm}|}{Course Description:}\\
\multicolumn{4}{|p{16cm}|}{}\\
\multicolumn{4}{|p{\dimexpr 16cm + 6\tabcolsep\relax}|}{Science communication; concepts and practical aspects of social media; scientific reading; concepts of journalistic text production.}\\
\multicolumn{4}{|p{16cm}|}{}\\
\hline
\multicolumn{4}{|p{16cm}|}{Course Content:}\\*
\multicolumn{4}{|p{\dimexpr 16cm + 6\tabcolsep\relax}|}{\begin{enumerate}[start=1]\item Introduction to Science Communication and its Social Relevance: the role of science in society and the importance of scientific communication; science communication as a tool for inclusion and access to knowledge; the impact of social media and digital media on the dissemination of science\end{enumerate}}\\
\multicolumn{4}{|p{\dimexpr 16cm + 6\tabcolsep\relax}|}{\begin{enumerate}[start=2]\item Communication Strategies and Dialogicity in Science Communication: techniques for making scientific content accessible to different audiences; interaction with the community to identify relevant topics; language and content adaptation for different media (video, text, audio)\end{enumerate}}\\
\multicolumn{4}{|p{\dimexpr 16cm + 6\tabcolsep\relax}|}{\begin{enumerate}[start=3]\item Co-creation and Production of Scientific Content for the Community: science communication workshops in schools and community spaces; production of teaching materials on technology and computing for diverse audiences; planning and execution of educational campaigns on topics such as digital privacy, artificial intelligence, and the social impact of computing\end{enumerate}}\\
\multicolumn{4}{|p{\dimexpr 16cm + 6\tabcolsep\relax}|}{\begin{enumerate}[start=4]\item Use of Digital Technologies for Science Communication: tools for producing videos, podcasts, and infographics; strategies for boosting content on social media; monitoring and analysis of the impact of outreach initiatives\end{enumerate}}\\
\multicolumn{4}{|p{\dimexpr 16cm + 6\tabcolsep\relax}|}{\begin{enumerate}[start=5]\item Impact Assessment and Critical Reflection on Science Communication: collection of feedback from the community served; reflection on transforming academic knowledge into accessible information; records of good practices and systematization of lessons learned\end{enumerate}}\\
\multicolumn{4}{|p{\dimexpr 16cm + 6\tabcolsep\relax}|}{\begin{enumerate}[start=6]\item \textbf{Community Profile:} The course BCC292 - Science Communication aims to bring the university closer to society through accessible communication of scientific and technological knowledge. Its target audience includes schools, non-governmental organizations (NGOs), cultural collectives, entrepreneurs, teachers, and communities with limited access to scientific information. Students will act as knowledge mediators, turning computer science concepts into content that is understandable and applicable to the daily life of the population. Interaction with the community will be participatory and dialogic, ensuring that the topics addressed in the course are defined based on the interests and needs of society. In this way, the course contributes to the democratization of knowledge and to the strengthening of digital and scientific inclusion.\end{enumerate}}\\
\multicolumn{4}{|p{\dimexpr 16cm + 6\tabcolsep\relax}|}{\begin{enumerate}[start=7]\item \textbf{Extension Objectives:} The course aims to (i) promote the democratization of scientific and technological knowledge, making computing accessible to diverse audiences through science communication strategies, (ii) foster dialogic interaction between the university and the community, allowing the topics addressed in the course to be built collectively based on social demands, (iii) train students to communicate scientific concepts in an accessible way, developing writing and speaking skills and the use of technologies to popularize science, (iv) encourage the production of educational content by promoting workshops, digital materials, and informational campaigns on relevant computing topics, and (v) assess the impact of science communication on the community, ensuring that the actions carried out have a concrete effect on the dissemination of knowledge and on social engagement.\end{enumerate}}\\
\multicolumn{4}{|p{\dimexpr 16cm + 6\tabcolsep\relax}|}{}\\
\hline
\multicolumn{4}{|p{16cm}|}{Basic Bibliography:}\\*
\multicolumn{4}{|p{\dimexpr 16cm + 6\tabcolsep\relax}|}{\begin{itemize}\item BOWATER, Laura; YEOMAN, Kay. Science Communication: A Practical Guide for Scientists. Chichester: Wiley-Blackwell, 2012.\end{itemize}}\\
\multicolumn{4}{|p{\dimexpr 16cm + 6\tabcolsep\relax}|}{\begin{itemize}\item TUTEN, Tracy L. Social Media Marketing. 5. ed. London: SAGE Publications, 2023.\end{itemize}}\\
\multicolumn{4}{|p{\dimexpr 16cm + 6\tabcolsep\relax}|}{\begin{itemize}\item CHAFFEY, Dave; ELLIS-CHADWICK, Fiona. Digital Marketing: Strategy, Implementation and Practice. 8. ed. Harlow: Pearson, 2022.\end{itemize}}\\
\multicolumn{4}{|p{\dimexpr 16cm + 6\tabcolsep\relax}|}{\begin{itemize}\item ALLEY, Michael. The Craft of Scientific Writing. 4. ed. New York: Springer, 2018.\end{itemize}}\\
\multicolumn{4}{|p{\dimexpr 16cm + 6\tabcolsep\relax}|}{\begin{itemize}\item BLUM, Deborah; KNUDSON, Mary; HENIG, Robin Marantz (ed.). A Field Guide for Science Writers: The Official Guide of the National Association of Science Writers. 2. ed. New York: Oxford University Press, 2006.\end{itemize}}\\
\multicolumn{4}{|p{16cm}|}{}\\
\hline
\multicolumn{4}{|p{16cm}|}{Complementary Bibliography:}\\*
\multicolumn{4}{|p{\dimexpr 16cm + 6\tabcolsep\relax}|}{\begin{itemize}\item BARON, Nancy. Escape from the Ivory Tower: A Guide to Making Your Science Matter. Washington, DC: Island Press, 2010.\end{itemize}}\\
\multicolumn{4}{|p{\dimexpr 16cm + 6\tabcolsep\relax}|}{\begin{itemize}\item SCHIMEL, Joshua. Writing Science: How to Write Papers That Get Cited and Proposals That Get Funded. New York: Oxford University Press, 2012.\end{itemize}}\\
\multicolumn{4}{|p{\dimexpr 16cm + 6\tabcolsep\relax}|}{\begin{itemize}\item GRAFF, Gerald; BIRKENSTEIN, Cathy. “They Say / I Say”: The Moves That Matter in Academic Writing. 5. ed. New York: W. W. Norton, 2021.\end{itemize}}\\
\hline
\end{longtable}
\end{center}

\clearpage
\begin{center}
\begin{tabular}{lccr}
 \multirow{3}{*}{\includegraphics[height=2.7cm]{brasao.png}} &
 \multicolumn{2}{c}{\bfseries FEDERAL UNIVERSITY OF OURO PRETO} &
 \ \ \ \ \multirow{3}{*}{\includegraphics[height=2.7cm]{ufop.png}} \\
 & \multicolumn{2}{c}{\bfseries DEAN'S OFFICE FOR UNDERGRADUATE STUDIES} & \\
 & \multicolumn{2}{c}{} & \\
 & \multicolumn{2}{c}{\Large\bfseries COURSE PROGRAM} & \\
\end{tabular}
\end{center}

\begin{center}
\begin{longtable}{|p{4cm}|p{4cm}|p{4cm}|p{4cm}|}
\hline
\multicolumn{3}{|p{12cm}|}{Course Name in English:} &
\multicolumn{1}{p{4cm}|}{Code:} \\ 
\multicolumn{3}{|p{12cm}|}{\textbf{Functional Programming}} &
\textbf{BCC191}\\ 
\multicolumn{3}{|p{12cm}|}{Course Name in Portuguese:} & \\ 
\multicolumn{3}{|p{12cm}|}{\textbf{Programação Funcional}} & \\ 
\hline
\multicolumn{3}{|p{12cm}|}{Department Name and Acronym:} & Academic Unit: \\ 
\multicolumn{3}{|p{12cm}|}{Department of Computer Science (DECOM)} & {ICEB} \\ 
\hline
\multicolumn{4}{|p{16cm}|}{Delivery Mode:
[X] in person \hspace{0.6cm}
[ ] blended \hspace{0.6cm}
[ ] distance learning}\\
\hline
\multicolumn{2}{|p{8cm}|}{Semester workload} &
\multicolumn{2}{p{8cm}|}{Weekly workload}\\
\hline
\multicolumn{1}{|p{4cm}|}{Total} &
\multicolumn{1}{p{4cm}|}{Extension} &
\multicolumn{1}{p{4cm}|}{Theoretical} &
\multicolumn{1}{p{4cm}|}{Practical} \\ 
\multicolumn{1}{|p{4cm}|}{60\,hours} &
\multicolumn{1}{p{4cm}|}{0\;hours} &
\multicolumn{1}{p{4cm}|}{2\;class hours} &
\multicolumn{1}{p{4cm}|}{2\;class hours} \\ 
\hline
\multicolumn{4}{|p{16cm}|}{Course Description:}\\
\multicolumn{4}{|p{16cm}|}{}\\
\multicolumn{4}{|p{\dimexpr 16cm + 6\tabcolsep\relax}|}{Characteristics of the main programming paradigms; principles of the functional programming paradigm; main characteristics of functional programming languages: recursion, functional abstraction, higher-order functions, algebraic data types, polymorphism, type inference, strict and lazy evaluation, overloading; demonstrating the correctness of functional programs.}\\
\multicolumn{4}{|p{16cm}|}{}\\
\hline
\multicolumn{4}{|p{16cm}|}{Course Content:}\\*
\multicolumn{4}{|p{\dimexpr 16cm + 6\tabcolsep\relax}|}{\begin{enumerate}[start=1]\item Introduction\end{enumerate}}\\
\multicolumn{4}{|p{\dimexpr 16cm + 6\tabcolsep\relax}|}{\begin{enumerate}[start=2]\item Programming paradigms\end{enumerate}}\\
\multicolumn{4}{|p{\dimexpr 16cm + 6\tabcolsep\relax}|}{\begin{enumerate}[start=3]\item Defining functions\end{enumerate}}\\
\multicolumn{4}{|p{\dimexpr 16cm + 6\tabcolsep\relax}|}{\begin{enumerate}[start=4]\item Recursive functions\end{enumerate}}\\
\multicolumn{4}{|p{\dimexpr 16cm + 6\tabcolsep\relax}|}{\begin{enumerate}[start=5]\item Tuples, lists, and parametric and overloading polymorphism\end{enumerate}}\\
\multicolumn{4}{|p{\dimexpr 16cm + 6\tabcolsep\relax}|}{\begin{enumerate}[start=6]\item Pattern matching\end{enumerate}}\\
\multicolumn{4}{|p{\dimexpr 16cm + 6\tabcolsep\relax}|}{\begin{enumerate}[start=7]\item Input and output\end{enumerate}}\\
\multicolumn{4}{|p{\dimexpr 16cm + 6\tabcolsep\relax}|}{\begin{enumerate}[start=8]\item Anonymous and higher-order functions\end{enumerate}}\\
\multicolumn{4}{|p{\dimexpr 16cm + 6\tabcolsep\relax}|}{\begin{enumerate}[start=9]\item Algebraic data types\end{enumerate}}\\
\multicolumn{4}{|p{\dimexpr 16cm + 6\tabcolsep\relax}|}{\begin{enumerate}[start=10]\item Functors and Monads\end{enumerate}}\\
\multicolumn{4}{|p{\dimexpr 16cm + 6\tabcolsep\relax}|}{\begin{enumerate}[start=11]\item Property-based testing\end{enumerate}}\\
\multicolumn{4}{|p{\dimexpr 16cm + 6\tabcolsep\relax}|}{\begin{enumerate}[start=12]\item Reasoning about functional programs\end{enumerate}}\\
\multicolumn{4}{|p{\dimexpr 16cm + 6\tabcolsep\relax}|}{\begin{enumerate}[start=13]\item Notions of dependent types in programming languages\end{enumerate}}\\
\multicolumn{4}{|p{\dimexpr 16cm + 6\tabcolsep\relax}|}{}\\
\hline
\multicolumn{4}{|p{16cm}|}{Basic Bibliography:}\\*
\multicolumn{4}{|p{\dimexpr 16cm + 6\tabcolsep\relax}|}{\begin{itemize}\item Lipovaca, Miran. Learn you a Haskell for the great good! No Starch Press. Disponível gratuitamente on-line: http://learnyouahaskell.com\end{itemize}}\\
\multicolumn{4}{|p{\dimexpr 16cm + 6\tabcolsep\relax}|}{\begin{itemize}\item O'Sullivan, Bryan; Stewart, Don; Goerzen, John. Real World Haskell. Disponível gratuitamente on-line: http://book.realworldhaskell.org\end{itemize}}\\
\multicolumn{4}{|p{\dimexpr 16cm + 6\tabcolsep\relax}|}{\begin{itemize}\item WADLER, Philip; KOKKE, Wen; SIEK, Jeremy. Programming Languages Foundations in Agda. Disponível gratuitamente on-line: https://plfa.github.io\end{itemize}}\\
\multicolumn{4}{|p{\dimexpr 16cm + 6\tabcolsep\relax}|}{\begin{itemize}\item MARLOW, Simon. Parallel and Concurrent Programming in Haskell. Disponível gratuitamente on-line: https://simonmar.github.io/pages/pcph.html\end{itemize}}\\
\multicolumn{4}{|p{16cm}|}{}\\
\hline
\multicolumn{4}{|p{16cm}|}{Complementary Bibliography:}\\*
\multicolumn{4}{|p{\dimexpr 16cm + 6\tabcolsep\relax}|}{\begin{itemize}\item Jhala, Ranjit; Seidel, Eric; Vazou, Niki. Programming with Refinement Types: An introduction to Liquid Haskell. Disponível gratuitamente on-line: http://ucsd-progsys.github.io/liquidhaskell-tutorial/\end{itemize}}\\
\multicolumn{4}{|p{\dimexpr 16cm + 6\tabcolsep\relax}|}{\begin{itemize}\item BIRD, Richard. Introduction to Functional Programming using Haskell. 2. ed. London: Prentice Hall Press, 1998.\end{itemize}}\\
\multicolumn{4}{|p{\dimexpr 16cm + 6\tabcolsep\relax}|}{\begin{itemize}\item MILNER, R.; TOFTE, Mads; HARPER, Robert. The definition of standard ML. Cambridge: MIT, 1997.\end{itemize}}\\
\multicolumn{4}{|p{\dimexpr 16cm + 6\tabcolsep\relax}|}{\begin{itemize}\item ULLMAN, Jeffrey D. Elements of ML programming. New Jersey: Prentice-Hall, 1998.\end{itemize}}\\
\multicolumn{4}{|p{\dimexpr 16cm + 6\tabcolsep\relax}|}{\begin{itemize}\item PAULSON, Laurence C. ML for the working programming. 2. ed. Cambridge: Cambridge University, 1996.\end{itemize}}\\
\multicolumn{4}{|p{\dimexpr 16cm + 6\tabcolsep\relax}|}{\begin{itemize}\item PEYTON JONES, Simon L. The implementation of functional programming languages. New York: Prentice-Hall, 1987.\end{itemize}}\\
\hline
\end{longtable}
\end{center}

\clearpage
\begin{center}
\begin{tabular}{lccr}
 \multirow{3}{*}{\includegraphics[height=2.7cm]{brasao.png}} &
 \multicolumn{2}{c}{\bfseries FEDERAL UNIVERSITY OF OURO PRETO} &
 \ \ \ \ \multirow{3}{*}{\includegraphics[height=2.7cm]{ufop.png}} \\
 & \multicolumn{2}{c}{\bfseries DEAN'S OFFICE FOR UNDERGRADUATE STUDIES} & \\
 & \multicolumn{2}{c}{} & \\
 & \multicolumn{2}{c}{\Large\bfseries COURSE PROGRAM} & \\
\end{tabular}
\end{center}

\begin{center}
\begin{longtable}{|p{4cm}|p{4cm}|p{4cm}|p{4cm}|}
\hline
\multicolumn{3}{|p{12cm}|}{Course Name in English:} &
\multicolumn{1}{p{4cm}|}{Code:} \\ 
\multicolumn{3}{|p{12cm}|}{\textbf{Software Engineering II}} &
\textbf{BCC194}\\ 
\multicolumn{3}{|p{12cm}|}{Course Name in Portuguese:} & \\ 
\multicolumn{3}{|p{12cm}|}{\textbf{Engenharia de Software II}} & \\ 
\hline
\multicolumn{3}{|p{12cm}|}{Department Name and Acronym:} & Academic Unit: \\ 
\multicolumn{3}{|p{12cm}|}{Department of Computer Science (DECOM)} & {ICEB} \\ 
\hline
\multicolumn{4}{|p{16cm}|}{Delivery Mode:
[X] in person \hspace{0.6cm}
[ ] blended \hspace{0.6cm}
[ ] distance learning}\\
\hline
\multicolumn{2}{|p{8cm}|}{Semester workload} &
\multicolumn{2}{p{8cm}|}{Weekly workload}\\
\hline
\multicolumn{1}{|p{4cm}|}{Total} &
\multicolumn{1}{p{4cm}|}{Extension} &
\multicolumn{1}{p{4cm}|}{Theoretical} &
\multicolumn{1}{p{4cm}|}{Practical} \\ 
\multicolumn{1}{|p{4cm}|}{60\,hours} &
\multicolumn{1}{p{4cm}|}{0\;hours} &
\multicolumn{1}{p{4cm}|}{4\;class hours} &
\multicolumn{1}{p{4cm}|}{0\;class hours} \\ 
\hline
\multicolumn{4}{|p{16cm}|}{Course Description:}\\
\multicolumn{4}{|p{16cm}|}{}\\
\multicolumn{4}{|p{\dimexpr 16cm + 6\tabcolsep\relax}|}{Software Development Processes. Software Requirements Engineering. Software Project Management. Software as a Service (SaaS). DevOps.}\\
\multicolumn{4}{|p{16cm}|}{}\\
\hline
\multicolumn{4}{|p{16cm}|}{Course Content:}\\*
\multicolumn{4}{|p{\dimexpr 16cm + 6\tabcolsep\relax}|}{\begin{enumerate}[start=1]\item Development Processes: importance of processes for team productivity and product quality assurance; types of processes (sequential or cyclical, incremental and/or evolutionary); plan-driven processes (waterfall, spiral, Rational Unified Process); agile processes (Agile Manifesto, main characteristics, main differences between XP, Scrum and Kanban, criticisms of agile methods); Extreme Programming (historical context, principles and values, customer representative and user stories, iterations and releases, effort estimation, programming practices, project management); Scrum (historical context, development cycles, roles, artifacts, ceremonies); Kanban (historical context, product backlog, Kanban board, WIP limits)\end{enumerate}}\\
\multicolumn{4}{|p{\dimexpr 16cm + 6\tabcolsep\relax}|}{\begin{enumerate}[start=2]\item Requirements Engineering: types of requirements (functional, non-functional, interface); object-oriented requirements elicitation and analysis techniques; requirements specification (behavior-driven development - BDD, user stories and test scenarios, use cases, UML use case diagrams); minimum viable product (MVP); A/B testing\end{enumerate}}\\
\multicolumn{4}{|p{\dimexpr 16cm + 6\tabcolsep\relax}|}{\begin{enumerate}[start=3]\item Project Management: methods and techniques for project management based on the 10 PMI - Project Management Institute knowledge areas (scope, time, cost, quality, human resources, communications, risk, procurement, stakeholders and integration); techniques for defining the product backlog and work breakdown structure (WBS); team mobilization (processes, organization, roles and responsibilities, traceability, problem solving); effort estimation and pricing; tools for monitoring and controlling software development processes (Jira and Git)\end{enumerate}}\\
\multicolumn{4}{|p{\dimexpr 16cm + 6\tabcolsep\relax}|}{\begin{enumerate}[start=4]\item Software as a Service (SaaS): introduction and historical context; basic concepts (Service-Oriented Architecture - SOA, Software as a Service, Cloud Computing); SaaS application architecture (client-server versus peer-to-peer architecture, web clients and scripting languages, communication via HTTP and URIs, representation via HTML and CSS, 3-tier architecture and horizontal scalability, Model-View-Controller - MVC architecture); SaaS application performance (responsiveness and availability, monitoring and bottleneck identification, use of caching, database queries); security: protection of customer data and the Brazilian General Data Protection Law (LGPD)\end{enumerate}}\\
\multicolumn{4}{|p{\dimexpr 16cm + 6\tabcolsep\relax}|}{\begin{enumerate}[start=5]\item DevOps: version control (concepts, tools, multi-repositories versus mono-repositories); Continuous Integration - CI (motivation, concepts such as master/trunk, branches, merges, integration conflicts, gitflow); Continuous Delivery - CD (deployment, delivery, feature flags, release flags and business flags); tools for source code version control and CI/CD pipeline automation (Git)\end{enumerate}}\\
\multicolumn{4}{|p{\dimexpr 16cm + 6\tabcolsep\relax}|}{}\\
\hline
\multicolumn{4}{|p{16cm}|}{Basic Bibliography:}\\*
\multicolumn{4}{|p{\dimexpr 16cm + 6\tabcolsep\relax}|}{\begin{itemize}\item Valente, M. T. Engenharia de Software Moderna: Princípios e Práticas para Desenvolvimento de Software com Produtividade. 1ª edição. 2020. ISBN 978-6500019506. Disponível on-line: https://engsoftmoderna.info/\end{itemize}}\\
\multicolumn{4}{|p{\dimexpr 16cm + 6\tabcolsep\relax}|}{\begin{itemize}\item Patterson, D. A.; Fox, A. Construindo Software como Serviço (SaaS): Uma abordagem Ágil Usando Computação em Nuvem. 1st ed. Strawberry Canyon LLC, 2015. ISBN 9780984881277. Disponível on-line: http://www.saasbook.info/translations/br ou http://www.saasbook.info/first-edition\end{itemize}}\\
\multicolumn{4}{|p{\dimexpr 16cm + 6\tabcolsep\relax}|}{\begin{itemize}\item HULL, Elizabeth; JACKSON, Ken; DICK, Jeremy. Requirements Engineering. 3. ed. Springer. ISBN 978184996404. https://pergamum.ufop.br/pergamum/biblioteca/fotos.php?cod\_acervo=1415964\end{itemize}}\\
\multicolumn{4}{|p{16cm}|}{}\\
\hline
\multicolumn{4}{|p{16cm}|}{Complementary Bibliography:}\\*
\multicolumn{4}{|p{\dimexpr 16cm + 6\tabcolsep\relax}|}{\begin{itemize}\item Sommerville, I. Engenharia de Software. 10º edição. São Paulo: Pearson, 2019. ISBN 9788543024974. https://plataforma.bvirtual.com.br/Leitor/Loader/168127/pdf/6?keep=True\end{itemize}}\\
\multicolumn{4}{|p{\dimexpr 16cm + 6\tabcolsep\relax}|}{\begin{itemize}\item Fowler, M. UML distilled: a brief guide to the standard object modeling language. 3nd ed. Boston: Addison Wesley, c2004. xxx, 175 p. (Object technology series). ISBN 9780321193681. https://integrada.minhabiblioteca.com.br/\#/books/9788560031382/cfi/6/2!/4/18/52@0:76.8\end{itemize}}\\
\multicolumn{4}{|p{\dimexpr 16cm + 6\tabcolsep\relax}|}{\begin{itemize}\item Gamma, E. Padrões de Projetos: Soluções Reutilizáveis de Software Orientados a Objetos / Erich Gamma... [et al.]. Reading: Addison Wesley, 1995. 395 p. (Addison-Wesley Professional Computing Series). ISBN 9788573076103. https://doku.pub/download/padroes-de-projeto-soluoes-reutilizaveis-de-software-orientado-a-objetos-erich-gammapdf-6lkvge8em804\end{itemize}}\\
\multicolumn{4}{|p{\dimexpr 16cm + 6\tabcolsep\relax}|}{\begin{itemize}\item PMBOK® Guide (2021). https://www.pmi.org/pmbok-guide-standards/\end{itemize}}\\
\multicolumn{4}{|p{\dimexpr 16cm + 6\tabcolsep\relax}|}{\begin{itemize}\item CRISPIN, Lisa; GREGORY, Janet. Agile testing: a practical guide for testers and agile teams. Upper Saddle River, NJ: Addison Wesley, c2009. xli, 533 p. (The Addison-Wesley signature series). ISBN 9780321534460.\end{itemize}}\\
\hline
\end{longtable}
\end{center}

\clearpage
\begin{center}
\begin{tabular}{lccr}
 \multirow{3}{*}{\includegraphics[height=2.7cm]{brasao.png}} &
 \multicolumn{2}{c}{\bfseries FEDERAL UNIVERSITY OF OURO PRETO} &
 \ \ \ \ \multirow{3}{*}{\includegraphics[height=2.7cm]{ufop.png}} \\
 & \multicolumn{2}{c}{\bfseries DEAN'S OFFICE FOR UNDERGRADUATE STUDIES} & \\
 & \multicolumn{2}{c}{} & \\
 & \multicolumn{2}{c}{\Large\bfseries COURSE PROGRAM} & \\
\end{tabular}
\end{center}

\begin{center}
\begin{longtable}{|p{4cm}|p{4cm}|p{4cm}|p{4cm}|}
\hline
\multicolumn{3}{|p{12cm}|}{Course Name in English:} &
\multicolumn{1}{p{4cm}|}{Code:} \\ 
\multicolumn{3}{|p{12cm}|}{\textbf{Distributed Systems}} &
\textbf{BCC195}\\ 
\multicolumn{3}{|p{12cm}|}{Course Name in Portuguese:} & \\ 
\multicolumn{3}{|p{12cm}|}{\textbf{Sistemas Distribuídos}} & \\ 
\hline
\multicolumn{3}{|p{12cm}|}{Department Name and Acronym:} & Academic Unit: \\ 
\multicolumn{3}{|p{12cm}|}{Department of Computer Science (DECOM)} & {ICEB} \\ 
\hline
\multicolumn{4}{|p{16cm}|}{Delivery Mode:
[X] in person \hspace{0.6cm}
[ ] blended \hspace{0.6cm}
[ ] distance learning}\\
\hline
\multicolumn{2}{|p{8cm}|}{Semester workload} &
\multicolumn{2}{p{8cm}|}{Weekly workload}\\
\hline
\multicolumn{1}{|p{4cm}|}{Total} &
\multicolumn{1}{p{4cm}|}{Extension} &
\multicolumn{1}{p{4cm}|}{Theoretical} &
\multicolumn{1}{p{4cm}|}{Practical} \\ 
\multicolumn{1}{|p{4cm}|}{60\,hours} &
\multicolumn{1}{p{4cm}|}{0\;hours} &
\multicolumn{1}{p{4cm}|}{4\;class hours} &
\multicolumn{1}{p{4cm}|}{0\;class hours} \\ 
\hline
\multicolumn{4}{|p{16cm}|}{Course Description:}\\
\multicolumn{4}{|p{16cm}|}{}\\
\multicolumn{4}{|p{\dimexpr 16cm + 6\tabcolsep\relax}|}{Concept of distributed systems. Architecture of distributed systems. Processes and inter-process communication. Naming service in distributed systems. Synchronization service in distributed systems. Consistency and replication service in distributed systems. Fault tolerance service in distributed systems. Distributed platforms and applications.}\\
\multicolumn{4}{|p{16cm}|}{}\\
\hline
\multicolumn{4}{|p{16cm}|}{Course Content:}\\*
\multicolumn{4}{|p{\dimexpr 16cm + 6\tabcolsep\relax}|}{\begin{enumerate}[start=1]\item Introduction to distributed systems: architectures, examples, evolution, and remaining open challenges\end{enumerate}}\\
\multicolumn{4}{|p{\dimexpr 16cm + 6\tabcolsep\relax}|}{\begin{enumerate}[start=2]\item Review: processes, inter-process communication, concurrent programming, and network programming\end{enumerate}}\\
\multicolumn{4}{|p{\dimexpr 16cm + 6\tabcolsep\relax}|}{\begin{enumerate}[start=3]\item Concept of virtualization and introduction to virtual environments\end{enumerate}}\\
\multicolumn{4}{|p{\dimexpr 16cm + 6\tabcolsep\relax}|}{\begin{enumerate}[start=4]\item Examples of distributed platforms operating in computational clouds\end{enumerate}}\\
\multicolumn{4}{|p{\dimexpr 16cm + 6\tabcolsep\relax}|}{\begin{enumerate}[start=5]\item How to name and locate services in distributed environments\end{enumerate}}\\
\multicolumn{4}{|p{\dimexpr 16cm + 6\tabcolsep\relax}|}{\begin{enumerate}[start=6]\item Mutual exclusion and distributed synchronization protocols\end{enumerate}}\\
\multicolumn{4}{|p{\dimexpr 16cm + 6\tabcolsep\relax}|}{\begin{enumerate}[start=7]\item The consistency problem and replication protocols\end{enumerate}}\\
\multicolumn{4}{|p{\dimexpr 16cm + 6\tabcolsep\relax}|}{\begin{enumerate}[start=8]\item Types of failure and protocols for fault tolerance and failure recovery\end{enumerate}}\\
\multicolumn{4}{|p{\dimexpr 16cm + 6\tabcolsep\relax}|}{\begin{enumerate}[start=9]\item Examples of distributed applications (routing, ordering, graph problems, information retrieval, data mining, databases, among others)\end{enumerate}}\\
\multicolumn{4}{|p{\dimexpr 16cm + 6\tabcolsep\relax}|}{}\\
\hline
\multicolumn{4}{|p{16cm}|}{Basic Bibliography:}\\*
\multicolumn{4}{|p{\dimexpr 16cm + 6\tabcolsep\relax}|}{\begin{itemize}\item TANENBAUM, Andrew S.; STEEN, Maarten Van. Sistemas distribuídos: princípios e paradigmas. 2. ed. São Paulo: Pearson Prentice Hall, 2007.\end{itemize}}\\
\multicolumn{4}{|p{\dimexpr 16cm + 6\tabcolsep\relax}|}{\begin{itemize}\item DOLLIMORE, Jean; KINDBERG, Tim; TORTELLO, João; CARISSIMI, Alexandre; COULOURIS, George. Sistemas distribuídos: conceitos e projeto. 5. ed. Porto Alegre: Bookman, 2013.\end{itemize}}\\
\multicolumn{4}{|p{16cm}|}{}\\
\hline
\multicolumn{4}{|p{16cm}|}{Complementary Bibliography:}\\*
\multicolumn{4}{|p{\dimexpr 16cm + 6\tabcolsep\relax}|}{\begin{itemize}\item KLEPPMANN, Martin. Designing Data-Intensive Applications: The Big Ideas Behind Reliable, Scalable, and Maintainable Systems. 1. ed. São Paulo: Oreilly \& Assoc, 2015.\end{itemize}}\\
\multicolumn{4}{|p{\dimexpr 16cm + 6\tabcolsep\relax}|}{\begin{itemize}\item VERAS, Manoel. Computação em Nuvem. 1. ed. Brasport, 2015. Disponível online em MinhaUFOP/BibliotecaDigital/BVirtualPearson: https://plataforma.bvirtual.com.br/Acervo/Publicacao/160695\end{itemize}}\\
\multicolumn{4}{|p{\dimexpr 16cm + 6\tabcolsep\relax}|}{\begin{itemize}\item SANTOS, Ricardo Ribeiro dos. FACOM – UFMS, 2020. Material de Sistemas Distribuídos. Disponível em: https://www.facom.ufms.br/~ricardo/Courses/DisSys/Material/ . Acesso em: 07/12/2020.\end{itemize}}\\
\multicolumn{4}{|p{\dimexpr 16cm + 6\tabcolsep\relax}|}{\begin{itemize}\item RODRIGUEZ, Noemi. PUC-RJ, 2020. Material de Sistemas Distribuídos. Disponível em: http://www.inf.puc-rio.br/~noemi/sd-19/ . Acesso em: 07/12/2020.\end{itemize}}\\
\multicolumn{4}{|p{\dimexpr 16cm + 6\tabcolsep\relax}|}{\begin{itemize}\item FERNANDEZ, Marcial Porto. UECE, 2020. Material de Sistemas Distribuídos. Disponível em: http://marcial.larces.uece.br/cursos/sistemas-distribuidos-2018-1 . Acesso em: 07/12/2020.\end{itemize}}\\
\multicolumn{4}{|p{\dimexpr 16cm + 6\tabcolsep\relax}|}{\begin{itemize}\item FIGUEIREDO, Daniel Ratton. UFRJ, 2020. Material de sistemas distribuídos. Disponível em: https://www.cos.ufrj.br/~daniel/sd/ . Acesso em: 07/12/2020.\end{itemize}}\\
\hline
\end{longtable}
\end{center}

\clearpage
\begin{center}
\begin{tabular}{lccr}
 \multirow{3}{*}{\includegraphics[height=2.7cm]{brasao.png}} &
 \multicolumn{2}{c}{\bfseries FEDERAL UNIVERSITY OF OURO PRETO} &
 \ \ \ \ \multirow{3}{*}{\includegraphics[height=2.7cm]{ufop.png}} \\
 & \multicolumn{2}{c}{\bfseries DEAN'S OFFICE FOR UNDERGRADUATE STUDIES} & \\
 & \multicolumn{2}{c}{} & \\
 & \multicolumn{2}{c}{\Large\bfseries COURSE PROGRAM} & \\
\end{tabular}
\end{center}

\begin{center}
\begin{longtable}{|p{4cm}|p{4cm}|p{4cm}|p{4cm}|}
\hline
\multicolumn{3}{|p{12cm}|}{Course Name in English:} &
\multicolumn{1}{p{4cm}|}{Code:} \\ 
\multicolumn{3}{|p{12cm}|}{\textbf{Databases II}} &
\textbf{BCC281}\\ 
\multicolumn{3}{|p{12cm}|}{Course Name in Portuguese:} & \\ 
\multicolumn{3}{|p{12cm}|}{\textbf{Banco de Dados II}} & \\ 
\hline
\multicolumn{3}{|p{12cm}|}{Department Name and Acronym:} & Academic Unit: \\ 
\multicolumn{3}{|p{12cm}|}{Department of Computer Science (DECOM)} & {ICEB} \\ 
\hline
\multicolumn{4}{|p{16cm}|}{Delivery Mode:
[X] in person \hspace{0.6cm}
[ ] blended \hspace{0.6cm}
[ ] distance learning}\\
\hline
\multicolumn{2}{|p{8cm}|}{Semester workload} &
\multicolumn{2}{p{8cm}|}{Weekly workload}\\
\hline
\multicolumn{1}{|p{4cm}|}{Total} &
\multicolumn{1}{p{4cm}|}{Extension} &
\multicolumn{1}{p{4cm}|}{Theoretical} &
\multicolumn{1}{p{4cm}|}{Practical} \\ 
\multicolumn{1}{|p{4cm}|}{60\,hours} &
\multicolumn{1}{p{4cm}|}{0\;hours} &
\multicolumn{1}{p{4cm}|}{4\;class hours} &
\multicolumn{1}{p{4cm}|}{0\;class hours} \\ 
\hline
\multicolumn{4}{|p{16cm}|}{Course Description:}\\
\multicolumn{4}{|p{16cm}|}{}\\
\multicolumn{4}{|p{\dimexpr 16cm + 6\tabcolsep\relax}|}{Transaction processing. Concurrency control. Failure recovery. Stored functions and triggers. Query optimization. Object-oriented databases. Introduction to XML. Advanced database topics.}\\
\multicolumn{4}{|p{16cm}|}{}\\
\hline
\multicolumn{4}{|p{16cm}|}{Course Content:}\\*
\multicolumn{4}{|p{\dimexpr 16cm + 6\tabcolsep\relax}|}{\begin{enumerate}[start=1]\item Transaction processing: theory of transaction processing, transaction execution schedules\end{enumerate}}\\
\multicolumn{4}{|p{\dimexpr 16cm + 6\tabcolsep\relax}|}{\begin{enumerate}[start=2]\item Concurrency control: theory of concurrency control, concurrency control techniques\end{enumerate}}\\
\multicolumn{4}{|p{\dimexpr 16cm + 6\tabcolsep\relax}|}{\begin{enumerate}[start=3]\item Failure recovery: theory of failure recovery, failure recovery techniques\end{enumerate}}\\
\multicolumn{4}{|p{\dimexpr 16cm + 6\tabcolsep\relax}|}{\begin{enumerate}[start=4]\item Stored functions and triggers: creation and use of stored functions, creation and use of triggers\end{enumerate}}\\
\multicolumn{4}{|p{\dimexpr 16cm + 6\tabcolsep\relax}|}{\begin{enumerate}[start=5]\item Query optimization: algorithms for query processing and optimization, processing cost estimation\end{enumerate}}\\
\multicolumn{4}{|p{\dimexpr 16cm + 6\tabcolsep\relax}|}{\begin{enumerate}[start=6]\item Object-oriented databases: theory of object-oriented databases, object definition language, object query language\end{enumerate}}\\
\multicolumn{4}{|p{\dimexpr 16cm + 6\tabcolsep\relax}|}{\begin{enumerate}[start=7]\item Introduction to XML: XML syntax, XML and semi-structured data, XML and relational databases\end{enumerate}}\\
\multicolumn{4}{|p{\dimexpr 16cm + 6\tabcolsep\relax}|}{\begin{enumerate}[start=8]\item Advanced database topics: data storage and indexing, spatial databases, distributed databases, text databases, non-relational databases, digital libraries, data warehouses\end{enumerate}}\\
\multicolumn{4}{|p{\dimexpr 16cm + 6\tabcolsep\relax}|}{}\\
\hline
\multicolumn{4}{|p{16cm}|}{Basic Bibliography:}\\*
\multicolumn{4}{|p{\dimexpr 16cm + 6\tabcolsep\relax}|}{\begin{itemize}\item ELMASRI, Ramez; NAVATHE, Shamkant B. Fundamentals of Database Systems. 7. ed. Hoboken, NJ: Pearson, 2016.\end{itemize}}\\
\multicolumn{4}{|p{\dimexpr 16cm + 6\tabcolsep\relax}|}{\begin{itemize}\item SILBERSCHATZ, Abraham; KORTH, Henry F.; SUDARSHAN, S. Database System Concepts. 7. ed. New York: McGraw-Hill, 2019.\end{itemize}}\\
\multicolumn{4}{|p{\dimexpr 16cm + 6\tabcolsep\relax}|}{\begin{itemize}\item DATE, Christopher J. An Introduction to Database Systems. 8. ed. Boston: Pearson/Addison Wesley, 2003.\end{itemize}}\\
\multicolumn{4}{|p{16cm}|}{}\\
\hline
\multicolumn{4}{|p{16cm}|}{Complementary Bibliography:}\\*
\multicolumn{4}{|p{\dimexpr 16cm + 6\tabcolsep\relax}|}{\begin{itemize}\item RAMAKRISHNAN, Raghu; GEHRKE, Johannes. Sistemas de gerenciamento de banco de dados. 3. ed. São Paulo: McGraw-Hill, 2008. 905 p.\end{itemize}}\\
\multicolumn{4}{|p{\dimexpr 16cm + 6\tabcolsep\relax}|}{\begin{itemize}\item ELMASRI, Ramez; NAVATHE, Sham. Sistemas de banco de dados. 6. ed. São Paulo: Pearson, Addison Wesley, 2011. 788 p.\end{itemize}}\\
\multicolumn{4}{|p{\dimexpr 16cm + 6\tabcolsep\relax}|}{\begin{itemize}\item AMADEU, Claudia V. Banco de Dados. 1. ed. São Paulo: Pearson, 2015. 208 p.\end{itemize}}\\
\multicolumn{4}{|p{\dimexpr 16cm + 6\tabcolsep\relax}|}{\begin{itemize}\item MEDEIROS, Luciano F. Banco de Dados: princípios e prática. 1. ed. Editora Intersaberes, 2013. 188 p.\end{itemize}}\\
\multicolumn{4}{|p{\dimexpr 16cm + 6\tabcolsep\relax}|}{\begin{itemize}\item PUGA, Sandra; FRANÇA, Edson; GOYA, Milton. Banco de dados: implementação em SQL, PL/SQL e Oracle 11g. 1. ed. São Paulo: Pearson, 2013. 356 p.\end{itemize}}\\
\multicolumn{4}{|p{\dimexpr 16cm + 6\tabcolsep\relax}|}{\begin{itemize}\item SETZER, Valdemar W.; SILVA, Flávio S. C. Bancos de dados. 1. ed. Editora Blucher, 2005. 393 p.\end{itemize}}\\
\multicolumn{4}{|p{\dimexpr 16cm + 6\tabcolsep\relax}|}{\begin{itemize}\item CASANOVA, Marco; CÂMARA, Gilberto; DAVIS, Clodoveu; VINHAS, Lúbia; QUEIROZ, Gilberto R. Bancos de Dados Geográficos. 1. ed. Curitiba: MundoGEO, 2005.\end{itemize}}\\
\hline
\end{longtable}
\end{center}

\clearpage
\begin{center}
\begin{tabular}{lccr}
 \multirow{3}{*}{\includegraphics[height=2.7cm]{brasao.png}} &
 \multicolumn{2}{c}{\bfseries FEDERAL UNIVERSITY OF OURO PRETO} &
 \ \ \ \ \multirow{3}{*}{\includegraphics[height=2.7cm]{ufop.png}} \\
 & \multicolumn{2}{c}{\bfseries DEAN'S OFFICE FOR UNDERGRADUATE STUDIES} & \\
 & \multicolumn{2}{c}{} & \\
 & \multicolumn{2}{c}{\Large\bfseries COURSE PROGRAM} & \\
\end{tabular}
\end{center}

\begin{center}
\begin{longtable}{|p{4cm}|p{4cm}|p{4cm}|p{4cm}|}
\hline
\multicolumn{3}{|p{12cm}|}{Course Name in English:} &
\multicolumn{1}{p{4cm}|}{Code:} \\ 
\multicolumn{3}{|p{12cm}|}{\textbf{Web Programming}} &
\textbf{BCC383}\\ 
\multicolumn{3}{|p{12cm}|}{Course Name in Portuguese:} & \\ 
\multicolumn{3}{|p{12cm}|}{\textbf{Programação Web}} & \\ 
\hline
\multicolumn{3}{|p{12cm}|}{Department Name and Acronym:} & Academic Unit: \\ 
\multicolumn{3}{|p{12cm}|}{Department of Computer Science (DECOM)} & {ICEB} \\ 
\hline
\multicolumn{4}{|p{16cm}|}{Delivery Mode:
[X] in person \hspace{0.6cm}
[ ] blended \hspace{0.6cm}
[ ] distance learning}\\
\hline
\multicolumn{2}{|p{8cm}|}{Semester workload} &
\multicolumn{2}{p{8cm}|}{Weekly workload}\\
\hline
\multicolumn{1}{|p{4cm}|}{Total} &
\multicolumn{1}{p{4cm}|}{Extension} &
\multicolumn{1}{p{4cm}|}{Theoretical} &
\multicolumn{1}{p{4cm}|}{Practical} \\ 
\multicolumn{1}{|p{4cm}|}{60\,hours} &
\multicolumn{1}{p{4cm}|}{0\;hours} &
\multicolumn{1}{p{4cm}|}{4\;class hours} &
\multicolumn{1}{p{4cm}|}{0\;class hours} \\ 
\hline
\multicolumn{4}{|p{16cm}|}{Course Description:}\\
\multicolumn{4}{|p{16cm}|}{}\\
\multicolumn{4}{|p{\dimexpr 16cm + 6\tabcolsep\relax}|}{Markup languages and style sheets; Front-end development topics; Server-side programming; How browsers work; Libraries and frameworks for Web application development.}\\
\multicolumn{4}{|p{16cm}|}{}\\
\hline
\multicolumn{4}{|p{16cm}|}{Course Content:}\\*
\multicolumn{4}{|p{\dimexpr 16cm + 6\tabcolsep\relax}|}{\begin{enumerate}[start=1]\item Markup languages and style sheets: HTML (basic concepts and best practices, semantic web, responsiveness); CSS (basic concepts and best practices, block vs inline, selectors, box model, flexbox, position, responsiveness, animations)\end{enumerate}}\\
\multicolumn{4}{|p{\dimexpr 16cm + 6\tabcolsep\relax}|}{\begin{enumerate}[start=2]\item Front-end development topics: basic scripting language concepts, Document Object Model (DOM), events, browser extensions, asynchronous events (fetch API, Promise), JSON\end{enumerate}}\\
\multicolumn{4}{|p{\dimexpr 16cm + 6\tabcolsep\relax}|}{\begin{enumerate}[start=3]\item Server-side programming: HTTP requests, CORS, routes, databases, authentication\end{enumerate}}\\
\multicolumn{4}{|p{\dimexpr 16cm + 6\tabcolsep\relax}|}{\begin{enumerate}[start=4]\item How browsers work: history, design principles\end{enumerate}}\\
\multicolumn{4}{|p{\dimexpr 16cm + 6\tabcolsep\relax}|}{\begin{enumerate}[start=5]\item Libraries and frameworks for Web application development\end{enumerate}}\\
\multicolumn{4}{|p{\dimexpr 16cm + 6\tabcolsep\relax}|}{}\\
\hline
\multicolumn{4}{|p{16cm}|}{Basic Bibliography:}\\*
\multicolumn{4}{|p{\dimexpr 16cm + 6\tabcolsep\relax}|}{\begin{itemize}\item FLANAGAN, David. JavaScript: o guia definitivo. Porto Alegre: Bookman, 2013. Disponível em: https://integrada.minhabiblioteca.com.br/\#/books/9788565837484/. Acesso em: 31 jul. 2020.\end{itemize}}\\
\multicolumn{4}{|p{\dimexpr 16cm + 6\tabcolsep\relax}|}{\begin{itemize}\item FLATSCHART, Fábio. HTML5 Embarque imediato. Rio de Janeiro: Brasport, 2011. Disponível em: https://plataforma.bvirtual.com.br/Leitor/Publicacao/160673/epub/0. Acesso em: 31 jul. 2020.\end{itemize}}\\
\multicolumn{4}{|p{\dimexpr 16cm + 6\tabcolsep\relax}|}{\begin{itemize}\item SEGURADO, Valquiria Santos. Projeto de interface com o usuário. 1. ed. São Paulo: Pearson Education do Brasil, 2016. Disponível em: https://plataforma.bvirtual.com.br/Leitor/Publicacao/124143/pdf/0. Acesso em: 31 jul. 2020.\end{itemize}}\\
\multicolumn{4}{|p{16cm}|}{}\\
\hline
\multicolumn{4}{|p{16cm}|}{Complementary Bibliography:}\\*
\multicolumn{4}{|p{\dimexpr 16cm + 6\tabcolsep\relax}|}{\begin{itemize}\item BANIN, Sérgio Luiz. Python 3: conceitos e aplicações: uma abordagem didática. São Paulo: Érica, 2018. Disponível em: https://integrada.minhabiblioteca.com.br/\#/books/9788536530253. Acesso em: 02 ago. 2020.\end{itemize}}\\
\multicolumn{4}{|p{\dimexpr 16cm + 6\tabcolsep\relax}|}{\begin{itemize}\item CAELUM. Apostila do curso: Desenvolvimento Web com HTML, CSS e JavaScript. 2020. Disponível em: https://www.caelum.com.br/apostila-html-css-javascript/. Acesso em: 31 jul. 2020.\end{itemize}}\\
\multicolumn{4}{|p{\dimexpr 16cm + 6\tabcolsep\relax}|}{\begin{itemize}\item HAVERBEKE, Marijn. Eloquent JavaScript: A Modern Introduction to Programming. 3. ed., 2018. Disponível em: https://eloquentjavascript.net/. Acesso em: 02 ago. 2020.\end{itemize}}\\
\multicolumn{4}{|p{\dimexpr 16cm + 6\tabcolsep\relax}|}{\begin{itemize}\item MILETTO, Evandro Manara; BERTAGNOLLI, Silvia de Castro. Desenvolvimento de Software II: Introdução ao Desenvolvimento Web com HTML, CSS, JavaScript e PHP. Porto Alegre: Bookman, 2014. Disponível em: https://integrada.minhabiblioteca.com.br/\#/books/9788582601969. Acesso em: 02 ago. 2020.\end{itemize}}\\
\multicolumn{4}{|p{\dimexpr 16cm + 6\tabcolsep\relax}|}{\begin{itemize}\item SOUSA, Roque Fernando Marcos. Canvas HTML5. Rio de Janeiro: Brasport, 2013. Disponível em: https://plataforma.bvirtual.com.br/Leitor/Publicacao/160686/epub/0. Acesso em: 02 ago. 2020.\end{itemize}}\\
\hline
\end{longtable}
\end{center}

\clearpage
\begin{center}
\begin{tabular}{lccr}
 \multirow{3}{*}{\includegraphics[height=2.7cm]{brasao.png}} &
 \multicolumn{2}{c}{\bfseries FEDERAL UNIVERSITY OF OURO PRETO} &
 \ \ \ \ \multirow{3}{*}{\includegraphics[height=2.7cm]{ufop.png}} \\
 & \multicolumn{2}{c}{\bfseries DEAN'S OFFICE FOR UNDERGRADUATE STUDIES} & \\
 & \multicolumn{2}{c}{} & \\
 & \multicolumn{2}{c}{\Large\bfseries COURSE PROGRAM} & \\
\end{tabular}
\end{center}

\begin{center}
\begin{longtable}{|p{4cm}|p{4cm}|p{4cm}|p{4cm}|}
\hline
\multicolumn{3}{|p{12cm}|}{Course Name in English:} &
\multicolumn{1}{p{4cm}|}{Code:} \\ 
\multicolumn{3}{|p{12cm}|}{\textbf{Complex Networks}} &
\textbf{BCC284}\\ 
\multicolumn{3}{|p{12cm}|}{Course Name in Portuguese:} & \\ 
\multicolumn{3}{|p{12cm}|}{\textbf{Redes Complexas}} & \\ 
\hline
\multicolumn{3}{|p{12cm}|}{Department Name and Acronym:} & Academic Unit: \\ 
\multicolumn{3}{|p{12cm}|}{Department of Computer Science (DECOM)} & {ICEB} \\ 
\hline
\multicolumn{4}{|p{16cm}|}{Delivery Mode:
[X] in person \hspace{0.6cm}
[ ] blended \hspace{0.6cm}
[ ] distance learning}\\
\hline
\multicolumn{2}{|p{8cm}|}{Semester workload} &
\multicolumn{2}{p{8cm}|}{Weekly workload}\\
\hline
\multicolumn{1}{|p{4cm}|}{Total} &
\multicolumn{1}{p{4cm}|}{Extension} &
\multicolumn{1}{p{4cm}|}{Theoretical} &
\multicolumn{1}{p{4cm}|}{Practical} \\ 
\multicolumn{1}{|p{4cm}|}{60\,hours} &
\multicolumn{1}{p{4cm}|}{0\;hours} &
\multicolumn{1}{p{4cm}|}{4\;class hours} &
\multicolumn{1}{p{4cm}|}{0\;class hours} \\ 
\hline
\multicolumn{4}{|p{16cm}|}{Course Description:}\\
\multicolumn{4}{|p{16cm}|}{}\\
\multicolumn{4}{|p{\dimexpr 16cm + 6\tabcolsep\relax}|}{Introduction and basic concepts; Directed and weighted networks; Models and algorithms for generating complex networks; Degree correlations; Network robustness; Community detection; Dynamic network phenomena; Advanced topics.}\\
\multicolumn{4}{|p{16cm}|}{}\\
\hline
\multicolumn{4}{|p{16cm}|}{Course Content:}\\*
\multicolumn{4}{|p{\dimexpr 16cm + 6\tabcolsep\relax}|}{\begin{enumerate}[start=1]\item Introduction and basic concepts: types of networks and applications; adjacency matrix, adjacency list, weighted networks; paths and distances, centrality measures, connectivity and transitivity\end{enumerate}}\\
\multicolumn{4}{|p{\dimexpr 16cm + 6\tabcolsep\relax}|}{\begin{enumerate}[start=2]\item Directed and weighted networks: directed networks, the Web, PageRank\end{enumerate}}\\
\multicolumn{4}{|p{\dimexpr 16cm + 6\tabcolsep\relax}|}{\begin{enumerate}[start=3]\item Models and algorithms for generating complex networks: random networks, small-world networks, scale-free property, Barabási-Albert model, dynamic networks\end{enumerate}}\\
\multicolumn{4}{|p{\dimexpr 16cm + 6\tabcolsep\relax}|}{\begin{enumerate}[start=4]\item Degree correlations: assortativity\end{enumerate}}\\
\multicolumn{4}{|p{\dimexpr 16cm + 6\tabcolsep\relax}|}{\begin{enumerate}[start=5]\item Network robustness: percolation, failures and attacks\end{enumerate}}\\
\multicolumn{4}{|p{\dimexpr 16cm + 6\tabcolsep\relax}|}{\begin{enumerate}[start=6]\item Community detection: concepts and algorithms\end{enumerate}}\\
\multicolumn{4}{|p{\dimexpr 16cm + 6\tabcolsep\relax}|}{\begin{enumerate}[start=7]\item Dynamic network phenomena: viral spreading and information diffusion, synchronization of coupled oscillators\end{enumerate}}\\
\multicolumn{4}{|p{\dimexpr 16cm + 6\tabcolsep\relax}|}{\begin{enumerate}[start=8]\item Advanced topics\end{enumerate}}\\
\multicolumn{4}{|p{\dimexpr 16cm + 6\tabcolsep\relax}|}{}\\
\hline
\multicolumn{4}{|p{16cm}|}{Basic Bibliography:}\\*
\multicolumn{4}{|p{\dimexpr 16cm + 6\tabcolsep\relax}|}{\begin{itemize}\item BARABÁSI, A.-L. Network science. Cambridge: Cambridge University Press, 2016. 475 p.\end{itemize}}\\
\multicolumn{4}{|p{\dimexpr 16cm + 6\tabcolsep\relax}|}{\begin{itemize}\item MENCZER, F.; FORTUNATO, S.; DAVIS, C. A First Course in Network Science. Cambridge: Cambridge University Press, 2020. 300 p.\end{itemize}}\\
\multicolumn{4}{|p{\dimexpr 16cm + 6\tabcolsep\relax}|}{\begin{itemize}\item NEWMAN, M. Networks. 2. ed. Oxford: Oxford University Press, 2018. 800 p.\end{itemize}}\\
\multicolumn{4}{|p{16cm}|}{}\\
\hline
\multicolumn{4}{|p{16cm}|}{Complementary Bibliography:}\\*
\multicolumn{4}{|p{\dimexpr 16cm + 6\tabcolsep\relax}|}{\begin{itemize}\item BARABÁSI, A.-L. Linked: how everything is connected to everything else and what it means for business, science and everyday life. New York: Plume, 2003. 294 p.\end{itemize}}\\
\multicolumn{4}{|p{\dimexpr 16cm + 6\tabcolsep\relax}|}{\begin{itemize}\item BARRAT, A.; MELY, M. B.; VESPIGNANI, A. Dynamical Processes on Complex Networks. Cambridge: Cambridge University Press, 2012. 361 p.\end{itemize}}\\
\multicolumn{4}{|p{\dimexpr 16cm + 6\tabcolsep\relax}|}{\begin{itemize}\item EASLEY, D.; KLEINBERG, J. Networks, Crowds, and Markets: Reasoning about a Highly Connected World. Cambridge: Cambridge University Press, 2010. 727 p.\end{itemize}}\\
\multicolumn{4}{|p{\dimexpr 16cm + 6\tabcolsep\relax}|}{\begin{itemize}\item ESTRADA, E. The structure of complex networks: theory and applications. Oxford: Oxford University Press, 2012. 478 p.\end{itemize}}\\
\multicolumn{4}{|p{\dimexpr 16cm + 6\tabcolsep\relax}|}{\begin{itemize}\item JACKSON, M. O. Social and Economic Networks. Princeton: Princeton University Press, 2010. 504 p.\end{itemize}}\\
\hline
\end{longtable}
\end{center}

\clearpage
\begin{center}
\begin{tabular}{lccr}
 \multirow{3}{*}{\includegraphics[height=2.7cm]{brasao.png}} &
 \multicolumn{2}{c}{\bfseries FEDERAL UNIVERSITY OF OURO PRETO} &
 \ \ \ \ \multirow{3}{*}{\includegraphics[height=2.7cm]{ufop.png}} \\
 & \multicolumn{2}{c}{\bfseries DEAN'S OFFICE FOR UNDERGRADUATE STUDIES} & \\
 & \multicolumn{2}{c}{} & \\
 & \multicolumn{2}{c}{\Large\bfseries COURSE PROGRAM} & \\
\end{tabular}
\end{center}

\begin{center}
\begin{longtable}{|p{4cm}|p{4cm}|p{4cm}|p{4cm}|}
\hline
\multicolumn{3}{|p{12cm}|}{Course Name in English:} &
\multicolumn{1}{p{4cm}|}{Code:} \\ 
\multicolumn{3}{|p{12cm}|}{\textbf{Data Engineering}} &
\textbf{BCC285}\\ 
\multicolumn{3}{|p{12cm}|}{Course Name in Portuguese:} & \\ 
\multicolumn{3}{|p{12cm}|}{\textbf{Engenharia de Dados}} & \\ 
\hline
\multicolumn{3}{|p{12cm}|}{Department Name and Acronym:} & Academic Unit: \\ 
\multicolumn{3}{|p{12cm}|}{Department of Computer Science (DECOM)} & {ICEB} \\ 
\hline
\multicolumn{4}{|p{16cm}|}{Delivery Mode:
[X] in person \hspace{0.6cm}
[ ] blended \hspace{0.6cm}
[ ] distance learning}\\
\hline
\multicolumn{2}{|p{8cm}|}{Semester workload} &
\multicolumn{2}{p{8cm}|}{Weekly workload}\\
\hline
\multicolumn{1}{|p{4cm}|}{Total} &
\multicolumn{1}{p{4cm}|}{Extension} &
\multicolumn{1}{p{4cm}|}{Theoretical} &
\multicolumn{1}{p{4cm}|}{Practical} \\ 
\multicolumn{1}{|p{4cm}|}{60\,hours} &
\multicolumn{1}{p{4cm}|}{0\;hours} &
\multicolumn{1}{p{4cm}|}{2\;class hours} &
\multicolumn{1}{p{4cm}|}{2\;class hours} \\ 
\hline
\multicolumn{4}{|p{16cm}|}{Course Description:}\\
\multicolumn{4}{|p{16cm}|}{}\\
\multicolumn{4}{|p{\dimexpr 16cm + 6\tabcolsep\relax}|}{Virtualization. Cloud computing. Clusters, virtual machines, containers, and container orchestrators. Data ingestion. Data modeling, storage, and querying. Data pipeline: design, implementation, and testing.}\\
\multicolumn{4}{|p{16cm}|}{}\\
\hline
\multicolumn{4}{|p{16cm}|}{Course Content:}\\*
\multicolumn{4}{|p{\dimexpr 16cm + 6\tabcolsep\relax}|}{\begin{enumerate}[start=1]\item Course presentation: syllabus, objectives, bibliography\end{enumerate}}\\
\multicolumn{4}{|p{\dimexpr 16cm + 6\tabcolsep\relax}|}{\begin{enumerate}[start=2]\item Virtualization concepts and overview of the main cloud environments\end{enumerate}}\\
\multicolumn{4}{|p{\dimexpr 16cm + 6\tabcolsep\relax}|}{\begin{enumerate}[start=3]\item Basic operations for installing and managing a cluster of virtual machines\end{enumerate}}\\
\multicolumn{4}{|p{\dimexpr 16cm + 6\tabcolsep\relax}|}{\begin{enumerate}[start=4]\item Basic operations for installing and managing containers and their orchestrators\end{enumerate}}\\
\multicolumn{4}{|p{\dimexpr 16cm + 6\tabcolsep\relax}|}{\begin{enumerate}[start=5]\item Review: basic Python concepts, the pub/sub pattern, and messaging\end{enumerate}}\\
\multicolumn{4}{|p{\dimexpr 16cm + 6\tabcolsep\relax}|}{\begin{enumerate}[start=6]\item Notebooks: concept, tools, and examples of best practices\end{enumerate}}\\
\multicolumn{4}{|p{\dimexpr 16cm + 6\tabcolsep\relax}|}{\begin{enumerate}[start=7]\item Data ingestion: data types, existing architectures, and supporting tools\end{enumerate}}\\
\multicolumn{4}{|p{\dimexpr 16cm + 6\tabcolsep\relax}|}{\begin{enumerate}[start=8]\item Data modeling: relational, non-relational, and multidimensional\end{enumerate}}\\
\multicolumn{4}{|p{\dimexpr 16cm + 6\tabcolsep\relax}|}{\begin{enumerate}[start=9]\item Data storage: data ingestion into data lakes or data lakehouses\end{enumerate}}\\
\multicolumn{4}{|p{\dimexpr 16cm + 6\tabcolsep\relax}|}{\begin{enumerate}[start=10]\item Data querying: tools, examples, and operationalization\end{enumerate}}\\
\multicolumn{4}{|p{\dimexpr 16cm + 6\tabcolsep\relax}|}{\begin{enumerate}[start=11]\item Data pipeline: concepts and stages that make up a data pipeline; main tools for building and managing a data pipeline on virtual clusters; data pipeline bottlenecks and how to mitigate them; data pipeline examples; design, implementation, and testing of some data pipelines\end{enumerate}}\\
\multicolumn{4}{|p{\dimexpr 16cm + 6\tabcolsep\relax}|}{}\\
\hline
\multicolumn{4}{|p{16cm}|}{Basic Bibliography:}\\*
\multicolumn{4}{|p{\dimexpr 16cm + 6\tabcolsep\relax}|}{\begin{itemize}\item KLEPPMANN, Martin. Designing Data-Intensive Applications: The Big Ideas Behind Reliable, Scalable, and Maintainable Systems. 1. ed. São Paulo: Oreilly \& Assoc, 2015.\end{itemize}}\\
\multicolumn{4}{|p{\dimexpr 16cm + 6\tabcolsep\relax}|}{\begin{itemize}\item KUKREJA, Manoj. Data Engineering with Apache Spark, Delta Lake, and Lakehouse: Create scalable pipelines that ingest, curate, and aggregate complex data in a timely and secure way. 1. ed. Packt Publishing, 2021.\end{itemize}}\\
\multicolumn{4}{|p{\dimexpr 16cm + 6\tabcolsep\relax}|}{\begin{itemize}\item EAGAR, Gareth. Data Engineering with AWS: Learn how to design and build cloud-based data transformation pipelines using AWS. 1. ed. Packt Publishing, 2021.\end{itemize}}\\
\multicolumn{4}{|p{\dimexpr 16cm + 6\tabcolsep\relax}|}{\begin{itemize}\item JÚNIOR, Armando Kolbe. Computação em Nuvem. 1. ed. Contentus, 2020. Disponível online em MinhaUFOP/BibliotecaDigital/BVirtualPearson: https://plataforma.bvirtual.com.br/Acervo/Publicacao/184851\end{itemize}}\\
\multicolumn{4}{|p{\dimexpr 16cm + 6\tabcolsep\relax}|}{\begin{itemize}\item VITALINO, Jeferson Fernando Noronha; CASTRO, Marcus André Nunes. Descomplicando o Docker. 2. ed. Brasport, 2018. Disponível online em MinhaUFOP/BibliotecaDigital/BVirtualPearson: https://plataforma.bvirtual.com.br/Acervo/Publicacao/164690\end{itemize}}\\
\multicolumn{4}{|p{16cm}|}{}\\
\hline
\multicolumn{4}{|p{16cm}|}{Complementary Bibliography:}\\*
\multicolumn{4}{|p{\dimexpr 16cm + 6\tabcolsep\relax}|}{\begin{itemize}\item HARENSLAK, Bas P.; RUITER, Julian Rutger. Data Pipelines with Apache Airflow. 1. ed. Manning Publications, 2021.\end{itemize}}\\
\multicolumn{4}{|p{\dimexpr 16cm + 6\tabcolsep\relax}|}{\begin{itemize}\item Vários tutoriais sobre Jupyter notebooks. Disponível em: https://docs.jupyter.org/en/latest/ . Acesso em: 07/12/2021.\end{itemize}}\\
\multicolumn{4}{|p{\dimexpr 16cm + 6\tabcolsep\relax}|}{\begin{itemize}\item Vários tutoriais sobre Kubernetes. Disponível em: https://kubernetes.io/pt-br/docs/tutorials/ . Acesso em: 07/12/2021.\end{itemize}}\\
\multicolumn{4}{|p{\dimexpr 16cm + 6\tabcolsep\relax}|}{\begin{itemize}\item KOEN, Semi. Architecting a Machine Learning Pipeline. TowardsDataScience, 2019. Disponível em: https://towardsdatascience.com/architecting-a-machine-learning-pipeline-a847f094d1c7 . Acesso em: 07/12/2021.\end{itemize}}\\
\multicolumn{4}{|p{\dimexpr 16cm + 6\tabcolsep\relax}|}{\begin{itemize}\item YOSKOVICH, Tomer. Build your Data-Lake using AWS S3 \& Athena. TowardsDataScience, 2019. Disponível em: https://medium.com/swlh/tutorial-build-your-data-lake-using-aws-s3-athena-150c1aaa44cf . Acesso em: 07/12/2021.\end{itemize}}\\
\multicolumn{4}{|p{\dimexpr 16cm + 6\tabcolsep\relax}|}{\begin{itemize}\item CARAMELLO, Ivan. Keeping Data Under (Your) Control with Apache Airflow. TowardsDataScience, 2022. Disponível em: https://medium.com/daitan-tech/keeping-data-under-your-control-with-apache-airflowa46f884f8c5d . Acesso em: 07/12/2021.\end{itemize}}\\
\multicolumn{4}{|p{\dimexpr 16cm + 6\tabcolsep\relax}|}{\begin{itemize}\item BLANCAS, Eduardo. On writing clean Jupyter notebooks. TowardsDataScience, 2021. Disponível em: https://medium.com/towards-data-science/on-writing-clean-jupyter-notebooks-abdf6c708c75 . Acesso em: 07/12/2021.\end{itemize}}\\
\hline
\end{longtable}
\end{center}

\clearpage
\begin{center}
\begin{tabular}{lccr}
 \multirow{3}{*}{\includegraphics[height=2.7cm]{brasao.png}} &
 \multicolumn{2}{c}{\bfseries FEDERAL UNIVERSITY OF OURO PRETO} &
 \ \ \ \ \multirow{3}{*}{\includegraphics[height=2.7cm]{ufop.png}} \\
 & \multicolumn{2}{c}{\bfseries DEAN'S OFFICE FOR UNDERGRADUATE STUDIES} & \\
 & \multicolumn{2}{c}{} & \\
 & \multicolumn{2}{c}{\Large\bfseries COURSE PROGRAM} & \\
\end{tabular}
\end{center}

\begin{center}
\begin{longtable}{|p{4cm}|p{4cm}|p{4cm}|p{4cm}|}
\hline
\multicolumn{3}{|p{12cm}|}{Course Name in English:} &
\multicolumn{1}{p{4cm}|}{Code:} \\ 
\multicolumn{3}{|p{12cm}|}{\textbf{Data Visualization}} &
\textbf{BCC287}\\ 
\multicolumn{3}{|p{12cm}|}{Course Name in Portuguese:} & \\ 
\multicolumn{3}{|p{12cm}|}{\textbf{Visualização de Dados}} & \\ 
\hline
\multicolumn{3}{|p{12cm}|}{Department Name and Acronym:} & Academic Unit: \\ 
\multicolumn{3}{|p{12cm}|}{Department of Computer Science (DECOM)} & {ICEB} \\ 
\hline
\multicolumn{4}{|p{16cm}|}{Delivery Mode:
[X] in person \hspace{0.6cm}
[ ] blended \hspace{0.6cm}
[ ] distance learning}\\
\hline
\multicolumn{2}{|p{8cm}|}{Semester workload} &
\multicolumn{2}{p{8cm}|}{Weekly workload}\\
\hline
\multicolumn{1}{|p{4cm}|}{Total} &
\multicolumn{1}{p{4cm}|}{Extension} &
\multicolumn{1}{p{4cm}|}{Theoretical} &
\multicolumn{1}{p{4cm}|}{Practical} \\ 
\multicolumn{1}{|p{4cm}|}{60\,hours} &
\multicolumn{1}{p{4cm}|}{0\;hours} &
\multicolumn{1}{p{4cm}|}{4\;class hours} &
\multicolumn{1}{p{4cm}|}{0\;class hours} \\ 
\hline
\multicolumn{4}{|p{16cm}|}{Course Description:}\\
\multicolumn{4}{|p{16cm}|}{}\\
\multicolumn{4}{|p{\dimexpr 16cm + 6\tabcolsep\relax}|}{Introduction to data visualization. Data representation. Scalar data visualization. Vector data visualization. Tensor data visualization. Information visualization.}\\
\multicolumn{4}{|p{16cm}|}{}\\
\hline
\multicolumn{4}{|p{16cm}|}{Course Content:}\\*
\multicolumn{4}{|p{\dimexpr 16cm + 6\tabcolsep\relax}|}{\begin{enumerate}[start=1]\item Introduction to data visualization: definitions and types of visualization; visualization pipeline; computational tools\end{enumerate}}\\
\multicolumn{4}{|p{\dimexpr 16cm + 6\tabcolsep\relax}|}{\begin{enumerate}[start=2]\item Data representation: continuous data; data sampling and reconstruction; discrete data sets (cell, mesh, and attribute)\end{enumerate}}\\
\multicolumn{4}{|p{\dimexpr 16cm + 6\tabcolsep\relax}|}{\begin{enumerate}[start=3]\item Scalar data visualization: color mapping; contouring; height plots\end{enumerate}}\\
\multicolumn{4}{|p{\dimexpr 16cm + 6\tabcolsep\relax}|}{\begin{enumerate}[start=4]\item Vector data visualization: vector glyphs; color coding; stream objects (line and tube)\end{enumerate}}\\
\multicolumn{4}{|p{\dimexpr 16cm + 6\tabcolsep\relax}|}{\begin{enumerate}[start=5]\item Tensor data visualization: dimensionality reduction techniques; tensor glyphs; fiber tracking\end{enumerate}}\\
\multicolumn{4}{|p{\dimexpr 16cm + 6\tabcolsep\relax}|}{\begin{enumerate}[start=6]\item Information visualization: relational data visualization; multidimensional data visualization\end{enumerate}}\\
\multicolumn{4}{|p{\dimexpr 16cm + 6\tabcolsep\relax}|}{}\\
\hline
\multicolumn{4}{|p{16cm}|}{Basic Bibliography:}\\*
\multicolumn{4}{|p{\dimexpr 16cm + 6\tabcolsep\relax}|}{\begin{itemize}\item FOLEY, James D. Computer graphics: principles and practice in C. 2nd ed. Reading, Mass.: Addison Wesley, 1996. 1175 p.\end{itemize}}\\
\multicolumn{4}{|p{\dimexpr 16cm + 6\tabcolsep\relax}|}{\begin{itemize}\item GONZALEZ, Rafael C.; WOODS, Richard E. Digital image processing. 3. ed. Upper Saddle River, NJ: Pearson Prentice Hall, 2008. 954 p.\end{itemize}}\\
\multicolumn{4}{|p{\dimexpr 16cm + 6\tabcolsep\relax}|}{\begin{itemize}\item HAN, Jiawei; KAMBER, Micheline; PEI, Jian. Data mining: concepts and techniques. 3. ed. Boston: Elsevier, 2012. 703 p.\end{itemize}}\\
\multicolumn{4}{|p{\dimexpr 16cm + 6\tabcolsep\relax}|}{\begin{itemize}\item PEDRINI, Hélio; SCHWARTZ, William Robson. Análise de imagens digitais: princípios, algoritmos e aplicações. São Paulo: Thomson Learning, 2008. 508 p.\end{itemize}}\\
\multicolumn{4}{|p{\dimexpr 16cm + 6\tabcolsep\relax}|}{\begin{itemize}\item SCHROEDER, Will; MARTIN, Ken; LORENSE, Bill. The visualization toolkit. New Jersey: Prentice-Hall, 1996. 826 p.\end{itemize}}\\
\multicolumn{4}{|p{16cm}|}{}\\
\hline
\multicolumn{4}{|p{16cm}|}{Complementary Bibliography:}\\*
\multicolumn{4}{|p{\dimexpr 16cm + 6\tabcolsep\relax}|}{\begin{itemize}\item HAYKIN, Simon S. Redes neurais: princípios e prática. 2. ed. Porto Alegre: Bookman, 2001. 900 p.\end{itemize}}\\
\multicolumn{4}{|p{\dimexpr 16cm + 6\tabcolsep\relax}|}{\begin{itemize}\item KIRK, Andy. Data visualisation: a handbook for data driven design. 2. ed. London: Sage Publications Ltd, 2019. 328 p.\end{itemize}}\\
\multicolumn{4}{|p{\dimexpr 16cm + 6\tabcolsep\relax}|}{\begin{itemize}\item LIM, Tee T.; SMITS, Alexander J. Flow visualization: techniques and examples. 2. ed. London: Imperial College Press, 2012. 442 p.\end{itemize}}\\
\multicolumn{4}{|p{\dimexpr 16cm + 6\tabcolsep\relax}|}{\begin{itemize}\item MUNZNER, Tamara. Visualization analysis and design. 1. ed. New York: A K Peters/CRC Press, 2014. 428 p.\end{itemize}}\\
\multicolumn{4}{|p{\dimexpr 16cm + 6\tabcolsep\relax}|}{\begin{itemize}\item SPENCE, Robert. Information visualization: an introduction. 3. ed. London: Springer International Publishing, 2014. 346 p.\end{itemize}}\\
\multicolumn{4}{|p{\dimexpr 16cm + 6\tabcolsep\relax}|}{\begin{itemize}\item TELEA, Alexandru C. Data visualization: principles and practice. 2. ed. New York: A K Peters/CRC Press, 2014. 617 p.\end{itemize}}\\
\multicolumn{4}{|p{\dimexpr 16cm + 6\tabcolsep\relax}|}{\begin{itemize}\item TIERNY, Julien. Topological data analysis for scientific visualization. 1. ed. Switzerland: Springer International Publishing AG, 2017. 165 p.\end{itemize}}\\
\hline
\end{longtable}
\end{center}

\clearpage
\begin{center}
\begin{tabular}{lccr}
 \multirow{3}{*}{\includegraphics[height=2.7cm]{brasao.png}} &
 \multicolumn{2}{c}{\bfseries FEDERAL UNIVERSITY OF OURO PRETO} &
 \ \ \ \ \multirow{3}{*}{\includegraphics[height=2.7cm]{ufop.png}} \\
 & \multicolumn{2}{c}{\bfseries DEAN'S OFFICE FOR UNDERGRADUATE STUDIES} & \\
 & \multicolumn{2}{c}{} & \\
 & \multicolumn{2}{c}{\Large\bfseries COURSE PROGRAM} & \\
\end{tabular}
\end{center}

\begin{center}
\begin{longtable}{|p{4cm}|p{4cm}|p{4cm}|p{4cm}|}
\hline
\multicolumn{3}{|p{12cm}|}{Course Name in English:} &
\multicolumn{1}{p{4cm}|}{Code:} \\ 
\multicolumn{3}{|p{12cm}|}{\textbf{Numerical Methods for Differential Equations}} &
\textbf{BCC294}\\ 
\multicolumn{3}{|p{12cm}|}{Course Name in Portuguese:} & \\ 
\multicolumn{3}{|p{12cm}|}{\textbf{Métodos Numéricos para Equações Diferenciais}} & \\ 
\hline
\multicolumn{3}{|p{12cm}|}{Department Name and Acronym:} & Academic Unit: \\ 
\multicolumn{3}{|p{12cm}|}{Department of Computer Science (DECOM)} & {ICEB} \\ 
\hline
\multicolumn{4}{|p{16cm}|}{Delivery Mode:
[X] in person \hspace{0.6cm}
[ ] blended \hspace{0.6cm}
[ ] distance learning}\\
\hline
\multicolumn{2}{|p{8cm}|}{Semester workload} &
\multicolumn{2}{p{8cm}|}{Weekly workload}\\
\hline
\multicolumn{1}{|p{4cm}|}{Total} &
\multicolumn{1}{p{4cm}|}{Extension} &
\multicolumn{1}{p{4cm}|}{Theoretical} &
\multicolumn{1}{p{4cm}|}{Practical} \\ 
\multicolumn{1}{|p{4cm}|}{60\,hours} &
\multicolumn{1}{p{4cm}|}{0\;hours} &
\multicolumn{1}{p{4cm}|}{4\;class hours} &
\multicolumn{1}{p{4cm}|}{0\;class hours} \\ 
\hline
\multicolumn{4}{|p{16cm}|}{Course Description:}\\
\multicolumn{4}{|p{16cm}|}{}\\
\multicolumn{4}{|p{\dimexpr 16cm + 6\tabcolsep\relax}|}{Initial value problems for ordinary differential equations. Boundary value problems for ordinary differential equations. Numerical solutions for partial differential equations.}\\
\multicolumn{4}{|p{16cm}|}{}\\
\hline
\multicolumn{4}{|p{16cm}|}{Course Content:}\\*
\multicolumn{4}{|p{\dimexpr 16cm + 6\tabcolsep\relax}|}{\begin{enumerate}[start=1]\item Initial value problems for ordinary differential equations: introduction to initial value problems; Euler's method; higher-order Taylor methods; Runge-Kutta methods; multistep methods; variable-step-size multistep methods; higher-order equations and systems of differential equations; application examples\end{enumerate}}\\
\multicolumn{4}{|p{\dimexpr 16cm + 6\tabcolsep\relax}|}{\begin{enumerate}[start=2]\item Boundary value problems for ordinary differential equations: introduction to boundary value problems; linear shooting method; shooting method for nonlinear problems; finite difference method for linear problems; finite difference method for nonlinear problems; application examples\end{enumerate}}\\
\multicolumn{4}{|p{\dimexpr 16cm + 6\tabcolsep\relax}|}{\begin{enumerate}[start=3]\item Numerical solutions for partial differential equations: introduction to partial differential equations; elliptic partial differential equations; parabolic partial differential equations; hyperbolic partial differential equations; application examples\end{enumerate}}\\
\multicolumn{4}{|p{\dimexpr 16cm + 6\tabcolsep\relax}|}{}\\
\hline
\multicolumn{4}{|p{16cm}|}{Basic Bibliography:}\\*
\multicolumn{4}{|p{\dimexpr 16cm + 6\tabcolsep\relax}|}{\begin{itemize}\item BURDEN, R. L.; FAIRES, J. D.; BURDEN, A. M. Análise numérica. 10. ed. São Paulo: Cengage Learning, c2016. 879 p. ISBN 9788522123407.\end{itemize}}\\
\multicolumn{4}{|p{\dimexpr 16cm + 6\tabcolsep\relax}|}{\begin{itemize}\item CAMPOS FILHO, F. F. Algoritmos numéricos. Rio de Janeiro: LTC, 2001. xv, 383 p. ISBN 8521612656.\end{itemize}}\\
\multicolumn{4}{|p{\dimexpr 16cm + 6\tabcolsep\relax}|}{\begin{itemize}\item RUGGIERO, M. A. G.; LOPES, V. L. R. Cálculo numérico: aspectos teóricos e computacionais. 2. ed. São Paulo: Makron Books, c1997. xvi, 406 p. ISBN 8534602042.\end{itemize}}\\
\multicolumn{4}{|p{\dimexpr 16cm + 6\tabcolsep\relax}|}{\begin{itemize}\item FRANCO, N. B. Cálculo numérico. São Paulo: Pearson Prentice Hall, 2006. 505 p. ISBN 8576050870.\end{itemize}}\\
\multicolumn{4}{|p{\dimexpr 16cm + 6\tabcolsep\relax}|}{\begin{itemize}\item DOERING, C. I.; LOPES, A. O. Equações diferenciais ordinárias. 2. ed. Rio de Janeiro: IMPA, 2007. 423 p. ISBN 9788524402395.\end{itemize}}\\
\multicolumn{4}{|p{\dimexpr 16cm + 6\tabcolsep\relax}|}{\begin{itemize}\item SMITH, G. D. Numerical solution of partial differential equations: finite difference methods. 3. ed. Oxford: Clarendon, 1992. 337 p. ISBN 0198596502.\end{itemize}}\\
\multicolumn{4}{|p{16cm}|}{}\\
\hline
\multicolumn{4}{|p{16cm}|}{Complementary Bibliography:}\\*
\multicolumn{4}{|p{\dimexpr 16cm + 6\tabcolsep\relax}|}{\begin{itemize}\item QUARTERONI, A.; SACCO, R.; SALERI, F. Numerical Mathematics. 2. ed. New York: Springer, 2010. 657 p. ISBN 3642071015.\end{itemize}}\\
\multicolumn{4}{|p{\dimexpr 16cm + 6\tabcolsep\relax}|}{\begin{itemize}\item QUARTERONI, A.; VALLI, A. Numerical Approximation of Partial Differential Equations. 1. ed. New York: Springer, 2008. 562 p. ISBN 3540852670.\end{itemize}}\\
\multicolumn{4}{|p{\dimexpr 16cm + 6\tabcolsep\relax}|}{\begin{itemize}\item GOCKENBACH, M. S. Partial differential equations: analytical and numerical methods. Philadelphia: Society for Industrial and Applied Mathematics, c2002. xxii, 614 p. ISBN 0898715180.\end{itemize}}\\
\hline
\end{longtable}
\end{center}

\clearpage
\begin{center}
\begin{tabular}{lccr}
 \multirow{3}{*}{\includegraphics[height=2.7cm]{brasao.png}} &
 \multicolumn{2}{c}{\bfseries FEDERAL UNIVERSITY OF OURO PRETO} &
 \ \ \ \ \multirow{3}{*}{\includegraphics[height=2.7cm]{ufop.png}} \\
 & \multicolumn{2}{c}{\bfseries DEAN'S OFFICE FOR UNDERGRADUATE STUDIES} & \\
 & \multicolumn{2}{c}{} & \\
 & \multicolumn{2}{c}{\Large\bfseries COURSE PROGRAM} & \\
\end{tabular}
\end{center}

\begin{center}
\begin{longtable}{|p{4cm}|p{4cm}|p{4cm}|p{4cm}|}
\hline
\multicolumn{3}{|p{12cm}|}{Course Name in English:} &
\multicolumn{1}{p{4cm}|}{Code:} \\ 
\multicolumn{3}{|p{12cm}|}{\textbf{Logics in Computer Science}} &
\textbf{BCC295}\\ 
\multicolumn{3}{|p{12cm}|}{Course Name in Portuguese:} & \\ 
\multicolumn{3}{|p{12cm}|}{\textbf{Lógica Aplicada à Computação}} & \\ 
\hline
\multicolumn{3}{|p{12cm}|}{Department Name and Acronym:} & Academic Unit: \\ 
\multicolumn{3}{|p{12cm}|}{Department of Computer Science (DECOM)} & {ICEB} \\ 
\hline
\multicolumn{4}{|p{16cm}|}{Delivery Mode:
[X] in person \hspace{0.6cm}
[ ] blended \hspace{0.6cm}
[ ] distance learning}\\
\hline
\multicolumn{2}{|p{8cm}|}{Semester workload} &
\multicolumn{2}{p{8cm}|}{Weekly workload}\\
\hline
\multicolumn{1}{|p{4cm}|}{Total} &
\multicolumn{1}{p{4cm}|}{Extension} &
\multicolumn{1}{p{4cm}|}{Theoretical} &
\multicolumn{1}{p{4cm}|}{Practical} \\ 
\multicolumn{1}{|p{4cm}|}{60\,hours} &
\multicolumn{1}{p{4cm}|}{0\;hours} &
\multicolumn{1}{p{4cm}|}{4\;class hours} &
\multicolumn{1}{p{4cm}|}{0\;class hours} \\ 
\hline
\multicolumn{4}{|p{16cm}|}{Course Description:}\\
\multicolumn{4}{|p{16cm}|}{}\\
\multicolumn{4}{|p{\dimexpr 16cm + 6\tabcolsep\relax}|}{First-order logic, the DPLL algorithm, and the use of SAT and SMT solvers. Intuitionistic logic and type theory. Use of proof assistants. Induction and co-induction and their use in program verification. Temporal logic and model checking.}\\
\multicolumn{4}{|p{16cm}|}{}\\
\hline
\multicolumn{4}{|p{16cm}|}{Course Content:}\\*
\multicolumn{4}{|p{\dimexpr 16cm + 6\tabcolsep\relax}|}{\begin{enumerate}[start=1]\item Course presentation: motivation and evaluation criteria\end{enumerate}}\\
\multicolumn{4}{|p{\dimexpr 16cm + 6\tabcolsep\relax}|}{\begin{enumerate}[start=2]\item Review of propositional logic: syntax, semantics, and natural deduction\end{enumerate}}\\
\multicolumn{4}{|p{\dimexpr 16cm + 6\tabcolsep\relax}|}{\begin{enumerate}[start=3]\item Review of predicate logic: syntax, semantics, and natural deduction\end{enumerate}}\\
\multicolumn{4}{|p{\dimexpr 16cm + 6\tabcolsep\relax}|}{\begin{enumerate}[start=4]\item The SAT problem and the DPLL algorithm\end{enumerate}}\\
\multicolumn{4}{|p{\dimexpr 16cm + 6\tabcolsep\relax}|}{\begin{enumerate}[start=5]\item Modeling problems using SAT and using SAT solvers\end{enumerate}}\\
\multicolumn{4}{|p{\dimexpr 16cm + 6\tabcolsep\relax}|}{\begin{enumerate}[start=6]\item Satisfiability Modulo Theories (SMT) and the DPLL-T algorithm\end{enumerate}}\\
\multicolumn{4}{|p{\dimexpr 16cm + 6\tabcolsep\relax}|}{\begin{enumerate}[start=7]\item Modeling problems using SMT and SMT solvers\end{enumerate}}\\
\multicolumn{4}{|p{\dimexpr 16cm + 6\tabcolsep\relax}|}{\begin{enumerate}[start=8]\item Intuitionistic logic and typed lambda calculus: the Curry-Howard isomorphism\end{enumerate}}\\
\multicolumn{4}{|p{\dimexpr 16cm + 6\tabcolsep\relax}|}{\begin{enumerate}[start=9]\item Introduction to the use of type-theory-based proof assistants\end{enumerate}}\\
\multicolumn{4}{|p{\dimexpr 16cm + 6\tabcolsep\relax}|}{\begin{enumerate}[start=10]\item Review: mathematical induction, proofs by structural induction\end{enumerate}}\\
\multicolumn{4}{|p{\dimexpr 16cm + 6\tabcolsep\relax}|}{\begin{enumerate}[start=11]\item Induction over well-founded relations\end{enumerate}}\\
\multicolumn{4}{|p{\dimexpr 16cm + 6\tabcolsep\relax}|}{\begin{enumerate}[start=12]\item Proofs by co-induction; use of co-induction in proof assistants\end{enumerate}}\\
\multicolumn{4}{|p{\dimexpr 16cm + 6\tabcolsep\relax}|}{\begin{enumerate}[start=13]\item Temporal logic and model checking\end{enumerate}}\\
\multicolumn{4}{|p{\dimexpr 16cm + 6\tabcolsep\relax}|}{}\\
\hline
\multicolumn{4}{|p{16cm}|}{Basic Bibliography:}\\*
\multicolumn{4}{|p{\dimexpr 16cm + 6\tabcolsep\relax}|}{\begin{itemize}\item HUTH, Michael; RYAN, Mark. Lógica para ciência da computação: modelagem e argumentação sobre sistemas. 2. ed. LTC, 2008.\end{itemize}}\\
\multicolumn{4}{|p{\dimexpr 16cm + 6\tabcolsep\relax}|}{\begin{itemize}\item BERTOT, Yves; CASTÉRAN, Pierre. Interactive Theorem Proving and Program Development: The Coq'Art – The calculus of inductive constructions. 1. ed. Springer-Verlag, 2002.\end{itemize}}\\
\multicolumn{4}{|p{\dimexpr 16cm + 6\tabcolsep\relax}|}{\begin{itemize}\item MITCHELL, John. Foundations for Programming Languages. MIT Press, 1996.\end{itemize}}\\
\multicolumn{4}{|p{16cm}|}{}\\
\hline
\multicolumn{4}{|p{16cm}|}{Complementary Bibliography:}\\*
\multicolumn{4}{|p{\dimexpr 16cm + 6\tabcolsep\relax}|}{\begin{itemize}\item BIERE, Armin; HEULE, Martin. Handbook of Satisfiability: Volume 185, Frontiers in Artificial Intelligence and Applications. 1. ed. IOS Press.\end{itemize}}\\
\multicolumn{4}{|p{\dimexpr 16cm + 6\tabcolsep\relax}|}{\begin{itemize}\item CLARKE, Edmund; GRUMBERG, Orna. Model Checking. 1. ed. MIT Press, 1999.\end{itemize}}\\
\multicolumn{4}{|p{\dimexpr 16cm + 6\tabcolsep\relax}|}{\begin{itemize}\item BAIER, Christel; KATOEN, Joost-Pieter. Principles of Model Checking. 1. ed. MIT Press, 2008.\end{itemize}}\\
\multicolumn{4}{|p{\dimexpr 16cm + 6\tabcolsep\relax}|}{\begin{itemize}\item SANGIORGI, Davide. Introduction to Bisimulation and Coinduction. 1. ed. Cambridge University Press, 2011.\end{itemize}}\\
\hline
\end{longtable}
\end{center}

\clearpage
\begin{center}
\begin{tabular}{lccr}
 \multirow{3}{*}{\includegraphics[height=2.7cm]{brasao.png}} &
 \multicolumn{2}{c}{\bfseries FEDERAL UNIVERSITY OF OURO PRETO} &
 \ \ \ \ \multirow{3}{*}{\includegraphics[height=2.7cm]{ufop.png}} \\
 & \multicolumn{2}{c}{\bfseries DEAN'S OFFICE FOR UNDERGRADUATE STUDIES} & \\
 & \multicolumn{2}{c}{} & \\
 & \multicolumn{2}{c}{\Large\bfseries COURSE PROGRAM} & \\
\end{tabular}
\end{center}

\begin{center}
\begin{longtable}{|p{4cm}|p{4cm}|p{4cm}|p{4cm}|}
\hline
\multicolumn{3}{|p{12cm}|}{Course Name in English:} &
\multicolumn{1}{p{4cm}|}{Code:} \\ 
\multicolumn{3}{|p{12cm}|}{\textbf{Real-Time Software Design}} &
\textbf{BCC822}\\ 
\multicolumn{3}{|p{12cm}|}{Course Name in Portuguese:} & \\ 
\multicolumn{3}{|p{12cm}|}{\textbf{Projeto de Software de Tempo Real}} & \\ 
\hline
\multicolumn{3}{|p{12cm}|}{Department Name and Acronym:} & Academic Unit: \\ 
\multicolumn{3}{|p{12cm}|}{Department of Computer Science (DECOM)} & {ICEB} \\ 
\hline
\multicolumn{4}{|p{16cm}|}{Delivery Mode:
[X] in person \hspace{0.6cm}
[ ] blended \hspace{0.6cm}
[ ] distance learning}\\
\hline
\multicolumn{2}{|p{8cm}|}{Semester workload} &
\multicolumn{2}{p{8cm}|}{Weekly workload}\\
\hline
\multicolumn{1}{|p{4cm}|}{Total} &
\multicolumn{1}{p{4cm}|}{Extension} &
\multicolumn{1}{p{4cm}|}{Theoretical} &
\multicolumn{1}{p{4cm}|}{Practical} \\ 
\multicolumn{1}{|p{4cm}|}{60\,hours} &
\multicolumn{1}{p{4cm}|}{0\;hours} &
\multicolumn{1}{p{4cm}|}{2\;class hours} &
\multicolumn{1}{p{4cm}|}{2\;class hours} \\ 
\hline
\multicolumn{4}{|p{16cm}|}{Course Description:}\\
\multicolumn{4}{|p{16cm}|}{}\\
\multicolumn{4}{|p{\dimexpr 16cm + 6\tabcolsep\relax}|}{Introduction to Operating Systems. Basic Operating System concepts. Introduction to Real-Time Operating Systems. Classification of Operating Systems. Concurrency. Scheduling for Operating Systems. Scheduling for Real-Time Operating Systems. Synchronization objects and Critical Section. Inter-task Communication. Time Management. Formal methods for Real Time: modeling, languages, verification. Development of real-time systems.}\\
\multicolumn{4}{|p{16cm}|}{}\\
\hline
\multicolumn{4}{|p{16cm}|}{Course Content:}\\*
\multicolumn{4}{|p{\dimexpr 16cm + 6\tabcolsep\relax}|}{\begin{enumerate}[start=1]\item Operating System Concepts: characteristics of operating systems, processes, tasks, schedulers, scheduling policies, synchronization between tasks/processes, task and process states\end{enumerate}}\\
\multicolumn{4}{|p{\dimexpr 16cm + 6\tabcolsep\relax}|}{\begin{enumerate}[start=2]\item Introduction to Real-Time Operating Systems: definitions and characteristics; motivation for real-time operating systems; periodic and aperiodic tasks (arrival time, ready time, execution time, period, deadline); priorities; classification (Hard and Soft RTOS); volatile and static global variables; preemptive and non-preemptive; real-time interrupts and exceptions; real-time task states; task control block (TCB); performance measures\end{enumerate}}\\
\multicolumn{4}{|p{\dimexpr 16cm + 6\tabcolsep\relax}|}{\begin{enumerate}[start=3]\item Concurrent Systems: concepts and problems\end{enumerate}}\\
\multicolumn{4}{|p{\dimexpr 16cm + 6\tabcolsep\relax}|}{\begin{enumerate}[start=4]\item Real-Time System Scheduling Policies\end{enumerate}}\\
\multicolumn{4}{|p{\dimexpr 16cm + 6\tabcolsep\relax}|}{\begin{enumerate}[start=5]\item Synchronization Objects and Critical Section\end{enumerate}}\\
\multicolumn{4}{|p{\dimexpr 16cm + 6\tabcolsep\relax}|}{\begin{enumerate}[start=6]\item Task Synchronization: busy waiting, semaphores, mutexes, deadlock\end{enumerate}}\\
\multicolumn{4}{|p{\dimexpr 16cm + 6\tabcolsep\relax}|}{\begin{enumerate}[start=7]\item Inter-task Communication: shared memory, message passing\end{enumerate}}\\
\multicolumn{4}{|p{\dimexpr 16cm + 6\tabcolsep\relax}|}{\begin{enumerate}[start=8]\item Real-Time Communication: architecture, access methods, real-time protocols\end{enumerate}}\\
\multicolumn{4}{|p{\dimexpr 16cm + 6\tabcolsep\relax}|}{\begin{enumerate}[start=9]\item Formal Methods for Real Time: modeling, language, verification\end{enumerate}}\\
\multicolumn{4}{|p{\dimexpr 16cm + 6\tabcolsep\relax}|}{\begin{enumerate}[start=10]\item Development of real-time systems\end{enumerate}}\\
\multicolumn{4}{|p{\dimexpr 16cm + 6\tabcolsep\relax}|}{}\\
\hline
\multicolumn{4}{|p{16cm}|}{Basic Bibliography:}\\*
\multicolumn{4}{|p{\dimexpr 16cm + 6\tabcolsep\relax}|}{\begin{itemize}\item BERNARDIN, Gustavo Weber; BARRIQUELLO, Carlos Henrique. Sistemas operacionais de tempo real e sua aplicação em sistemas embarcados. Editora Blucher, 2019. ISBN 9788521213970. Disponível na Minha Biblioteca: https://plataforma.bvirtual.com.br/Leitor/Publicacao/169968/pdf/ (acessível via Minha UFOP – Biblioteca Digital).\end{itemize}}\\
\multicolumn{4}{|p{\dimexpr 16cm + 6\tabcolsep\relax}|}{\begin{itemize}\item LABROSSE, J. J. MicroC/OS-III, The Real-Time Kernel, or a High Performance, Scalable, ROMable, Preemptive, Multitasking Kernel for Microprocessors, Microcontrollers \& DSPs. Micrium Press, 2009. ISBN 978-0982337530 / 9788587918574. Disponível em: https://www.micrium.com/books/ucosiii/\end{itemize}}\\
\multicolumn{4}{|p{\dimexpr 16cm + 6\tabcolsep\relax}|}{\begin{itemize}\item TANENBAUM, Andrew S. Sistemas Operacionais Modernos. Editora Pearson, 2003. Disponível na Minha Biblioteca: https://plataforma.bvirtual.com.br/Leitor/Publicacao/36876/pdf/ (acessível via Minha UFOP – Biblioteca Digital).\end{itemize}}\\
\multicolumn{4}{|p{16cm}|}{}\\
\hline
\multicolumn{4}{|p{16cm}|}{Complementary Bibliography:}\\*
\multicolumn{4}{|p{\dimexpr 16cm + 6\tabcolsep\relax}|}{\begin{itemize}\item SILBERSCHATZ, Abraham; GAGNE, Greg; GALVIN, Peter B. Fundamentos de Sistemas Operacionais. John Wiley \& Sons, 2008. ISBN 9780470128725. Disponível na Minha Biblioteca: https://integrada.minhabiblioteca.com.br/\#/books/978-85-216-3001-2 (acessível via Minha UFOP – Biblioteca Digital).\end{itemize}}\\
\multicolumn{4}{|p{\dimexpr 16cm + 6\tabcolsep\relax}|}{\begin{itemize}\item DEITEL, H. M.; DEITEL, P. J.; CHOFFNES, D. R. Sistemas Operacionais. Editora Pearson, 2005. ISBN 9788576050117. Disponível na Minha Biblioteca: https://plataforma.bvirtual.com.br/Leitor/Loader/315/pdf (acessível via Minha UFOP – Biblioteca Digital Virtual Pearson).\end{itemize}}\\
\multicolumn{4}{|p{\dimexpr 16cm + 6\tabcolsep\relax}|}{\begin{itemize}\item TANENBAUM, Andrew S. Organização estruturada de computadores. 6. ed. Disponível na Biblioteca Pearson: https://plataforma.bvirtual.com.br/Acervo/Publicacao/3825 (acessível via Minha UFOP – Biblioteca Digital).\end{itemize}}\\
\multicolumn{4}{|p{\dimexpr 16cm + 6\tabcolsep\relax}|}{\begin{itemize}\item HALLINAN, Christopher. Embedded Linux Primer: A Practical Real-World Approach. Prentice-Hall, 2010. ISBN 9780137017836.\end{itemize}}\\
\multicolumn{4}{|p{\dimexpr 16cm + 6\tabcolsep\relax}|}{\begin{itemize}\item ABBOTT, Doug. Linux for Embedded and Real-time Applications. Newnes, 2006. ISBN 9780750679329.\end{itemize}}\\
\multicolumn{4}{|p{\dimexpr 16cm + 6\tabcolsep\relax}|}{\begin{itemize}\item BURNS, Alan; WELLINGS, Andy. Real-Time Systems \& Programming Languages. Addison-Wesley, 1. ed. ISBN 9780201729887.\end{itemize}}\\
\hline
\end{longtable}
\end{center}

\clearpage
\begin{center}
\begin{tabular}{lccr}
 \multirow{3}{*}{\includegraphics[height=2.7cm]{brasao.png}} &
 \multicolumn{2}{c}{\bfseries FEDERAL UNIVERSITY OF OURO PRETO} &
 \ \ \ \ \multirow{3}{*}{\includegraphics[height=2.7cm]{ufop.png}} \\
 & \multicolumn{2}{c}{\bfseries DEAN'S OFFICE FOR UNDERGRADUATE STUDIES} & \\
 & \multicolumn{2}{c}{} & \\
 & \multicolumn{2}{c}{\Large\bfseries COURSE PROGRAM} & \\
\end{tabular}
\end{center}

\begin{center}
\begin{longtable}{|p{4cm}|p{4cm}|p{4cm}|p{4cm}|}
\hline
\multicolumn{3}{|p{12cm}|}{Course Name in English:} &
\multicolumn{1}{p{4cm}|}{Code:} \\ 
\multicolumn{3}{|p{12cm}|}{\textbf{Data Structures II}} &
\textbf{BCC283}\\ 
\multicolumn{3}{|p{12cm}|}{Course Name in Portuguese:} & \\ 
\multicolumn{3}{|p{12cm}|}{\textbf{Estrutura de Dados II}} & \\ 
\hline
\multicolumn{3}{|p{12cm}|}{Department Name and Acronym:} & Academic Unit: \\ 
\multicolumn{3}{|p{12cm}|}{Department of Computer Science (DECOM)} & {ICEB} \\ 
\hline
\multicolumn{4}{|p{16cm}|}{Delivery Mode:
[X] in person \hspace{0.6cm}
[ ] blended \hspace{0.6cm}
[ ] distance learning}\\
\hline
\multicolumn{2}{|p{8cm}|}{Semester workload} &
\multicolumn{2}{p{8cm}|}{Weekly workload}\\
\hline
\multicolumn{1}{|p{4cm}|}{Total} &
\multicolumn{1}{p{4cm}|}{Extension} &
\multicolumn{1}{p{4cm}|}{Theoretical} &
\multicolumn{1}{p{4cm}|}{Practical} \\ 
\multicolumn{1}{|p{4cm}|}{60\,hours} &
\multicolumn{1}{p{4cm}|}{0\;hours} &
\multicolumn{1}{p{4cm}|}{4\;class hours} &
\multicolumn{1}{p{4cm}|}{0\;class hours} \\ 
\hline
\multicolumn{4}{|p{16cm}|}{Course Description:}\\
\multicolumn{4}{|p{16cm}|}{}\\
\multicolumn{4}{|p{\dimexpr 16cm + 6\tabcolsep\relax}|}{Secondary memory storage. Secondary memory sorting. Secondary memory searching. String matching. Text compression. Spatial data structures.}\\
\multicolumn{4}{|p{16cm}|}{}\\
\hline
\multicolumn{4}{|p{16cm}|}{Course Content:}\\*
\multicolumn{4}{|p{\dimexpr 16cm + 6\tabcolsep\relax}|}{\begin{enumerate}[start=1]\item Secondary memory storage: storage devices, sequential and random access\end{enumerate}}\\
\multicolumn{4}{|p{\dimexpr 16cm + 6\tabcolsep\relax}|}{\begin{enumerate}[start=2]\item Secondary memory sorting: balanced multiway merging, replacement selection, polyphase merging, external quicksort\end{enumerate}}\\
\multicolumn{4}{|p{\dimexpr 16cm + 6\tabcolsep\relax}|}{\begin{enumerate}[start=3]\item Secondary memory searching: indexed sequential access, B-tree, B*-tree\end{enumerate}}\\
\multicolumn{4}{|p{\dimexpr 16cm + 6\tabcolsep\relax}|}{\begin{enumerate}[start=4]\item String matching: exact string matching, approximate string matching\end{enumerate}}\\
\multicolumn{4}{|p{\dimexpr 16cm + 6\tabcolsep\relax}|}{\begin{enumerate}[start=5]\item Text compression: natural language text compression, Huffman coding\end{enumerate}}\\
\multicolumn{4}{|p{\dimexpr 16cm + 6\tabcolsep\relax}|}{\begin{enumerate}[start=6]\item Spatial data structures: quadtree, kD-tree, R-tree\end{enumerate}}\\
\multicolumn{4}{|p{\dimexpr 16cm + 6\tabcolsep\relax}|}{}\\
\hline
\multicolumn{4}{|p{16cm}|}{Basic Bibliography:}\\*
\multicolumn{4}{|p{\dimexpr 16cm + 6\tabcolsep\relax}|}{\begin{itemize}\item SEDGEWICK, Robert. Algorithms. 4. ed. Upper Saddle River: Addison Wesley, 2011.\end{itemize}}\\
\multicolumn{4}{|p{\dimexpr 16cm + 6\tabcolsep\relax}|}{\begin{itemize}\item ZIVIANI, Nivio. Projeto de algoritmos: com implementações em Pascal e C. 5. ed. São Paulo: Pioneira, 2000.\end{itemize}}\\
\multicolumn{4}{|p{\dimexpr 16cm + 6\tabcolsep\relax}|}{\begin{itemize}\item CORMEN, Thomas H. ... [et al.]; tradução Vandenberg D. de Souza. Algoritmos: teoria e prática. Rio de Janeiro: Campus, 2002.\end{itemize}}\\
\multicolumn{4}{|p{16cm}|}{}\\
\hline
\multicolumn{4}{|p{16cm}|}{Complementary Bibliography:}\\*
\multicolumn{4}{|p{\dimexpr 16cm + 6\tabcolsep\relax}|}{\begin{itemize}\item ZIVIANI, Nivio; BOTELHO, Fabiano C. Projeto de algoritmos: com implementações em Java e C++. São Paulo: Cengage Learning, 2015. (item inexistente na UFOP)\end{itemize}}\\
\multicolumn{4}{|p{\dimexpr 16cm + 6\tabcolsep\relax}|}{\begin{itemize}\item Bancos de dados geográficos. Curitiba: MundoGeo, 2005.\end{itemize}}\\
\multicolumn{4}{|p{\dimexpr 16cm + 6\tabcolsep\relax}|}{\begin{itemize}\item DEITEL, Harvey M. C++ how to program. Harvey M. Deitel, Paul J. Deitel. 6. ed. Boston: Prentice Hall, c2008.\end{itemize}}\\
\multicolumn{4}{|p{\dimexpr 16cm + 6\tabcolsep\relax}|}{\begin{itemize}\item MIZRAHI, Victorine Viviane. Treinamento em linguagem C. 2. ed. São Paulo: Pearson Education, c2010.\end{itemize}}\\
\multicolumn{4}{|p{\dimexpr 16cm + 6\tabcolsep\relax}|}{\begin{itemize}\item TENENBAUM, Aaron M. Estruturas de dados usando C. São Paulo: Makron Books, 1995.\end{itemize}}\\
\hline
\end{longtable}
\end{center}

\clearpage
\begin{center}
\begin{tabular}{lccr}
 \multirow{3}{*}{\includegraphics[height=2.7cm]{brasao.png}} &
 \multicolumn{2}{c}{\bfseries FEDERAL UNIVERSITY OF OURO PRETO} &
 \ \ \ \ \multirow{3}{*}{\includegraphics[height=2.7cm]{ufop.png}} \\
 & \multicolumn{2}{c}{\bfseries DEAN'S OFFICE FOR UNDERGRADUATE STUDIES} & \\
 & \multicolumn{2}{c}{} & \\
 & \multicolumn{2}{c}{\Large\bfseries COURSE PROGRAM} & \\
\end{tabular}
\end{center}

\begin{center}
\begin{longtable}{|p{4cm}|p{4cm}|p{4cm}|p{4cm}|}
\hline
\multicolumn{3}{|p{12cm}|}{Course Name in English:} &
\multicolumn{1}{p{4cm}|}{Code:} \\ 
\multicolumn{3}{|p{12cm}|}{\textbf{Theory of Computation}} &
\textbf{BCC884}\\ 
\multicolumn{3}{|p{12cm}|}{Course Name in Portuguese:} & \\ 
\multicolumn{3}{|p{12cm}|}{\textbf{Teoria da Computação}} & \\ 
\hline
\multicolumn{3}{|p{12cm}|}{Department Name and Acronym:} & Academic Unit: \\ 
\multicolumn{3}{|p{12cm}|}{Department of Computer Science (DECOM)} & {ICEB} \\ 
\hline
\multicolumn{4}{|p{16cm}|}{Delivery Mode:
[X] in person \hspace{0.6cm}
[ ] blended \hspace{0.6cm}
[ ] distance learning}\\
\hline
\multicolumn{2}{|p{8cm}|}{Semester workload} &
\multicolumn{2}{p{8cm}|}{Weekly workload}\\
\hline
\multicolumn{1}{|p{4cm}|}{Total} &
\multicolumn{1}{p{4cm}|}{Extension} &
\multicolumn{1}{p{4cm}|}{Theoretical} &
\multicolumn{1}{p{4cm}|}{Practical} \\ 
\multicolumn{1}{|p{4cm}|}{60\,hours} &
\multicolumn{1}{p{4cm}|}{0\;hours} &
\multicolumn{1}{p{4cm}|}{4\;class hours} &
\multicolumn{1}{p{4cm}|}{0\;class hours} \\ 
\hline
\multicolumn{4}{|p{16cm}|}{Course Description:}\\
\multicolumn{4}{|p{16cm}|}{}\\
\multicolumn{4}{|p{\dimexpr 16cm + 6\tabcolsep\relax}|}{Regular languages, regular expressions, finite-state automata. Context-free languages and grammars, pushdown automata. Context-sensitive languages and grammars. Turing machines. The Church-Turing thesis. Computability and decidability. The Chomsky hierarchy.}\\
\multicolumn{4}{|p{16cm}|}{}\\
\hline
\multicolumn{4}{|p{16cm}|}{Course Content:}\\*
\multicolumn{4}{|p{\dimexpr 16cm + 6\tabcolsep\relax}|}{\begin{enumerate}[start=1]\item Introduction: alphabets, strings, and languages\end{enumerate}}\\
\multicolumn{4}{|p{\dimexpr 16cm + 6\tabcolsep\relax}|}{\begin{enumerate}[start=2]\item Deterministic and Nondeterministic Finite-State Automata\end{enumerate}}\\
\multicolumn{4}{|p{\dimexpr 16cm + 6\tabcolsep\relax}|}{\begin{enumerate}[start=3]\item Regular Expressions\end{enumerate}}\\
\multicolumn{4}{|p{\dimexpr 16cm + 6\tabcolsep\relax}|}{\begin{enumerate}[start=4]\item Minimization of Finite Automata\end{enumerate}}\\
\multicolumn{4}{|p{\dimexpr 16cm + 6\tabcolsep\relax}|}{\begin{enumerate}[start=5]\item Properties of Regular Languages\end{enumerate}}\\
\multicolumn{4}{|p{\dimexpr 16cm + 6\tabcolsep\relax}|}{\begin{enumerate}[start=6]\item Pumping Lemma for Regular Languages (RLs)\end{enumerate}}\\
\multicolumn{4}{|p{\dimexpr 16cm + 6\tabcolsep\relax}|}{\begin{enumerate}[start=7]\item Context-Free Grammars and Languages (CFLs)\end{enumerate}}\\
\multicolumn{4}{|p{\dimexpr 16cm + 6\tabcolsep\relax}|}{\begin{enumerate}[start=8]\item Ambiguity\end{enumerate}}\\
\multicolumn{4}{|p{\dimexpr 16cm + 6\tabcolsep\relax}|}{\begin{enumerate}[start=9]\item Properties of CFLs\end{enumerate}}\\
\multicolumn{4}{|p{\dimexpr 16cm + 6\tabcolsep\relax}|}{\begin{enumerate}[start=10]\item Pushdown Automata\end{enumerate}}\\
\multicolumn{4}{|p{\dimexpr 16cm + 6\tabcolsep\relax}|}{\begin{enumerate}[start=11]\item Chomsky Normal Form\end{enumerate}}\\
\multicolumn{4}{|p{\dimexpr 16cm + 6\tabcolsep\relax}|}{\begin{enumerate}[start=12]\item Regular Grammars and Context-Sensitive Grammars\end{enumerate}}\\
\multicolumn{4}{|p{\dimexpr 16cm + 6\tabcolsep\relax}|}{\begin{enumerate}[start=13]\item Pumping Lemma for CFLs\end{enumerate}}\\
\multicolumn{4}{|p{\dimexpr 16cm + 6\tabcolsep\relax}|}{\begin{enumerate}[start=14]\item Turing Machines\end{enumerate}}\\
\multicolumn{4}{|p{\dimexpr 16cm + 6\tabcolsep\relax}|}{\begin{enumerate}[start=15]\item The Church-Turing Thesis\end{enumerate}}\\
\multicolumn{4}{|p{\dimexpr 16cm + 6\tabcolsep\relax}|}{\begin{enumerate}[start=16]\item Decision Problems\end{enumerate}}\\
\multicolumn{4}{|p{\dimexpr 16cm + 6\tabcolsep\relax}|}{\begin{enumerate}[start=17]\item Undecidability of the Halting Problem\end{enumerate}}\\
\multicolumn{4}{|p{\dimexpr 16cm + 6\tabcolsep\relax}|}{\begin{enumerate}[start=18]\item Decidable and undecidable problems about languages\end{enumerate}}\\
\multicolumn{4}{|p{\dimexpr 16cm + 6\tabcolsep\relax}|}{}\\
\hline
\multicolumn{4}{|p{16cm}|}{Basic Bibliography:}\\*
\multicolumn{4}{|p{\dimexpr 16cm + 6\tabcolsep\relax}|}{\begin{itemize}\item VIEIRA, Newton José. Introdução aos fundamentos da computação: linguagens e máquinas. São Paulo: Thomson Pioneira, 2006.\end{itemize}}\\
\multicolumn{4}{|p{\dimexpr 16cm + 6\tabcolsep\relax}|}{\begin{itemize}\item SIPSER, Michael. Introdução à teoria da computação. São Paulo: Thomson Learning, 2007.\end{itemize}}\\
\multicolumn{4}{|p{\dimexpr 16cm + 6\tabcolsep\relax}|}{\begin{itemize}\item SUDKAMP, Thomas A. Languages and machines: an introduction to the theory of computer science. 2. ed. Reading: Addison Wesley, c1997.\end{itemize}}\\
\multicolumn{4}{|p{16cm}|}{}\\
\hline
\multicolumn{4}{|p{16cm}|}{Complementary Bibliography:}\\*
\multicolumn{4}{|p{\dimexpr 16cm + 6\tabcolsep\relax}|}{\begin{itemize}\item MENDELSON, Elliott. Introduction to mathematical logic. 3rd ed. Pacific Grove, CA: Wadsworth \& Brooks/Cole Advanced Books \& Software, 1987.\end{itemize}}\\
\multicolumn{4}{|p{\dimexpr 16cm + 6\tabcolsep\relax}|}{\begin{itemize}\item LEWIS, Harry R. Elementos de teoria da computação. 2. ed. Porto Alegre: Bookman, 2000.\end{itemize}}\\
\multicolumn{4}{|p{\dimexpr 16cm + 6\tabcolsep\relax}|}{\begin{itemize}\item EPSTEIN, Richard L. Computability: computable functions, logic and the foundations of mathematics. Pacific Grove: Wadsworth, c1989.\end{itemize}}\\
\multicolumn{4}{|p{\dimexpr 16cm + 6\tabcolsep\relax}|}{\begin{itemize}\item DAVIS, Martin. Computability and unsolvability. New York: Dover, c1982.\end{itemize}}\\
\multicolumn{4}{|p{\dimexpr 16cm + 6\tabcolsep\relax}|}{\begin{itemize}\item Handbook of theoretical computer science. Amsterdã: Elsevier.\end{itemize}}\\
\hline
\end{longtable}
\end{center}

\clearpage
\begin{center}
\begin{tabular}{lccr}
 \multirow{3}{*}{\includegraphics[height=2.7cm]{brasao.png}} &
 \multicolumn{2}{c}{\bfseries FEDERAL UNIVERSITY OF OURO PRETO} &
 \ \ \ \ \multirow{3}{*}{\includegraphics[height=2.7cm]{ufop.png}} \\
 & \multicolumn{2}{c}{\bfseries DEAN'S OFFICE FOR UNDERGRADUATE STUDIES} & \\
 & \multicolumn{2}{c}{} & \\
 & \multicolumn{2}{c}{\Large\bfseries COURSE PROGRAM} & \\
\end{tabular}
\end{center}

\begin{center}
\begin{longtable}{|p{4cm}|p{4cm}|p{4cm}|p{4cm}|}
\hline
\multicolumn{3}{|p{12cm}|}{Course Name in English:} &
\multicolumn{1}{p{4cm}|}{Code:} \\ 
\multicolumn{3}{|p{12cm}|}{\textbf{Operating Systems}} &
\textbf{BCC344}\\ 
\multicolumn{3}{|p{12cm}|}{Course Name in Portuguese:} & \\ 
\multicolumn{3}{|p{12cm}|}{\textbf{Sistemas Operacionais}} & \\ 
\hline
\multicolumn{3}{|p{12cm}|}{Department Name and Acronym:} & Academic Unit: \\ 
\multicolumn{3}{|p{12cm}|}{Department of Computer Science (DECOM)} & {ICEB} \\ 
\hline
\multicolumn{4}{|p{16cm}|}{Delivery Mode:
[X] in person \hspace{0.6cm}
[ ] blended \hspace{0.6cm}
[ ] distance learning}\\
\hline
\multicolumn{2}{|p{8cm}|}{Semester workload} &
\multicolumn{2}{p{8cm}|}{Weekly workload}\\
\hline
\multicolumn{1}{|p{4cm}|}{Total} &
\multicolumn{1}{p{4cm}|}{Extension} &
\multicolumn{1}{p{4cm}|}{Theoretical} &
\multicolumn{1}{p{4cm}|}{Practical} \\ 
\multicolumn{1}{|p{4cm}|}{60\,hours} &
\multicolumn{1}{p{4cm}|}{0\;hours} &
\multicolumn{1}{p{4cm}|}{4\;class hours} &
\multicolumn{1}{p{4cm}|}{0\;class hours} \\ 
\hline
\multicolumn{4}{|p{16cm}|}{Course Description:}\\
\multicolumn{4}{|p{16cm}|}{}\\
\multicolumn{4}{|p{\dimexpr 16cm + 6\tabcolsep\relax}|}{Overview and history of operating systems; operating system principles; process management; use of thread APIs (Application Programming Interface); memory management; device management; security and protection; file systems.}\\
\multicolumn{4}{|p{16cm}|}{}\\
\hline
\multicolumn{4}{|p{16cm}|}{Course Content:}\\*
\multicolumn{4}{|p{\dimexpr 16cm + 6\tabcolsep\relax}|}{\begin{enumerate}[start=1]\item Introduction: definition, history, basic concepts, system calls, and structure\end{enumerate}}\\
\multicolumn{4}{|p{\dimexpr 16cm + 6\tabcolsep\relax}|}{\begin{enumerate}[start=2]\item Processes and Threads: inter-process communication (IPC), scheduling, and deadlocks\end{enumerate}}\\
\multicolumn{4}{|p{\dimexpr 16cm + 6\tabcolsep\relax}|}{\begin{enumerate}[start=3]\item Memory management: swapping, virtual memory, replacement algorithms, and segmentation\end{enumerate}}\\
\multicolumn{4}{|p{\dimexpr 16cm + 6\tabcolsep\relax}|}{\begin{enumerate}[start=4]\item Input/Output: disks, clocks, networks, and terminals\end{enumerate}}\\
\multicolumn{4}{|p{\dimexpr 16cm + 6\tabcolsep\relax}|}{\begin{enumerate}[start=5]\item File Systems: directories, file system implementation, and examples\end{enumerate}}\\
\multicolumn{4}{|p{\dimexpr 16cm + 6\tabcolsep\relax}|}{\begin{enumerate}[start=6]\item Network security: basic concepts of cryptography, authentication, attacks, and protection mechanisms\end{enumerate}}\\
\multicolumn{4}{|p{\dimexpr 16cm + 6\tabcolsep\relax}|}{\begin{enumerate}[start=7]\item Case Studies: overview, processes, memory management, I/O, file system, and security\end{enumerate}}\\
\multicolumn{4}{|p{\dimexpr 16cm + 6\tabcolsep\relax}|}{}\\
\hline
\multicolumn{4}{|p{16cm}|}{Basic Bibliography:}\\*
\multicolumn{4}{|p{\dimexpr 16cm + 6\tabcolsep\relax}|}{\begin{itemize}\item SILBERSCHATZ, Abraham. Fundamentos de sistemas operacionais. 8. ed. Rio de Janeiro: LTC, c2010.\end{itemize}}\\
\multicolumn{4}{|p{\dimexpr 16cm + 6\tabcolsep\relax}|}{\begin{itemize}\item SILBERSCHATZ, Abraham; GALVIN, Peter Baer; tradução e revisão técnica Carlos Camarão de Figueiredo. Sistemas operacionais: conceitos. São Paulo: Prentice-Hall, 2000.\end{itemize}}\\
\multicolumn{4}{|p{\dimexpr 16cm + 6\tabcolsep\relax}|}{\begin{itemize}\item TANENBAUM, Andrew S. Sistemas operacionais modernos. 2. ed. São Paulo: Prentice-Hall, 2003.\end{itemize}}\\
\multicolumn{4}{|p{16cm}|}{}\\
\hline
\multicolumn{4}{|p{16cm}|}{Complementary Bibliography:}\\*
\multicolumn{4}{|p{\dimexpr 16cm + 6\tabcolsep\relax}|}{\begin{itemize}\item BACH, Maurice J. The design of the UNIX operating system. Englewood Cliffs, N.J.: Prentice-Hall, c1986.\end{itemize}}\\
\multicolumn{4}{|p{\dimexpr 16cm + 6\tabcolsep\relax}|}{\begin{itemize}\item STEVENS, W. Richard. Advanced programming in the UNIX environment. Reading, MA: Addison Wesley, c1992.\end{itemize}}\\
\multicolumn{4}{|p{\dimexpr 16cm + 6\tabcolsep\relax}|}{\begin{itemize}\item MACHADO, Francis B. Arquitetura de sistemas operacionais. 3. ed. Rio de Janeiro: LTC, 2002.\end{itemize}}\\
\multicolumn{4}{|p{\dimexpr 16cm + 6\tabcolsep\relax}|}{\begin{itemize}\item CÔRTES, Pedro Luiz. Sistemas operacionais: fundamentos. São Paulo: Érica, 2003.\end{itemize}}\\
\multicolumn{4}{|p{\dimexpr 16cm + 6\tabcolsep\relax}|}{\begin{itemize}\item DEITEL, Harvey M. An Introduction to operating systems. 2nd ed. Reading: Addison Wesley, 1990.\end{itemize}}\\
\hline
\end{longtable}
\end{center}

\clearpage
\begin{center}
\begin{tabular}{lccr}
 \multirow{3}{*}{\includegraphics[height=2.7cm]{brasao.png}} &
 \multicolumn{2}{c}{\bfseries FEDERAL UNIVERSITY OF OURO PRETO} &
 \ \ \ \ \multirow{3}{*}{\includegraphics[height=2.7cm]{ufop.png}} \\
 & \multicolumn{2}{c}{\bfseries DEAN'S OFFICE FOR UNDERGRADUATE STUDIES} & \\
 & \multicolumn{2}{c}{} & \\
 & \multicolumn{2}{c}{\Large\bfseries COURSE PROGRAM} & \\
\end{tabular}
\end{center}

\begin{center}
\begin{longtable}{|p{4cm}|p{4cm}|p{4cm}|p{4cm}|}
\hline
\multicolumn{3}{|p{12cm}|}{Course Name in English:} &
\multicolumn{1}{p{4cm}|}{Code:} \\ 
\multicolumn{3}{|p{12cm}|}{\textbf{Human-Computer Interaction}} &
\textbf{BCC885}\\ 
\multicolumn{3}{|p{12cm}|}{Course Name in Portuguese:} & \\ 
\multicolumn{3}{|p{12cm}|}{\textbf{Interação Humano-Computador}} & \\ 
\hline
\multicolumn{3}{|p{12cm}|}{Department Name and Acronym:} & Academic Unit: \\ 
\multicolumn{3}{|p{12cm}|}{Department of Computer Science (DECOM)} & {ICEB} \\ 
\hline
\multicolumn{4}{|p{16cm}|}{Delivery Mode:
[X] in person \hspace{0.6cm}
[ ] blended \hspace{0.6cm}
[ ] distance learning}\\
\hline
\multicolumn{2}{|p{8cm}|}{Semester workload} &
\multicolumn{2}{p{8cm}|}{Weekly workload}\\
\hline
\multicolumn{1}{|p{4cm}|}{Total} &
\multicolumn{1}{p{4cm}|}{Extension} &
\multicolumn{1}{p{4cm}|}{Theoretical} &
\multicolumn{1}{p{4cm}|}{Practical} \\ 
\multicolumn{1}{|p{4cm}|}{60\,hours} &
\multicolumn{1}{p{4cm}|}{0\;hours} &
\multicolumn{1}{p{4cm}|}{4\;class hours} &
\multicolumn{1}{p{4cm}|}{0\;class hours} \\ 
\hline
\multicolumn{4}{|p{16cm}|}{Course Description:}\\
\multicolumn{4}{|p{16cm}|}{}\\
\multicolumn{4}{|p{\dimexpr 16cm + 6\tabcolsep\relax}|}{User interface and interaction. Quality in HCI. Theoretical approaches in HCI. Introduction to the interactive systems design process. Introduction to the interactive systems evaluation process.}\\
\multicolumn{4}{|p{16cm}|}{}\\
\hline
\multicolumn{4}{|p{16cm}|}{Course Content:}\\*
\multicolumn{4}{|p{\dimexpr 16cm + 6\tabcolsep\relax}|}{\begin{enumerate}[start=1]\item User Interface and Interaction: concept of interface and interaction, importance of the HCI field in software development, history, HCI as a multidisciplinary field, benefits of HCI, interaction problems\end{enumerate}}\\
\multicolumn{4}{|p{\dimexpr 16cm + 6\tabcolsep\relax}|}{\begin{enumerate}[start=2]\item Quality in HCI: usability, communicability, accessibility, user experience\end{enumerate}}\\
\multicolumn{4}{|p{\dimexpr 16cm + 6\tabcolsep\relax}|}{\begin{enumerate}[start=3]\item Theoretical Approaches in HCI: experimental psychology (Hick-Hyman Law and Fitts's Law), cognitive engineering, semiotic engineering\end{enumerate}}\\
\multicolumn{4}{|p{\dimexpr 16cm + 6\tabcolsep\relax}|}{\begin{enumerate}[start=4]\item Introduction to the Interactive Systems Design Process: task modeling, interaction modeling, prototyping, principles and guidelines for HCI design, interaction styles\end{enumerate}}\\
\multicolumn{4}{|p{\dimexpr 16cm + 6\tabcolsep\relax}|}{\begin{enumerate}[start=5]\item Introduction to the Interactive Systems Evaluation Process: HCI evaluation planning, inspection-based evaluation, user testing\end{enumerate}}\\
\multicolumn{4}{|p{\dimexpr 16cm + 6\tabcolsep\relax}|}{}\\
\hline
\multicolumn{4}{|p{16cm}|}{Basic Bibliography:}\\*
\multicolumn{4}{|p{\dimexpr 16cm + 6\tabcolsep\relax}|}{\begin{itemize}\item PREECE, Jennifer; ROGERS, Yvonne; SHARP, Helen; tradução de Viviane Possamai. Design de interação: além da interação homem-computador. Porto Alegre: Bookman, 2005.\end{itemize}}\\
\multicolumn{4}{|p{\dimexpr 16cm + 6\tabcolsep\relax}|}{\begin{itemize}\item BARBOSA, Simone Diniz Junqueira. Interação humano-computador. Rio de Janeiro: Campus, 2010.\end{itemize}}\\
\multicolumn{4}{|p{\dimexpr 16cm + 6\tabcolsep\relax}|}{\begin{itemize}\item SOUZA, Clarisse da Silva Vieira Camelo de. The semiotic engineering of human-computer interaction. Cambridge, Mass.: MIT Press, 2005.\end{itemize}}\\
\multicolumn{4}{|p{16cm}|}{}\\
\hline
\multicolumn{4}{|p{16cm}|}{Complementary Bibliography:}\\*
\multicolumn{4}{|p{\dimexpr 16cm + 6\tabcolsep\relax}|}{\begin{itemize}\item PREECE, Jenny. Human-computer interaction. Harlow: Addison Wesley, c1994.\end{itemize}}\\
\multicolumn{4}{|p{\dimexpr 16cm + 6\tabcolsep\relax}|}{\begin{itemize}\item CYBIS, Walter; BETIOL, Adriana Holtz; FAUST, Richard. Ergonomia e usabilidade: conhecimentos, métodos e aplicações. 2. ed. São Paulo: Novatec, 2007.\end{itemize}}\\
\multicolumn{4}{|p{\dimexpr 16cm + 6\tabcolsep\relax}|}{\begin{itemize}\item BENYON, David. Interação humano-computador. 2. ed. São Paulo: Pearson, 2011.\end{itemize}}\\
\multicolumn{4}{|p{\dimexpr 16cm + 6\tabcolsep\relax}|}{\begin{itemize}\item SHNEIDERMAN, Ben; PLAISANT, Catherine. Designing the user interface: strategies for effective human-computer interaction. 5th ed. Boston: Addison Wesley, 2010.\end{itemize}}\\
\multicolumn{4}{|p{\dimexpr 16cm + 6\tabcolsep\relax}|}{\begin{itemize}\item JACKO, Julie A. (ed.). The human-computer interaction handbook: fundamentals, evolving technologies, and emerging applications. 3rd ed. Boca Raton, FL: CRC Press, c2012.\end{itemize}}\\
\hline
\end{longtable}
\end{center}

\clearpage
\begin{center}
\begin{tabular}{lccr}
 \multirow{3}{*}{\includegraphics[height=2.7cm]{brasao.png}} &
 \multicolumn{2}{c}{\bfseries FEDERAL UNIVERSITY OF OURO PRETO} &
 \ \ \ \ \multirow{3}{*}{\includegraphics[height=2.7cm]{ufop.png}} \\
 & \multicolumn{2}{c}{\bfseries DEAN'S OFFICE FOR UNDERGRADUATE STUDIES} & \\
 & \multicolumn{2}{c}{} & \\
 & \multicolumn{2}{c}{\Large\bfseries COURSE PROGRAM} & \\
\end{tabular}
\end{center}

\begin{center}
\begin{longtable}{|p{4cm}|p{4cm}|p{4cm}|p{4cm}|}
\hline
\multicolumn{3}{|p{12cm}|}{Course Name in English:} &
\multicolumn{1}{p{4cm}|}{Code:} \\ 
\multicolumn{3}{|p{12cm}|}{\textbf{Computer Graphics}} &
\textbf{BCC887}\\ 
\multicolumn{3}{|p{12cm}|}{Course Name in Portuguese:} & \\ 
\multicolumn{3}{|p{12cm}|}{\textbf{Computação Gráfica}} & \\ 
\hline
\multicolumn{3}{|p{12cm}|}{Department Name and Acronym:} & Academic Unit: \\ 
\multicolumn{3}{|p{12cm}|}{Department of Computer Science (DECOM)} & {ICEB} \\ 
\hline
\multicolumn{4}{|p{16cm}|}{Delivery Mode:
[X] in person \hspace{0.6cm}
[ ] blended \hspace{0.6cm}
[ ] distance learning}\\
\hline
\multicolumn{2}{|p{8cm}|}{Semester workload} &
\multicolumn{2}{p{8cm}|}{Weekly workload}\\
\hline
\multicolumn{1}{|p{4cm}|}{Total} &
\multicolumn{1}{p{4cm}|}{Extension} &
\multicolumn{1}{p{4cm}|}{Theoretical} &
\multicolumn{1}{p{4cm}|}{Practical} \\ 
\multicolumn{1}{|p{4cm}|}{60\,hours} &
\multicolumn{1}{p{4cm}|}{0\;hours} &
\multicolumn{1}{p{4cm}|}{4\;class hours} &
\multicolumn{1}{p{4cm}|}{0\;class hours} \\ 
\hline
\multicolumn{4}{|p{16cm}|}{Course Description:}\\
\multicolumn{4}{|p{16cm}|}{}\\
\multicolumn{4}{|p{\dimexpr 16cm + 6\tabcolsep\relax}|}{Graphics systems; interactive graphics programming; modeling of three-dimensional objects and scenes; geometric transformations; viewing; lighting and shading; transparency and opacity; texture mapping; clipping; hidden surface removal; rasterization.}\\
\multicolumn{4}{|p{16cm}|}{}\\
\hline
\multicolumn{4}{|p{16cm}|}{Course Content:}\\*
\multicolumn{4}{|p{\dimexpr 16cm + 6\tabcolsep\relax}|}{\begin{enumerate}[start=1]\item Graphics systems\end{enumerate}}\\
\multicolumn{4}{|p{\dimexpr 16cm + 6\tabcolsep\relax}|}{\begin{enumerate}[start=2]\item Input and output devices\end{enumerate}}\\
\multicolumn{4}{|p{\dimexpr 16cm + 6\tabcolsep\relax}|}{\begin{enumerate}[start=3]\item Graphics system architecture\end{enumerate}}\\
\multicolumn{4}{|p{\dimexpr 16cm + 6\tabcolsep\relax}|}{\begin{enumerate}[start=4]\item Color\end{enumerate}}\\
\multicolumn{4}{|p{\dimexpr 16cm + 6\tabcolsep\relax}|}{\begin{enumerate}[start=5]\item Graphics programming: application programmer interface, primitives and attributes, interactive techniques\end{enumerate}}\\
\multicolumn{4}{|p{\dimexpr 16cm + 6\tabcolsep\relax}|}{\begin{enumerate}[start=6]\item Modeling of three-dimensional objects and scenes: elementary geometric objects, operations on geometric data, three-dimensional primitives, coordinate systems, object modeling\end{enumerate}}\\
\multicolumn{4}{|p{\dimexpr 16cm + 6\tabcolsep\relax}|}{\begin{enumerate}[start=7]\item Data structures\end{enumerate}}\\
\multicolumn{4}{|p{\dimexpr 16cm + 6\tabcolsep\relax}|}{\begin{enumerate}[start=8]\item Affine transformations\end{enumerate}}\\
\multicolumn{4}{|p{\dimexpr 16cm + 6\tabcolsep\relax}|}{\begin{enumerate}[start=9]\item Object movement\end{enumerate}}\\
\multicolumn{4}{|p{\dimexpr 16cm + 6\tabcolsep\relax}|}{\begin{enumerate}[start=10]\item Viewing: positioning the observer relative to the scene, viewing transformation, projection types, projection transformation\end{enumerate}}\\
\multicolumn{4}{|p{\dimexpr 16cm + 6\tabcolsep\relax}|}{\begin{enumerate}[start=11]\item Lighting and shading: interaction between light and surfaces, light sources, Phong reflection model, shading of polygonal surfaces, transparency, texture mapping\end{enumerate}}\\
\multicolumn{4}{|p{\dimexpr 16cm + 6\tabcolsep\relax}|}{\begin{enumerate}[start=12]\item Clipping: of line segments, of polygons, of other primitives\end{enumerate}}\\
\multicolumn{4}{|p{\dimexpr 16cm + 6\tabcolsep\relax}|}{\begin{enumerate}[start=13]\item Hidden surface removal: back-face culling\end{enumerate}}\\
\multicolumn{4}{|p{\dimexpr 16cm + 6\tabcolsep\relax}|}{\begin{enumerate}[start=14]\item Z-buffer sorting algorithm: depth sorting\end{enumerate}}\\
\multicolumn{4}{|p{\dimexpr 16cm + 6\tabcolsep\relax}|}{\begin{enumerate}[start=15]\item Rasterization: of line segments, of polygons\end{enumerate}}\\
\multicolumn{4}{|p{\dimexpr 16cm + 6\tabcolsep\relax}|}{\begin{enumerate}[start=16]\item Aliasing\end{enumerate}}\\
\multicolumn{4}{|p{\dimexpr 16cm + 6\tabcolsep\relax}|}{}\\
\hline
\multicolumn{4}{|p{16cm}|}{Basic Bibliography:}\\*
\multicolumn{4}{|p{\dimexpr 16cm + 6\tabcolsep\relax}|}{\begin{itemize}\item CONCI, Aura; AZEVEDO, Eduardo. Computação gráfica: geração de imagens. 8. ed. Rio de Janeiro: Campus, 2003.\end{itemize}}\\
\multicolumn{4}{|p{\dimexpr 16cm + 6\tabcolsep\relax}|}{\begin{itemize}\item ANGEL, Edward. Interactive computer graphics: a top-down approach with OpenGL. 6th ed. Boston: Addison Wesley, 2012.\end{itemize}}\\
\multicolumn{4}{|p{\dimexpr 16cm + 6\tabcolsep\relax}|}{\begin{itemize}\item HILL, F. S. Jr. Computer graphics: using OpenGL. 2nd ed. Upper Saddle River: Prentice-Hall, 2001.\end{itemize}}\\
\multicolumn{4}{|p{16cm}|}{}\\
\hline
\multicolumn{4}{|p{16cm}|}{Complementary Bibliography:}\\*
\multicolumn{4}{|p{\dimexpr 16cm + 6\tabcolsep\relax}|}{\begin{itemize}\item GOMES, Jonas de Miranda. Fundamentos da computação gráfica. Rio de Janeiro: Instituto Nacional de Matemática Pura e Aplicada, 2008.\end{itemize}}\\
\multicolumn{4}{|p{\dimexpr 16cm + 6\tabcolsep\relax}|}{\begin{itemize}\item HEARN, Donald. Computer graphics with OpenGL. 4th ed. Boston: Prentice-Hall, 2011.\end{itemize}}\\
\multicolumn{4}{|p{\dimexpr 16cm + 6\tabcolsep\relax}|}{\begin{itemize}\item FOLEY, James D. Fundamentals of interactive computer graphics. Reading, Mass.: Addison Wesley, 1984.\end{itemize}}\\
\multicolumn{4}{|p{\dimexpr 16cm + 6\tabcolsep\relax}|}{\begin{itemize}\item TORI, Romero et al. Fundamentos da computação gráfica: Compugrafia. São Paulo: Livros Técnicos e Científicos.\end{itemize}}\\
\multicolumn{4}{|p{\dimexpr 16cm + 6\tabcolsep\relax}|}{\begin{itemize}\item AZEVEDO, Eduardo; CONCI, Aura. Computação gráfica: teoria e prática. Rio de Janeiro: Elsevier.\end{itemize}}\\
\hline
\end{longtable}
\end{center}

\clearpage
\begin{center}
\begin{tabular}{lccr}
 \multirow{3}{*}{\includegraphics[height=2.7cm]{brasao.png}} &
 \multicolumn{2}{c}{\bfseries FEDERAL UNIVERSITY OF OURO PRETO} &
 \ \ \ \ \multirow{3}{*}{\includegraphics[height=2.7cm]{ufop.png}} \\
 & \multicolumn{2}{c}{\bfseries DEAN'S OFFICE FOR UNDERGRADUATE STUDIES} & \\
 & \multicolumn{2}{c}{} & \\
 & \multicolumn{2}{c}{\Large\bfseries COURSE PROGRAM} & \\
\end{tabular}
\end{center}

\begin{center}
\begin{longtable}{|p{4cm}|p{4cm}|p{4cm}|p{4cm}|}
\hline
\multicolumn{3}{|p{12cm}|}{Course Name in English:} &
\multicolumn{1}{p{4cm}|}{Code:} \\ 
\multicolumn{3}{|p{12cm}|}{\textbf{Compiler Construction I}} &
\textbf{BCC888}\\ 
\multicolumn{3}{|p{12cm}|}{Course Name in Portuguese:} & \\ 
\multicolumn{3}{|p{12cm}|}{\textbf{Construção de Compiladores I}} & \\ 
\hline
\multicolumn{3}{|p{12cm}|}{Department Name and Acronym:} & Academic Unit: \\ 
\multicolumn{3}{|p{12cm}|}{Department of Computer Science (DECOM)} & {ICEB} \\ 
\hline
\multicolumn{4}{|p{16cm}|}{Delivery Mode:
[X] in person \hspace{0.6cm}
[ ] blended \hspace{0.6cm}
[ ] distance learning}\\
\hline
\multicolumn{2}{|p{8cm}|}{Semester workload} &
\multicolumn{2}{p{8cm}|}{Weekly workload}\\
\hline
\multicolumn{1}{|p{4cm}|}{Total} &
\multicolumn{1}{p{4cm}|}{Extension} &
\multicolumn{1}{p{4cm}|}{Theoretical} &
\multicolumn{1}{p{4cm}|}{Practical} \\ 
\multicolumn{1}{|p{4cm}|}{60\,hours} &
\multicolumn{1}{p{4cm}|}{0\;hours} &
\multicolumn{1}{p{4cm}|}{4\;class hours} &
\multicolumn{1}{p{4cm}|}{0\;class hours} \\ 
\hline
\multicolumn{4}{|p{16cm}|}{Course Description:}\\
\multicolumn{4}{|p{16cm}|}{}\\
\multicolumn{4}{|p{\dimexpr 16cm + 6\tabcolsep\relax}|}{Implementation of programming languages: compilation and interpretation; lexical analysis; syntax analysis; semantic analysis.}\\
\multicolumn{4}{|p{16cm}|}{}\\
\hline
\multicolumn{4}{|p{16cm}|}{Course Content:}\\*
\multicolumn{4}{|p{\dimexpr 16cm + 6\tabcolsep\relax}|}{\begin{enumerate}[start=1]\item Implementation of programming languages: compilation and interpretation\end{enumerate}}\\
\multicolumn{4}{|p{\dimexpr 16cm + 6\tabcolsep\relax}|}{\begin{enumerate}[start=2]\item Lexical analysis: lexical analysis, ad hoc lexical analyzers, regular expressions, finite automata, and lexical analyzer generators\end{enumerate}}\\
\multicolumn{4}{|p{\dimexpr 16cm + 6\tabcolsep\relax}|}{\begin{enumerate}[start=3]\item Syntax analysis: syntax analysis, context-free grammars, top-down parsing, bottom-up parsing, parser generators, and semantic actions\end{enumerate}}\\
\multicolumn{4}{|p{\dimexpr 16cm + 6\tabcolsep\relax}|}{\begin{enumerate}[start=4]\item Semantic analysis: abstract syntax trees, symbol tables, type checking, and semantic rules\end{enumerate}}\\
\multicolumn{4}{|p{\dimexpr 16cm + 6\tabcolsep\relax}|}{}\\
\hline
\multicolumn{4}{|p{16cm}|}{Basic Bibliography:}\\*
\multicolumn{4}{|p{\dimexpr 16cm + 6\tabcolsep\relax}|}{\begin{itemize}\item APPEL, Andrew W. Modern compiler implementation in Java. 2nd ed. Cambridge: Cambridge University Press, c2002.\end{itemize}}\\
\multicolumn{4}{|p{\dimexpr 16cm + 6\tabcolsep\relax}|}{\begin{itemize}\item AHO, Alfred V. Compiladores, princípios, técnicas e ferramentas. São Paulo: Pearson Addison-Wesley, 2007.\end{itemize}}\\
\multicolumn{4}{|p{\dimexpr 16cm + 6\tabcolsep\relax}|}{\begin{itemize}\item APPEL, Andrew W. Modern compiler implementation in ML. New York: Cambridge University Press, c1998.\end{itemize}}\\
\multicolumn{4}{|p{16cm}|}{}\\
\hline
\multicolumn{4}{|p{16cm}|}{Complementary Bibliography:}\\*
\multicolumn{4}{|p{\dimexpr 16cm + 6\tabcolsep\relax}|}{\begin{itemize}\item MOGENSEN, Torben Ægidius. Introduction to Compiler Design. 2. ed. Springer, 2011. (item inexistente na UFOP)\end{itemize}}\\
\multicolumn{4}{|p{\dimexpr 16cm + 6\tabcolsep\relax}|}{\begin{itemize}\item GRÜNE, Dick; REEUWIJK, Kees van; BAL, Henri E.; JACOBS, Ceriel J. H.; LANGENDOEN, Koen. Modern Compiler Design. 2. ed. Springer, 2012. (item inexistente na UFOP)\end{itemize}}\\
\multicolumn{4}{|p{\dimexpr 16cm + 6\tabcolsep\relax}|}{\begin{itemize}\item WILHELM, Reinhard; SEIDL, Helmut; HACK, Sebastian. Compiler Design: Syntactic and Semantic Analysis. Springer, 2013. (item inexistente na UFOP)\end{itemize}}\\
\multicolumn{4}{|p{\dimexpr 16cm + 6\tabcolsep\relax}|}{\begin{itemize}\item FISCHER, Charles N.; CYTRON, Ron K.; LEBLANC, Richard J. Jr. Crafting a Compiler. Pearson, 2010. (item inexistente na UFOP)\end{itemize}}\\
\multicolumn{4}{|p{\dimexpr 16cm + 6\tabcolsep\relax}|}{\begin{itemize}\item COOPER, Keith D.; TORCZON, Linda. Engineering a Compiler. 2. ed. Elsevier, 2011. (item inexistente na UFOP)\end{itemize}}\\
\hline
\end{longtable}
\end{center}

\clearpage
\begin{center}
\begin{tabular}{lccr}
 \multirow{3}{*}{\includegraphics[height=2.7cm]{brasao.png}} &
 \multicolumn{2}{c}{\bfseries FEDERAL UNIVERSITY OF OURO PRETO} &
 \ \ \ \ \multirow{3}{*}{\includegraphics[height=2.7cm]{ufop.png}} \\
 & \multicolumn{2}{c}{\bfseries DEAN'S OFFICE FOR UNDERGRADUATE STUDIES} & \\
 & \multicolumn{2}{c}{} & \\
 & \multicolumn{2}{c}{\Large\bfseries COURSE PROGRAM} & \\
\end{tabular}
\end{center}

\begin{center}
\begin{longtable}{|p{4cm}|p{4cm}|p{4cm}|p{4cm}|}
\hline
\multicolumn{3}{|p{12cm}|}{Course Name in English:} &
\multicolumn{1}{p{4cm}|}{Code:} \\ 
\multicolumn{3}{|p{12cm}|}{\textbf{Design of Integrated Logic Circuits Using HDL (Hardware Description Language)}} &
\textbf{BCC568}\\ 
\multicolumn{3}{|p{12cm}|}{Course Name in Portuguese:} & \\ 
\multicolumn{3}{|p{12cm}|}{\textbf{Projeto de Circuitos Lógicos Integrados usando HDL (Hardware Description Language)}} & \\ 
\hline
\multicolumn{3}{|p{12cm}|}{Department Name and Acronym:} & Academic Unit: \\ 
\multicolumn{3}{|p{12cm}|}{Department of Computer Science (DECOM)} & {ICEB} \\ 
\hline
\multicolumn{4}{|p{16cm}|}{Delivery Mode:
[X] in person \hspace{0.6cm}
[ ] blended \hspace{0.6cm}
[ ] distance learning}\\
\hline
\multicolumn{2}{|p{8cm}|}{Semester workload} &
\multicolumn{2}{p{8cm}|}{Weekly workload}\\
\hline
\multicolumn{1}{|p{4cm}|}{Total} &
\multicolumn{1}{p{4cm}|}{Extension} &
\multicolumn{1}{p{4cm}|}{Theoretical} &
\multicolumn{1}{p{4cm}|}{Practical} \\ 
\multicolumn{1}{|p{4cm}|}{60\,hours} &
\multicolumn{1}{p{4cm}|}{0\;hours} &
\multicolumn{1}{p{4cm}|}{1\;class hours} &
\multicolumn{1}{p{4cm}|}{3\;class hours} \\ 
\hline
\multicolumn{4}{|p{16cm}|}{Course Description:}\\
\multicolumn{4}{|p{16cm}|}{}\\
\multicolumn{4}{|p{\dimexpr 16cm + 6\tabcolsep\relax}|}{Overview of the Digital Integrated Circuit Design flow and its importance for the country's economic development. Concepts and applicability of application-specific integrated circuits (ASICs). Main hardware description languages (HDLs). Integrated Circuit Design techniques. Use of HDL for integrated circuit design.}\\
\multicolumn{4}{|p{16cm}|}{}\\
\hline
\multicolumn{4}{|p{16cm}|}{Course Content:}\\*
\multicolumn{4}{|p{\dimexpr 16cm + 6\tabcolsep\relax}|}{\begin{enumerate}[start=1]\item Introduction: digital integrated circuit design flow (architecture, design, verification, logic synthesis, and physical synthesis)\end{enumerate}}\\
\multicolumn{4}{|p{\dimexpr 16cm + 6\tabcolsep\relax}|}{\begin{enumerate}[start=2]\item Introduction to FPGA (gate arrays, capacity, Look-Up Tables - LUTs)\end{enumerate}}\\
\multicolumn{4}{|p{\dimexpr 16cm + 6\tabcolsep\relax}|}{\begin{enumerate}[start=3]\item Digital Integrated Circuit Architecture (processors, communication buses for System-on-Chip - SoC, peripherals, and memory access)\end{enumerate}}\\
\multicolumn{4}{|p{\dimexpr 16cm + 6\tabcolsep\relax}|}{\begin{enumerate}[start=4]\item Hardware Description Language: Verilog 2001\end{enumerate}}\\
\multicolumn{4}{|p{\dimexpr 16cm + 6\tabcolsep\relax}|}{\begin{enumerate}[start=5]\item Description of basic combinational and sequential digital integrated circuits (logic gates, half adder, full adder, multiplexers and demultiplexers, encoders and decoders, counters, registers and flip-flops) using Verilog 2001\end{enumerate}}\\
\multicolumn{4}{|p{\dimexpr 16cm + 6\tabcolsep\relax}|}{\begin{enumerate}[start=6]\item CISC Architecture and RISC Architecture\end{enumerate}}\\
\multicolumn{4}{|p{\dimexpr 16cm + 6\tabcolsep\relax}|}{\begin{enumerate}[start=7]\item Integrated circuit design projects using HDL\end{enumerate}}\\
\multicolumn{4}{|p{\dimexpr 16cm + 6\tabcolsep\relax}|}{}\\
\hline
\multicolumn{4}{|p{16cm}|}{Basic Bibliography:}\\*
\multicolumn{4}{|p{\dimexpr 16cm + 6\tabcolsep\relax}|}{\begin{itemize}\item PEDRONI, Volnei A. Eletrônica digital moderna e VHDL. Rio de Janeiro: Elsevier, 2010. ISBN 9788535234657.\end{itemize}}\\
\multicolumn{4}{|p{\dimexpr 16cm + 6\tabcolsep\relax}|}{\begin{itemize}\item COFFMAN, Ken. Real world FPGA design with Verilog. Pearson Education, 1999. ISBN 0130998516.\end{itemize}}\\
\multicolumn{4}{|p{\dimexpr 16cm + 6\tabcolsep\relax}|}{\begin{itemize}\item CHU, Pong P. FPGA prototyping by VHDL examples: Xilinx Spartan-3 version. John Wiley \& Sons, 2011. ISBN 9780470185315.\end{itemize}}\\
\multicolumn{4}{|p{16cm}|}{}\\
\hline
\multicolumn{4}{|p{16cm}|}{Complementary Bibliography:}\\*
\multicolumn{4}{|p{\dimexpr 16cm + 6\tabcolsep\relax}|}{\begin{itemize}\item ERCEGOVAC, Miloš Dragutin; LANG, Tomás; MORENO, Jaime H. Introdução aos sistemas digitais. Bookman, 2000. ISBN 0471527998.\end{itemize}}\\
\multicolumn{4}{|p{\dimexpr 16cm + 6\tabcolsep\relax}|}{\begin{itemize}\item TOCCI, Ronald J.; WIDMER, Neal S.; MOSS, Gregory L. Sistemas Digitais: Princípios e Aplicações. 11. ed. Rio de Janeiro: MZ Editora, 2007. ISBN 9788576059226.\end{itemize}}\\
\multicolumn{4}{|p{\dimexpr 16cm + 6\tabcolsep\relax}|}{\begin{itemize}\item VAHID, Frank. Sistemas Digitais. Bookman Editora, 2009. ISBN 9788577801909.\end{itemize}}\\
\multicolumn{4}{|p{\dimexpr 16cm + 6\tabcolsep\relax}|}{\begin{itemize}\item TANENBAUM, Andrew S. Organização estruturada de computadores. 3. ed. Rio de Janeiro: Prentice-Hall do Brasil, 1992. 460 p. ISBN 857054040X.\end{itemize}}\\
\multicolumn{4}{|p{\dimexpr 16cm + 6\tabcolsep\relax}|}{\begin{itemize}\item STALLINGS, William. Arquitetura e organização de computadores. 8. ed. São Paulo: Pearson, c2010. 624 p. ISBN 9788576055648.\end{itemize}}\\
\hline
\end{longtable}
\end{center}

\clearpage
\begin{center}
\begin{tabular}{lccr}
 \multirow{3}{*}{\includegraphics[height=2.7cm]{brasao.png}} &
 \multicolumn{2}{c}{\bfseries FEDERAL UNIVERSITY OF OURO PRETO} &
 \ \ \ \ \multirow{3}{*}{\includegraphics[height=2.7cm]{ufop.png}} \\
 & \multicolumn{2}{c}{\bfseries DEAN'S OFFICE FOR UNDERGRADUATE STUDIES} & \\
 & \multicolumn{2}{c}{} & \\
 & \multicolumn{2}{c}{\Large\bfseries COURSE PROGRAM} & \\
\end{tabular}
\end{center}

\begin{center}
\begin{longtable}{|p{4cm}|p{4cm}|p{4cm}|p{4cm}|}
\hline
\multicolumn{3}{|p{12cm}|}{Course Name in English:} &
\multicolumn{1}{p{4cm}|}{Code:} \\ 
\multicolumn{3}{|p{12cm}|}{\textbf{Mobile Computing}} &
\textbf{BCC581}\\ 
\multicolumn{3}{|p{12cm}|}{Course Name in Portuguese:} & \\ 
\multicolumn{3}{|p{12cm}|}{\textbf{Computação Móvel}} & \\ 
\hline
\multicolumn{3}{|p{12cm}|}{Department Name and Acronym:} & Academic Unit: \\ 
\multicolumn{3}{|p{12cm}|}{Department of Computer Science (DECOM)} & {ICEB} \\ 
\hline
\multicolumn{4}{|p{16cm}|}{Delivery Mode:
[X] in person \hspace{0.6cm}
[ ] blended \hspace{0.6cm}
[ ] distance learning}\\
\hline
\multicolumn{2}{|p{8cm}|}{Semester workload} &
\multicolumn{2}{p{8cm}|}{Weekly workload}\\
\hline
\multicolumn{1}{|p{4cm}|}{Total} &
\multicolumn{1}{p{4cm}|}{Extension} &
\multicolumn{1}{p{4cm}|}{Theoretical} &
\multicolumn{1}{p{4cm}|}{Practical} \\ 
\multicolumn{1}{|p{4cm}|}{60\,hours} &
\multicolumn{1}{p{4cm}|}{0\;hours} &
\multicolumn{1}{p{4cm}|}{4\;class hours} &
\multicolumn{1}{p{4cm}|}{0\;class hours} \\ 
\hline
\multicolumn{4}{|p{16cm}|}{Course Description:}\\
\multicolumn{4}{|p{16cm}|}{}\\
\multicolumn{4}{|p{\dimexpr 16cm + 6\tabcolsep\relax}|}{Signal propagation, wireless communication networks, location-based services, communication protocols, information management, structured and unstructured networks. Development of mobile applications on current platforms.}\\
\multicolumn{4}{|p{16cm}|}{}\\
\hline
\multicolumn{4}{|p{16cm}|}{Course Content:}\\*
\multicolumn{4}{|p{\dimexpr 16cm + 6\tabcolsep\relax}|}{\begin{enumerate}[start=1]\item Signal propagation\end{enumerate}}\\
\multicolumn{4}{|p{\dimexpr 16cm + 6\tabcolsep\relax}|}{\begin{enumerate}[start=2]\item Wireless communication networks\end{enumerate}}\\
\multicolumn{4}{|p{\dimexpr 16cm + 6\tabcolsep\relax}|}{\begin{enumerate}[start=3]\item Location-based services\end{enumerate}}\\
\multicolumn{4}{|p{\dimexpr 16cm + 6\tabcolsep\relax}|}{\begin{enumerate}[start=4]\item Communication protocols\end{enumerate}}\\
\multicolumn{4}{|p{\dimexpr 16cm + 6\tabcolsep\relax}|}{\begin{enumerate}[start=5]\item Information management\end{enumerate}}\\
\multicolumn{4}{|p{\dimexpr 16cm + 6\tabcolsep\relax}|}{\begin{enumerate}[start=6]\item Structured and unstructured networks\end{enumerate}}\\
\multicolumn{4}{|p{\dimexpr 16cm + 6\tabcolsep\relax}|}{\begin{enumerate}[start=7]\item Mobile application development: Android platform, Windows platform, iOS platform\end{enumerate}}\\
\multicolumn{4}{|p{\dimexpr 16cm + 6\tabcolsep\relax}|}{}\\
\hline
\multicolumn{4}{|p{16cm}|}{Basic Bibliography:}\\*
\multicolumn{4}{|p{\dimexpr 16cm + 6\tabcolsep\relax}|}{\begin{itemize}\item RAPPAPORT, Theodore. Wireless Personal Communications. 2. ed. Kluwer Academic, 2007. (item inexistente na UFOP)\end{itemize}}\\
\multicolumn{4}{|p{\dimexpr 16cm + 6\tabcolsep\relax}|}{\begin{itemize}\item MATEUS, G. R. Introdução à Computação Móvel. 11. ed. Rio de Janeiro: Escola de Computação, 1998. (item inexistente na UFOP)\end{itemize}}\\
\multicolumn{4}{|p{\dimexpr 16cm + 6\tabcolsep\relax}|}{\begin{itemize}\item TERUEL, E. Web Mobile: Desenvolva Sites para Dispositivos Móveis com Tecnologias de Uso Livre. Ciência Moderna, 2010. (item inexistente na UFOP)\end{itemize}}\\
\multicolumn{4}{|p{\dimexpr 16cm + 6\tabcolsep\relax}|}{\begin{itemize}\item SCHILLER, Jochen H. Mobile communications. 2nd ed. London: Addison Wesley, 2003.\end{itemize}}\\
\multicolumn{4}{|p{16cm}|}{}\\
\hline
\multicolumn{4}{|p{16cm}|}{Complementary Bibliography:}\\*
\multicolumn{4}{|p{\dimexpr 16cm + 6\tabcolsep\relax}|}{\begin{itemize}\item LECHETA, Ricardo R. Google Android: aprenda a criar aplicações para dispositivos móveis com o Android SDK. 2. ed. São Paulo: Novatec Editora, 2010. (item inexistente na UFOP)\end{itemize}}\\
\multicolumn{4}{|p{\dimexpr 16cm + 6\tabcolsep\relax}|}{\begin{itemize}\item BORGES JÚNIOR, M. Aplicativos Móveis: Aplicativos para Dispositivos Móveis usando C\#.Net. Rio de Janeiro: Ciência Moderna, 2005. (item inexistente na UFOP)\end{itemize}}\\
\multicolumn{4}{|p{\dimexpr 16cm + 6\tabcolsep\relax}|}{\begin{itemize}\item PETERSON, Larry L.; DAVIE, Bruce S.; tradução Daniel Vieira; revisão técnica Marcos Simplício. Redes de computadores: uma abordagem de sistemas. Rio de Janeiro: Campus.\end{itemize}}\\
\multicolumn{4}{|p{\dimexpr 16cm + 6\tabcolsep\relax}|}{\begin{itemize}\item OPPENHEIM, Alan V. Discrete-time signal processing. 3rd ed. New Delhi: Pearson, c2010.\end{itemize}}\\
\hline
\end{longtable}
\end{center}

\clearpage
\begin{center}
\begin{tabular}{lccr}
 \multirow{3}{*}{\includegraphics[height=2.7cm]{brasao.png}} &
 \multicolumn{2}{c}{\bfseries FEDERAL UNIVERSITY OF OURO PRETO} &
 \ \ \ \ \multirow{3}{*}{\includegraphics[height=2.7cm]{ufop.png}} \\
 & \multicolumn{2}{c}{\bfseries DEAN'S OFFICE FOR UNDERGRADUATE STUDIES} & \\
 & \multicolumn{2}{c}{} & \\
 & \multicolumn{2}{c}{\Large\bfseries COURSE PROGRAM} & \\
\end{tabular}
\end{center}

\begin{center}
\begin{longtable}{|p{4cm}|p{4cm}|p{4cm}|p{4cm}|}
\hline
\multicolumn{3}{|p{12cm}|}{Course Name in English:} &
\multicolumn{1}{p{4cm}|}{Code:} \\ 
\multicolumn{3}{|p{12cm}|}{\textbf{Cloud Computing}} &
\textbf{BCC582}\\ 
\multicolumn{3}{|p{12cm}|}{Course Name in Portuguese:} & \\ 
\multicolumn{3}{|p{12cm}|}{\textbf{Computação nas Nuvens}} & \\ 
\hline
\multicolumn{3}{|p{12cm}|}{Department Name and Acronym:} & Academic Unit: \\ 
\multicolumn{3}{|p{12cm}|}{Department of Computer Science (DECOM)} & {ICEB} \\ 
\hline
\multicolumn{4}{|p{16cm}|}{Delivery Mode:
[X] in person \hspace{0.6cm}
[ ] blended \hspace{0.6cm}
[ ] distance learning}\\
\hline
\multicolumn{2}{|p{8cm}|}{Semester workload} &
\multicolumn{2}{p{8cm}|}{Weekly workload}\\
\hline
\multicolumn{1}{|p{4cm}|}{Total} &
\multicolumn{1}{p{4cm}|}{Extension} &
\multicolumn{1}{p{4cm}|}{Theoretical} &
\multicolumn{1}{p{4cm}|}{Practical} \\ 
\multicolumn{1}{|p{4cm}|}{60\,hours} &
\multicolumn{1}{p{4cm}|}{0\;hours} &
\multicolumn{1}{p{4cm}|}{4\;class hours} &
\multicolumn{1}{p{4cm}|}{0\;class hours} \\ 
\hline
\multicolumn{4}{|p{16cm}|}{Course Description:}\\
\multicolumn{4}{|p{16cm}|}{}\\
\multicolumn{4}{|p{\dimexpr 16cm + 6\tabcolsep\relax}|}{Fundamentals of Cloud Computing: terminology and concepts; virtualization; elasticity, resilience, on-demand and metered use; benefits, challenges, and risks of platforms and services; administration and governance; Software as a Service (SaaS), Platform as a Service (PaaS), and Infrastructure as a Service (IaaS) Cloud Delivery Models; Service Level Agreements (SLAs) for Cloud-based IT Resources.}\\
\multicolumn{4}{|p{16cm}|}{}\\
\hline
\multicolumn{4}{|p{16cm}|}{Course Content:}\\*
\multicolumn{4}{|p{\dimexpr 16cm + 6\tabcolsep\relax}|}{\begin{enumerate}[start=1]\item Fundamentals of Cloud Computing: terminology and concepts\end{enumerate}}\\
\multicolumn{4}{|p{\dimexpr 16cm + 6\tabcolsep\relax}|}{\begin{enumerate}[start=2]\item Virtualization: XEN, QEMU, Eucalyptus, Private Cloud\end{enumerate}}\\
\multicolumn{4}{|p{\dimexpr 16cm + 6\tabcolsep\relax}|}{\begin{enumerate}[start=3]\item Elasticity, Resilience, On-Demand, and Metered Use\end{enumerate}}\\
\multicolumn{4}{|p{\dimexpr 16cm + 6\tabcolsep\relax}|}{\begin{enumerate}[start=4]\item Software as a Service (SaaS), Platform as a Service (PaaS), and Infrastructure as a Service (IaaS) Cloud Delivery\end{enumerate}}\\
\multicolumn{4}{|p{\dimexpr 16cm + 6\tabcolsep\relax}|}{\begin{enumerate}[start=5]\item Benefits, challenges, and risks of platforms and services\end{enumerate}}\\
\multicolumn{4}{|p{\dimexpr 16cm + 6\tabcolsep\relax}|}{\begin{enumerate}[start=6]\item Creating a Cloud solution: Ruby on Rails as PaaS; Software-Defined Routing as IaaS; SOA as SaaS\end{enumerate}}\\
\multicolumn{4}{|p{\dimexpr 16cm + 6\tabcolsep\relax}|}{\begin{enumerate}[start=7]\item Administration and Governance\end{enumerate}}\\
\multicolumn{4}{|p{\dimexpr 16cm + 6\tabcolsep\relax}|}{\begin{enumerate}[start=8]\item Service Level Agreements (SLAs) for Cloud-based IT Resources\end{enumerate}}\\
\multicolumn{4}{|p{\dimexpr 16cm + 6\tabcolsep\relax}|}{\begin{enumerate}[start=9]\item Performance analysis\end{enumerate}}\\
\multicolumn{4}{|p{\dimexpr 16cm + 6\tabcolsep\relax}|}{}\\
\hline
\multicolumn{4}{|p{16cm}|}{Basic Bibliography:}\\*
\multicolumn{4}{|p{\dimexpr 16cm + 6\tabcolsep\relax}|}{\begin{itemize}\item BIRMAN, Kenneth. Guide to Reliable Distributed Systems: Building High-Assurance Applications and Cloud-Hosted Services. Springer, 2012. (item inexistente na UFOP)\end{itemize}}\\
\multicolumn{4}{|p{\dimexpr 16cm + 6\tabcolsep\relax}|}{\begin{itemize}\item DEITEL, Paul J. Internet Applications e desenvolvimento Web para programadores. São Paulo: Prentice-Hall, 2009. (item inexistente na UFOP)\end{itemize}}\\
\multicolumn{4}{|p{\dimexpr 16cm + 6\tabcolsep\relax}|}{\begin{itemize}\item KUROSE, James F. Redes de computadores e a Internet: uma abordagem top-down. 5. ed. São Paulo: Pearson Addison Wesley, 2010.\end{itemize}}\\
\multicolumn{4}{|p{16cm}|}{}\\
\hline
\multicolumn{4}{|p{16cm}|}{Complementary Bibliography:}\\*
\multicolumn{4}{|p{\dimexpr 16cm + 6\tabcolsep\relax}|}{\begin{itemize}\item COMER, Douglas E.; tradução ARX Publicações. Interligação em redes com TCP/IP, volume 1: princípios, protocolos e arquitetura. 6. ed. Rio de Janeiro: Campus, c2015.\end{itemize}}\\
\multicolumn{4}{|p{\dimexpr 16cm + 6\tabcolsep\relax}|}{\begin{itemize}\item STALLINGS, William. Arquitetura e organização de computadores. 8. ed. São Paulo: Pearson, c2010.\end{itemize}}\\
\multicolumn{4}{|p{\dimexpr 16cm + 6\tabcolsep\relax}|}{\begin{itemize}\item ERL, Thomas; tradução Carlos Schafranski e Edson Furmankiewicz; revisão técnica Francisco Calmon N. da Gama, Ricardo A. de Castro Barbosa. SOA: princípios de design de serviços. São Paulo: Prentice Hall, c2009.\end{itemize}}\\
\multicolumn{4}{|p{\dimexpr 16cm + 6\tabcolsep\relax}|}{\begin{itemize}\item PETERSON, Larry L.; DAVIE, Bruce S.; tradução Daniel Vieira; revisão técnica Marcos Simplício. Redes de computadores: uma abordagem de sistemas. Rio de Janeiro: Campus.\end{itemize}}\\
\hline
\end{longtable}
\end{center}

\clearpage
\begin{center}
\begin{tabular}{lccr}
 \multirow{3}{*}{\includegraphics[height=2.7cm]{brasao.png}} &
 \multicolumn{2}{c}{\bfseries FEDERAL UNIVERSITY OF OURO PRETO} &
 \ \ \ \ \multirow{3}{*}{\includegraphics[height=2.7cm]{ufop.png}} \\
 & \multicolumn{2}{c}{\bfseries DEAN'S OFFICE FOR UNDERGRADUATE STUDIES} & \\
 & \multicolumn{2}{c}{} & \\
 & \multicolumn{2}{c}{\Large\bfseries COURSE PROGRAM} & \\
\end{tabular}
\end{center}

\begin{center}
\begin{longtable}{|p{4cm}|p{4cm}|p{4cm}|p{4cm}|}
\hline
\multicolumn{3}{|p{12cm}|}{Course Name in English:} &
\multicolumn{1}{p{4cm}|}{Code:} \\ 
\multicolumn{3}{|p{12cm}|}{\textbf{Cryptography and Security Systems}} &
\textbf{BCC583}\\ 
\multicolumn{3}{|p{12cm}|}{Course Name in Portuguese:} & \\ 
\multicolumn{3}{|p{12cm}|}{\textbf{Criptografia e Segurança de Sistemas}} & \\ 
\hline
\multicolumn{3}{|p{12cm}|}{Department Name and Acronym:} & Academic Unit: \\ 
\multicolumn{3}{|p{12cm}|}{Department of Computer Science (DECOM)} & {ICEB} \\ 
\hline
\multicolumn{4}{|p{16cm}|}{Delivery Mode:
[X] in person \hspace{0.6cm}
[ ] blended \hspace{0.6cm}
[ ] distance learning}\\
\hline
\multicolumn{2}{|p{8cm}|}{Semester workload} &
\multicolumn{2}{p{8cm}|}{Weekly workload}\\
\hline
\multicolumn{1}{|p{4cm}|}{Total} &
\multicolumn{1}{p{4cm}|}{Extension} &
\multicolumn{1}{p{4cm}|}{Theoretical} &
\multicolumn{1}{p{4cm}|}{Practical} \\ 
\multicolumn{1}{|p{4cm}|}{60\,hours} &
\multicolumn{1}{p{4cm}|}{0\;hours} &
\multicolumn{1}{p{4cm}|}{4\;class hours} &
\multicolumn{1}{p{4cm}|}{0\;class hours} \\ 
\hline
\multicolumn{4}{|p{16cm}|}{Course Description:}\\
\multicolumn{4}{|p{16cm}|}{}\\
\multicolumn{4}{|p{\dimexpr 16cm + 6\tabcolsep\relax}|}{Network, system, and information security, cybersecurity. Cryptography. Digital signatures, digital certificates, attribute certificates. Network security. Cryptographic media: biometric identifiers. Cybersecurity and impacts on contemporary society.}\\
\multicolumn{4}{|p{16cm}|}{}\\
\hline
\multicolumn{4}{|p{16cm}|}{Course Content:}\\*
\multicolumn{4}{|p{\dimexpr 16cm + 6\tabcolsep\relax}|}{\begin{enumerate}[start=1]\item Introduction to security: information security, network security, system security, and cybersecurity\end{enumerate}}\\
\multicolumn{4}{|p{\dimexpr 16cm + 6\tabcolsep\relax}|}{\begin{enumerate}[start=2]\item Symmetric and asymmetric cryptography\end{enumerate}}\\
\multicolumn{4}{|p{\dimexpr 16cm + 6\tabcolsep\relax}|}{\begin{enumerate}[start=3]\item Cryptographic algorithms\end{enumerate}}\\
\multicolumn{4}{|p{\dimexpr 16cm + 6\tabcolsep\relax}|}{\begin{enumerate}[start=4]\item Digital signatures: digital certificates and the PKI (ICP) standard\end{enumerate}}\\
\multicolumn{4}{|p{\dimexpr 16cm + 6\tabcolsep\relax}|}{\begin{enumerate}[start=5]\item ICP-Brasil and ICP-Edu\end{enumerate}}\\
\multicolumn{4}{|p{\dimexpr 16cm + 6\tabcolsep\relax}|}{\begin{enumerate}[start=6]\item Network security: monitoring and network intrusion\end{enumerate}}\\
\multicolumn{4}{|p{\dimexpr 16cm + 6\tabcolsep\relax}|}{\begin{enumerate}[start=7]\item Principles of hacking\end{enumerate}}\\
\multicolumn{4}{|p{\dimexpr 16cm + 6\tabcolsep\relax}|}{\begin{enumerate}[start=8]\item Cryptographic media and biometric identifiers\end{enumerate}}\\
\multicolumn{4}{|p{\dimexpr 16cm + 6\tabcolsep\relax}|}{\begin{enumerate}[start=9]\item ISO27000\end{enumerate}}\\
\multicolumn{4}{|p{\dimexpr 16cm + 6\tabcolsep\relax}|}{\begin{enumerate}[start=10]\item ITIL and COBIT within the information security context\end{enumerate}}\\
\multicolumn{4}{|p{\dimexpr 16cm + 6\tabcolsep\relax}|}{\begin{enumerate}[start=11]\item IISP certifications, among others\end{enumerate}}\\
\multicolumn{4}{|p{\dimexpr 16cm + 6\tabcolsep\relax}|}{\begin{enumerate}[start=12]\item Cybersecurity and impacts on contemporary society\end{enumerate}}\\
\multicolumn{4}{|p{\dimexpr 16cm + 6\tabcolsep\relax}|}{}\\
\hline
\multicolumn{4}{|p{16cm}|}{Basic Bibliography:}\\*
\multicolumn{4}{|p{\dimexpr 16cm + 6\tabcolsep\relax}|}{\begin{itemize}\item STALLINGS, William. Cryptography and Network Security: Principles and Practice. 7. ed. N.Y.: Pearson, 2016. (item inexistente na UFOP)\end{itemize}}\\
\multicolumn{4}{|p{\dimexpr 16cm + 6\tabcolsep\relax}|}{\begin{itemize}\item GOLDREICH, Oded. Foundations of cryptography. Cambridge: Cambridge University Press, 2009. (item inexistente na UFOP)\end{itemize}}\\
\multicolumn{4}{|p{\dimexpr 16cm + 6\tabcolsep\relax}|}{\begin{itemize}\item MENEZES, A. J. Handbook of applied cryptography. Boca Raton: CRC Press, c1997.\end{itemize}}\\
\multicolumn{4}{|p{16cm}|}{}\\
\hline
\multicolumn{4}{|p{16cm}|}{Complementary Bibliography:}\\*
\multicolumn{4}{|p{\dimexpr 16cm + 6\tabcolsep\relax}|}{\begin{itemize}\item TRAPPE, Wade; WASHINGTON, Lawrence C. Introduction to cryptography: with coding theory. 2. ed. Upper Saddle River, N.J.: Pearson Prentice Hall, c2006.\end{itemize}}\\
\multicolumn{4}{|p{\dimexpr 16cm + 6\tabcolsep\relax}|}{\begin{itemize}\item COUTINHO, S. C. Números inteiros e criptografia RSA. 2. ed. Rio de Janeiro: IMPA, 2009.\end{itemize}}\\
\multicolumn{4}{|p{\dimexpr 16cm + 6\tabcolsep\relax}|}{\begin{itemize}\item KATZ, Jonathan. Introduction to modern cryptography. New York: Chapman \& Hall/CRC, 2008.\end{itemize}}\\
\multicolumn{4}{|p{\dimexpr 16cm + 6\tabcolsep\relax}|}{\begin{itemize}\item GOLDREICH, Oded. Foundations of cryptography: basic tools. New York: Cambridge University Press, 2001.\end{itemize}}\\
\hline
\end{longtable}
\end{center}

\clearpage
\begin{center}
\begin{tabular}{lccr}
 \multirow{3}{*}{\includegraphics[height=2.7cm]{brasao.png}} &
 \multicolumn{2}{c}{\bfseries FEDERAL UNIVERSITY OF OURO PRETO} &
 \ \ \ \ \multirow{3}{*}{\includegraphics[height=2.7cm]{ufop.png}} \\
 & \multicolumn{2}{c}{\bfseries DEAN'S OFFICE FOR UNDERGRADUATE STUDIES} & \\
 & \multicolumn{2}{c}{} & \\
 & \multicolumn{2}{c}{\Large\bfseries COURSE PROGRAM} & \\
\end{tabular}
\end{center}

\begin{center}
\begin{longtable}{|p{4cm}|p{4cm}|p{4cm}|p{4cm}|}
\hline
\multicolumn{3}{|p{12cm}|}{Course Name in English:} &
\multicolumn{1}{p{4cm}|}{Code:} \\ 
\multicolumn{3}{|p{12cm}|}{\textbf{Embedded Systems}} &
\textbf{BCC585}\\ 
\multicolumn{3}{|p{12cm}|}{Course Name in Portuguese:} & \\ 
\multicolumn{3}{|p{12cm}|}{\textbf{Sistemas Embutidos}} & \\ 
\hline
\multicolumn{3}{|p{12cm}|}{Department Name and Acronym:} & Academic Unit: \\ 
\multicolumn{3}{|p{12cm}|}{Department of Computer Science (DECOM)} & {ICEB} \\ 
\hline
\multicolumn{4}{|p{16cm}|}{Delivery Mode:
[X] in person \hspace{0.6cm}
[ ] blended \hspace{0.6cm}
[ ] distance learning}\\
\hline
\multicolumn{2}{|p{8cm}|}{Semester workload} &
\multicolumn{2}{p{8cm}|}{Weekly workload}\\
\hline
\multicolumn{1}{|p{4cm}|}{Total} &
\multicolumn{1}{p{4cm}|}{Extension} &
\multicolumn{1}{p{4cm}|}{Theoretical} &
\multicolumn{1}{p{4cm}|}{Practical} \\ 
\multicolumn{1}{|p{4cm}|}{60\,hours} &
\multicolumn{1}{p{4cm}|}{0\;hours} &
\multicolumn{1}{p{4cm}|}{2\;class hours} &
\multicolumn{1}{p{4cm}|}{2\;class hours} \\ 
\hline
\multicolumn{4}{|p{16cm}|}{Course Description:}\\
\multicolumn{4}{|p{16cm}|}{}\\
\multicolumn{4}{|p{\dimexpr 16cm + 6\tabcolsep\relax}|}{Fixed-point and floating-point integer arithmetic. Computer and microcontroller architecture. Microcontroller programming in C and Assembly. Input and output concepts for microcontrollers. Use of finite-state machines to solve programming problems. Programming applications with keypads, liquid crystal and seven-segment displays, interfacing with actuators (motors). Analog-to-digital (ADC) and digital-to-analog (DAC) converters. Interrupt and Timer concepts. Data transmission. Embedded Systems Design.}\\
\multicolumn{4}{|p{16cm}|}{}\\
\hline
\multicolumn{4}{|p{16cm}|}{Course Content:}\\*
\multicolumn{4}{|p{\dimexpr 16cm + 6\tabcolsep\relax}|}{\begin{enumerate}[start=1]\item Course plan presentation and Introduction: what is an embedded system, types of embedded systems\end{enumerate}}\\
\multicolumn{4}{|p{\dimexpr 16cm + 6\tabcolsep\relax}|}{\begin{enumerate}[start=2]\item Numerical representation: integer, fixed-point, floating-point\end{enumerate}}\\
\multicolumn{4}{|p{\dimexpr 16cm + 6\tabcolsep\relax}|}{\begin{enumerate}[start=3]\item C language and tools: compiler, linker, loader\end{enumerate}}\\
\multicolumn{4}{|p{\dimexpr 16cm + 6\tabcolsep\relax}|}{\begin{enumerate}[start=4]\item Review of computer architecture and introduction to microcontroller architecture: architecture, instruction set, programming, Assembly for microcontrollers\end{enumerate}}\\
\multicolumn{4}{|p{\dimexpr 16cm + 6\tabcolsep\relax}|}{\begin{enumerate}[start=5]\item Basic Input/Output interfacing: buttons, LEDs, displays, segments, LCD; basic input/output interfacing with actuators (DC motor, stepper motor, servo motor); analog interfacing (ADC, DAC, PWM)\end{enumerate}}\\
\multicolumn{4}{|p{\dimexpr 16cm + 6\tabcolsep\relax}|}{\begin{enumerate}[start=6]\item Interrupts and Timers: concept, programming, resource-sharing issues\end{enumerate}}\\
\multicolumn{4}{|p{\dimexpr 16cm + 6\tabcolsep\relax}|}{\begin{enumerate}[start=7]\item Serial Communication: RS232, I2C, SPI, CAN\end{enumerate}}\\
\multicolumn{4}{|p{\dimexpr 16cm + 6\tabcolsep\relax}|}{\begin{enumerate}[start=8]\item Embedded systems design: modeling methodology, finite-state machines, software development cycle, data flow diagrams, statecharts\end{enumerate}}\\
\multicolumn{4}{|p{\dimexpr 16cm + 6\tabcolsep\relax}|}{}\\
\hline
\multicolumn{4}{|p{16cm}|}{Basic Bibliography:}\\*
\multicolumn{4}{|p{\dimexpr 16cm + 6\tabcolsep\relax}|}{\begin{itemize}\item LEE, Edward Ashford. Introduction to embedded systems: A cyber-physical systems approach. MIT Press, 2016. (item inexistente na UFOP)\end{itemize}}\\
\multicolumn{4}{|p{\dimexpr 16cm + 6\tabcolsep\relax}|}{\begin{itemize}\item NOERGAARD, Tammy. Embedded systems architecture: a comprehensive guide for engineers and programmers. Amsterdam.\end{itemize}}\\
\multicolumn{4}{|p{\dimexpr 16cm + 6\tabcolsep\relax}|}{\begin{itemize}\item BALL, Stuart R. Embedded microprocessor systems: real world design. 2. ed. Boston: Newnes, 2002.\end{itemize}}\\
\multicolumn{4}{|p{16cm}|}{}\\
\hline
\multicolumn{4}{|p{16cm}|}{Complementary Bibliography:}\\*
\multicolumn{4}{|p{\dimexpr 16cm + 6\tabcolsep\relax}|}{\begin{itemize}\item WILMSHURST, Tim. Designing embedded systems with PIC microcontrollers: principles and applications. Amsterdã.\end{itemize}}\\
\multicolumn{4}{|p{\dimexpr 16cm + 6\tabcolsep\relax}|}{\begin{itemize}\item PONT, Michael J. Embedded C. Boston, MA: Addison Wesley, 2002.\end{itemize}}\\
\multicolumn{4}{|p{\dimexpr 16cm + 6\tabcolsep\relax}|}{\begin{itemize}\item PEREIRA, Fábio. Microcontroladores PIC: programação em C. 7. ed. São Paulo: Érica, 2007.\end{itemize}}\\
\multicolumn{4}{|p{\dimexpr 16cm + 6\tabcolsep\relax}|}{\begin{itemize}\item HALLINAN, Christopher. Embedded Linux Primer: a practical, real-world approach. 2. ed. New York: Prentice Hall, 2011.\end{itemize}}\\
\multicolumn{4}{|p{\dimexpr 16cm + 6\tabcolsep\relax}|}{\begin{itemize}\item SOUZA, David José de. Desbravando o PIC: ampliado e atualizado para PIC16F628A. 9. ed. São Paulo: Érica, 2005.\end{itemize}}\\
\hline
\end{longtable}
\end{center}

\clearpage
\begin{center}
\begin{tabular}{lccr}
 \multirow{3}{*}{\includegraphics[height=2.7cm]{brasao.png}} &
 \multicolumn{2}{c}{\bfseries FEDERAL UNIVERSITY OF OURO PRETO} &
 \ \ \ \ \multirow{3}{*}{\includegraphics[height=2.7cm]{ufop.png}} \\
 & \multicolumn{2}{c}{\bfseries DEAN'S OFFICE FOR UNDERGRADUATE STUDIES} & \\
 & \multicolumn{2}{c}{} & \\
 & \multicolumn{2}{c}{\Large\bfseries COURSE PROGRAM} & \\
\end{tabular}
\end{center}

\begin{center}
\begin{longtable}{|p{4cm}|p{4cm}|p{4cm}|p{4cm}|}
\hline
\multicolumn{3}{|p{12cm}|}{Course Name in English:} &
\multicolumn{1}{p{4cm}|}{Code:} \\ 
\multicolumn{3}{|p{12cm}|}{\textbf{Geoprocessing and Geographic Information Systems}} &
\textbf{BCC593}\\ 
\multicolumn{3}{|p{12cm}|}{Course Name in Portuguese:} & \\ 
\multicolumn{3}{|p{12cm}|}{\textbf{Geoprocessamento e Sistemas de Informação Geográfica}} & \\ 
\hline
\multicolumn{3}{|p{12cm}|}{Department Name and Acronym:} & Academic Unit: \\ 
\multicolumn{3}{|p{12cm}|}{Department of Computer Science (DECOM)} & {ICEB} \\ 
\hline
\multicolumn{4}{|p{16cm}|}{Delivery Mode:
[X] in person \hspace{0.6cm}
[ ] blended \hspace{0.6cm}
[ ] distance learning}\\
\hline
\multicolumn{2}{|p{8cm}|}{Semester workload} &
\multicolumn{2}{p{8cm}|}{Weekly workload}\\
\hline
\multicolumn{1}{|p{4cm}|}{Total} &
\multicolumn{1}{p{4cm}|}{Extension} &
\multicolumn{1}{p{4cm}|}{Theoretical} &
\multicolumn{1}{p{4cm}|}{Practical} \\ 
\multicolumn{1}{|p{4cm}|}{60\,hours} &
\multicolumn{1}{p{4cm}|}{0\;hours} &
\multicolumn{1}{p{4cm}|}{4\;class hours} &
\multicolumn{1}{p{4cm}|}{0\;class hours} \\ 
\hline
\multicolumn{4}{|p{16cm}|}{Course Description:}\\
\multicolumn{4}{|p{16cm}|}{}\\
\multicolumn{4}{|p{\dimexpr 16cm + 6\tabcolsep\relax}|}{Computational representations of geographic space; introduction to cartography for geoinformatics; introduction to Geographic Information Systems; introduction to numerical terrain modeling; introduction to the global positioning system; introduction to remote sensing; introduction to Geographic Database modeling; introduction to Geographic Data Analysis; applications in Geoinformatics.}\\
\multicolumn{4}{|p{16cm}|}{}\\
\hline
\multicolumn{4}{|p{16cm}|}{Course Content:}\\*
\multicolumn{4}{|p{\dimexpr 16cm + 6\tabcolsep\relax}|}{\begin{enumerate}[start=1]\item Computational representations of geographic space: mathematical-computational modeling concepts, time, space, scale (extent and resolution) and processes of change; conceptual field and object models; raster data structures (numerical grids and images) and vector data structures (points, lines, and polygons)\end{enumerate}}\\
\multicolumn{4}{|p{\dimexpr 16cm + 6\tabcolsep\relax}|}{\begin{enumerate}[start=2]\item Introduction to Cartography for Geoinformatics: datum, cartographic projections, geographic coordinate systems, standardized geographic reference systems, cartographic scale\end{enumerate}}\\
\multicolumn{4}{|p{\dimexpr 16cm + 6\tabcolsep\relax}|}{\begin{enumerate}[start=3]\item Introduction to Geographic Information Systems: architecture (dual and integrated), typology (desktop, client-server, web), use of free tools (TerraView and QGIS); data layers, views, themes, and legends; vector and raster data acquisition; interoperability (DXF, SHP, GML); alphanumeric and spatial queries\end{enumerate}}\\
\multicolumn{4}{|p{\dimexpr 16cm + 6\tabcolsep\relax}|}{\begin{enumerate}[start=4]\item Introduction to Remote Sensing: concepts of satellites and remote sensors; types of imagery (optical, radar, and thermal); bands and false-color composition (RGB); image registration and georeferencing; digital image processing techniques (contrast, resampling, filtering, segmentation, and classification)\end{enumerate}}\\
\multicolumn{4}{|p{\dimexpr 16cm + 6\tabcolsep\relax}|}{\begin{enumerate}[start=5]\item Introduction to Geographic Database Modeling: OMT-G notation; implementation of geographic databases using free tools; spatial queries (reports) in free tools; spatial data indexing (Quadtree and R-tree)\end{enumerate}}\\
\multicolumn{4}{|p{\dimexpr 16cm + 6\tabcolsep\relax}|}{\begin{enumerate}[start=6]\item Introduction to Numerical Terrain Modeling: digital elevation model (DEM) and triangulated irregular networks (TIN); construction of digital terrain models from contour lines, drainage lines, and ridge lines\end{enumerate}}\\
\multicolumn{4}{|p{\dimexpr 16cm + 6\tabcolsep\relax}|}{\begin{enumerate}[start=7]\item Introduction to the Global Positioning System: GPS and Galileo systems; calibration and use of GPS equipment; sources of error and lack of precision; types of GPS equipment (navigation, mapping, differential)\end{enumerate}}\\
\multicolumn{4}{|p{\dimexpr 16cm + 6\tabcolsep\relax}|}{\begin{enumerate}[start=8]\item Introduction to Geographic Data Analysis: spatial analysis (point and area data), geostatistics (ordinary kriging), and map algebra (Tomlin and Egenhofer operators)\end{enumerate}}\\
\multicolumn{4}{|p{\dimexpr 16cm + 6\tabcolsep\relax}|}{\begin{enumerate}[start=9]\item Applications in Geoinformatics: land use and land cover change, disease spread, and multipurpose technical cadastre\end{enumerate}}\\
\multicolumn{4}{|p{\dimexpr 16cm + 6\tabcolsep\relax}|}{}\\
\hline
\multicolumn{4}{|p{16cm}|}{Basic Bibliography:}\\*
\multicolumn{4}{|p{\dimexpr 16cm + 6\tabcolsep\relax}|}{\begin{itemize}\item BLASCHKE, Thomas; KUX, Hermann (org.); tradução Hermann Kux. Sensoriamento remoto e SIG avançados: novos sistemas sensores, métodos inovadores. São Paulo: Oficina de Textos, c2005.\end{itemize}}\\
\multicolumn{4}{|p{\dimexpr 16cm + 6\tabcolsep\relax}|}{\begin{itemize}\item Anatomia de sistemas de informação geográfica. Campinas: UNICAMP, 1996.\end{itemize}}\\
\multicolumn{4}{|p{\dimexpr 16cm + 6\tabcolsep\relax}|}{\begin{itemize}\item GONZALEZ, Rafael C. Digital image processing. 3rd ed. Upper Saddle River, NJ: Pearson Prentice Hall, c2008.\end{itemize}}\\
\multicolumn{4}{|p{\dimexpr 16cm + 6\tabcolsep\relax}|}{\begin{itemize}\item JENSEN, John R.; tradução José Carlos Neves Epiphanio ... [et al.]. Sensoriamento remoto do ambiente: uma perspectiva em recursos terrestres. São José dos Campos: Parêntese, 2011.\end{itemize}}\\
\multicolumn{4}{|p{\dimexpr 16cm + 6\tabcolsep\relax}|}{\begin{itemize}\item CÂMARA, Gilberto. Geoprocessamento para projetos ambientais. São José dos Campos: INPE, 1996.\end{itemize}}\\
\multicolumn{4}{|p{\dimexpr 16cm + 6\tabcolsep\relax}|}{\begin{itemize}\item Conceitos básicos de sistemas de informação geográfica e cartografia aplicados à saúde. Brasília: Organização Pan-Americana de Saúde, 2000.\end{itemize}}\\
\multicolumn{4}{|p{16cm}|}{}\\
\hline
\multicolumn{4}{|p{16cm}|}{Complementary Bibliography:}\\*
\multicolumn{4}{|p{\dimexpr 16cm + 6\tabcolsep\relax}|}{\begin{itemize}\item Bancos de dados geográficos. Curitiba: MundoGeo, 2005.\end{itemize}}\\
\multicolumn{4}{|p{\dimexpr 16cm + 6\tabcolsep\relax}|}{\begin{itemize}\item Análise espacial de dados geográficos. Planaltina, DF: EMBRAPA Cerrados, 2004.\end{itemize}}\\
\multicolumn{4}{|p{\dimexpr 16cm + 6\tabcolsep\relax}|}{\begin{itemize}\item CANTY, Morton J. Image analysis, classification and change detection in remote sensing: with algorithms for ENVI/IDL. Boca Raton, FL: CRC/Taylor \& Francis, 2007.\end{itemize}}\\
\multicolumn{4}{|p{\dimexpr 16cm + 6\tabcolsep\relax}|}{\begin{itemize}\item NOVO, Evlyn M. L. de Moraes. Sensoriamento remoto: princípios e aplicações. 2. ed. São Paulo: E. Blucher, c1992.\end{itemize}}\\
\multicolumn{4}{|p{\dimexpr 16cm + 6\tabcolsep\relax}|}{\begin{itemize}\item MOURA, Ana Clara Mourão. Geoprocessamento na gestão e planejamento urbano. Belo Horizonte: Difusora, 2003.\end{itemize}}\\
\hline
\end{longtable}
\end{center}

\clearpage
\begin{center}
\begin{tabular}{lccr}
 \multirow{3}{*}{\includegraphics[height=2.7cm]{brasao.png}} &
 \multicolumn{2}{c}{\bfseries FEDERAL UNIVERSITY OF OURO PRETO} &
 \ \ \ \ \multirow{3}{*}{\includegraphics[height=2.7cm]{ufop.png}} \\
 & \multicolumn{2}{c}{\bfseries DEAN'S OFFICE FOR UNDERGRADUATE STUDIES} & \\
 & \multicolumn{2}{c}{} & \\
 & \multicolumn{2}{c}{\Large\bfseries COURSE PROGRAM} & \\
\end{tabular}
\end{center}

\begin{center}
\begin{longtable}{|p{4cm}|p{4cm}|p{4cm}|p{4cm}|}
\hline
\multicolumn{3}{|p{12cm}|}{Course Name in English:} &
\multicolumn{1}{p{4cm}|}{Code:} \\ 
\multicolumn{3}{|p{12cm}|}{\textbf{Modeling and Simulation of Land Systems}} &
\textbf{BCC595}\\ 
\multicolumn{3}{|p{12cm}|}{Course Name in Portuguese:} & \\ 
\multicolumn{3}{|p{12cm}|}{\textbf{Modelagem e Simulação de Sistemas Terrestres}} & \\ 
\hline
\multicolumn{3}{|p{12cm}|}{Department Name and Acronym:} & Academic Unit: \\ 
\multicolumn{3}{|p{12cm}|}{Department of Computer Science (DECOM)} & {ICEB} \\ 
\hline
\multicolumn{4}{|p{16cm}|}{Delivery Mode:
[X] in person \hspace{0.6cm}
[ ] blended \hspace{0.6cm}
[ ] distance learning}\\
\hline
\multicolumn{2}{|p{8cm}|}{Semester workload} &
\multicolumn{2}{p{8cm}|}{Weekly workload}\\
\hline
\multicolumn{1}{|p{4cm}|}{Total} &
\multicolumn{1}{p{4cm}|}{Extension} &
\multicolumn{1}{p{4cm}|}{Theoretical} &
\multicolumn{1}{p{4cm}|}{Practical} \\ 
\multicolumn{1}{|p{4cm}|}{60\,hours} &
\multicolumn{1}{p{4cm}|}{0\;hours} &
\multicolumn{1}{p{4cm}|}{4\;class hours} &
\multicolumn{1}{p{4cm}|}{0\;class hours} \\ 
\hline
\multicolumn{4}{|p{16cm}|}{Course Description:}\\
\multicolumn{4}{|p{16cm}|}{}\\
\multicolumn{4}{|p{\dimexpr 16cm + 6\tabcolsep\relax}|}{Introduction to computational modeling of geographic phenomena. Theoretical foundations of Earth system science. Concepts and computational representations of scale, space, time, and behavior. Introduction to spatially explicit dynamic models. Introduction to computational simulation. Introduction to Earth system modeling paradigms. Free tools for Earth system modeling. Multi-scale Earth system modeling. Applications in Earth system modeling and simulation. Mathematical methods applied to Earth system modeling.}\\
\multicolumn{4}{|p{16cm}|}{}\\
\hline
\multicolumn{4}{|p{16cm}|}{Course Content:}\\*
\multicolumn{4}{|p{\dimexpr 16cm + 6\tabcolsep\relax}|}{\begin{enumerate}[start=1]\item Introduction to Computational Modeling of Geographic Phenomena\end{enumerate}}\\
\multicolumn{4}{|p{\dimexpr 16cm + 6\tabcolsep\relax}|}{\begin{enumerate}[start=2]\item Theoretical foundations of Earth system science: mathematical-computational models, the modeling process, process simulation, energy stocks and flows, feedbacks, Earth system science\end{enumerate}}\\
\multicolumn{4}{|p{\dimexpr 16cm + 6\tabcolsep\relax}|}{\begin{enumerate}[start=3]\item Spatially explicit dynamic models: definition, taxonomy, development methodology, verification, calibration, validation\end{enumerate}}\\
\multicolumn{4}{|p{\dimexpr 16cm + 6\tabcolsep\relax}|}{\begin{enumerate}[start=4]\item Introduction to computational simulation: discrete simulation and discrete-event simulation, advantages and disadvantages, implementation of simulators\end{enumerate}}\\
\multicolumn{4}{|p{\dimexpr 16cm + 6\tabcolsep\relax}|}{\begin{enumerate}[start=5]\item Concepts and computational representation of scale, space, time, and behavior\end{enumerate}}\\
\multicolumn{4}{|p{\dimexpr 16cm + 6\tabcolsep\relax}|}{\begin{enumerate}[start=6]\item Earth system modeling paradigms: system dynamics, cellular automata, multi-agent systems\end{enumerate}}\\
\multicolumn{4}{|p{\dimexpr 16cm + 6\tabcolsep\relax}|}{\begin{enumerate}[start=7]\item Free tools for environmental modeling: Vensim, NetLogo, Repast, TerraME\end{enumerate}}\\
\multicolumn{4}{|p{\dimexpr 16cm + 6\tabcolsep\relax}|}{\begin{enumerate}[start=8]\item Multi-scale Earth system modeling: software requirements, integration with geographic information systems, development platforms\end{enumerate}}\\
\multicolumn{4}{|p{\dimexpr 16cm + 6\tabcolsep\relax}|}{\begin{enumerate}[start=9]\item Applications: hydrological models, climate models, population dynamics models, land use and land cover change models, species dispersal models, atmospheric models, epidemiological models, econometric models\end{enumerate}}\\
\multicolumn{4}{|p{\dimexpr 16cm + 6\tabcolsep\relax}|}{\begin{enumerate}[start=10]\item Mathematical methods applied to Earth system modeling: statistical inference, probability, and spatial analysis\end{enumerate}}\\
\multicolumn{4}{|p{\dimexpr 16cm + 6\tabcolsep\relax}|}{}\\
\hline
\multicolumn{4}{|p{16cm}|}{Basic Bibliography:}\\*
\multicolumn{4}{|p{\dimexpr 16cm + 6\tabcolsep\relax}|}{\begin{itemize}\item AXELROD, Robert M. The complexity of cooperation: agent-based models of competition and collaboration. Princeton, N.J.: Princeton University Press, 1997. (item inexistente na UFOP)\end{itemize}}\\
\multicolumn{4}{|p{\dimexpr 16cm + 6\tabcolsep\relax}|}{\begin{itemize}\item GOODCHILD, Michael F.; STEYAERT, Louis T.; PARKS, Bradley O.; JOHNSTON, Carol ... [et al.]. GIS and environmental modeling: progress and research issues. Fort Collins: GIS World Books, c1996.\end{itemize}}\\
\multicolumn{4}{|p{\dimexpr 16cm + 6\tabcolsep\relax}|}{\begin{itemize}\item WAINWRIGHT, John; MULLIGAN, Mark (eds.). Environmental modelling: finding simplicity in complexity. Chichester, West Sussex, England.\end{itemize}}\\
\multicolumn{4}{|p{\dimexpr 16cm + 6\tabcolsep\relax}|}{\begin{itemize}\item ZEIGLER, Bernard P.; PRAEHOFER, Herbert; KIM, Tag Gon. Theory of modeling and simulation: integrating discrete event and continuous complex dynamic systems. 2nd ed. San Diego: Academic, c2000.\end{itemize}}\\
\multicolumn{4}{|p{\dimexpr 16cm + 6\tabcolsep\relax}|}{\begin{itemize}\item FORD, Andrew. Modeling the environment: an introduction to system dynamics models of environmental systems. 2. ed. Washington, D.C.: Island Press, c1999.\end{itemize}}\\
\multicolumn{4}{|p{\dimexpr 16cm + 6\tabcolsep\relax}|}{\begin{itemize}\item EPSTEIN, Joshua M. Generative social science: studies in agent-based computational modeling. Princeton, NJ.\end{itemize}}\\
\multicolumn{4}{|p{\dimexpr 16cm + 6\tabcolsep\relax}|}{\begin{itemize}\item Integrating geographic information systems and agent-based modeling techniques for simulating social and ecological processes. Oxford.\end{itemize}}\\
\multicolumn{4}{|p{\dimexpr 16cm + 6\tabcolsep\relax}|}{\begin{itemize}\item MEADOWS, Donella H.; ed. Diana Wright. Thinking in systems: a primer. White River Junction, Vt.: Chelsea Green Pub., c2008.\end{itemize}}\\
\multicolumn{4}{|p{\dimexpr 16cm + 6\tabcolsep\relax}|}{\begin{itemize}\item LANTUÉJOUL, Christian. Geostatistical simulation: models and algorithms. Berlin: Springer-Verlag, 2002.\end{itemize}}\\
\multicolumn{4}{|p{16cm}|}{}\\
\hline
\multicolumn{4}{|p{16cm}|}{Complementary Bibliography:}\\*
\multicolumn{4}{|p{\dimexpr 16cm + 6\tabcolsep\relax}|}{\begin{itemize}\item BATTY, Michael. Cities and complexity: understanding cities with cellular automata, agent-based models, and fractals. Cambridge, Mass.: MIT Press, 2005. (item inexistente na UFOP)\end{itemize}}\\
\multicolumn{4}{|p{\dimexpr 16cm + 6\tabcolsep\relax}|}{\begin{itemize}\item SKIDMORE, Andrew. Environmental modelling with GIS and remote sensing. London: Taylor \& Francis, 2002. (item inexistente na UFOP)\end{itemize}}\\
\multicolumn{4}{|p{\dimexpr 16cm + 6\tabcolsep\relax}|}{\begin{itemize}\item SMITH, Jo U. Introduction to environmental modelling. Oxford: Oxford University Press, 2007.\end{itemize}}\\
\hline
\end{longtable}
\end{center}

\clearpage
\begin{center}
\begin{tabular}{lccr}
 \multirow{3}{*}{\includegraphics[height=2.7cm]{brasao.png}} &
 \multicolumn{2}{c}{\bfseries FEDERAL UNIVERSITY OF OURO PRETO} &
 \ \ \ \ \multirow{3}{*}{\includegraphics[height=2.7cm]{ufop.png}} \\
 & \multicolumn{2}{c}{\bfseries DEAN'S OFFICE FOR UNDERGRADUATE STUDIES} & \\
 & \multicolumn{2}{c}{} & \\
 & \multicolumn{2}{c}{\Large\bfseries COURSE PROGRAM} & \\
\end{tabular}
\end{center}

\begin{center}
\begin{longtable}{|p{4cm}|p{4cm}|p{4cm}|p{4cm}|}
\hline
\multicolumn{3}{|p{12cm}|}{Course Name in English:} &
\multicolumn{1}{p{4cm}|}{Code:} \\ 
\multicolumn{3}{|p{12cm}|}{\textbf{Logic Programming}} &
\textbf{BCC596}\\ 
\multicolumn{3}{|p{12cm}|}{Course Name in Portuguese:} & \\ 
\multicolumn{3}{|p{12cm}|}{\textbf{Programação em Lógica}} & \\ 
\hline
\multicolumn{3}{|p{12cm}|}{Department Name and Acronym:} & Academic Unit: \\ 
\multicolumn{3}{|p{12cm}|}{Department of Computer Science (DECOM)} & {ICEB} \\ 
\hline
\multicolumn{4}{|p{16cm}|}{Delivery Mode:
[X] in person \hspace{0.6cm}
[ ] blended \hspace{0.6cm}
[ ] distance learning}\\
\hline
\multicolumn{2}{|p{8cm}|}{Semester workload} &
\multicolumn{2}{p{8cm}|}{Weekly workload}\\
\hline
\multicolumn{1}{|p{4cm}|}{Total} &
\multicolumn{1}{p{4cm}|}{Extension} &
\multicolumn{1}{p{4cm}|}{Theoretical} &
\multicolumn{1}{p{4cm}|}{Practical} \\ 
\multicolumn{1}{|p{4cm}|}{60\,hours} &
\multicolumn{1}{p{4cm}|}{0\;hours} &
\multicolumn{1}{p{4cm}|}{4\;class hours} &
\multicolumn{1}{p{4cm}|}{0\;class hours} \\ 
\hline
\multicolumn{4}{|p{16cm}|}{Course Description:}\\
\multicolumn{4}{|p{16cm}|}{}\\
\multicolumn{4}{|p{\dimexpr 16cm + 6\tabcolsep\relax}|}{Theoretical foundations of logic programming: unification and resolution; concepts of logic programming languages; study of a logic programming language; applications.}\\
\multicolumn{4}{|p{16cm}|}{}\\
\hline
\multicolumn{4}{|p{16cm}|}{Course Content:}\\*
\multicolumn{4}{|p{\dimexpr 16cm + 6\tabcolsep\relax}|}{\begin{enumerate}[start=1]\item Declarative programming\end{enumerate}}\\
\multicolumn{4}{|p{\dimexpr 16cm + 6\tabcolsep\relax}|}{\begin{enumerate}[start=2]\item Introduction to Prolog\end{enumerate}}\\
\multicolumn{4}{|p{\dimexpr 16cm + 6\tabcolsep\relax}|}{\begin{enumerate}[start=3]\item Syntax and Semantics of Prolog programs\end{enumerate}}\\
\multicolumn{4}{|p{\dimexpr 16cm + 6\tabcolsep\relax}|}{\begin{enumerate}[start=4]\item Lists, operators, arithmetic\end{enumerate}}\\
\multicolumn{4}{|p{\dimexpr 16cm + 6\tabcolsep\relax}|}{\begin{enumerate}[start=5]\item Structures\end{enumerate}}\\
\multicolumn{4}{|p{\dimexpr 16cm + 6\tabcolsep\relax}|}{\begin{enumerate}[start=6]\item Backtracking\end{enumerate}}\\
\multicolumn{4}{|p{\dimexpr 16cm + 6\tabcolsep\relax}|}{\begin{enumerate}[start=7]\item Input and output\end{enumerate}}\\
\multicolumn{4}{|p{\dimexpr 16cm + 6\tabcolsep\relax}|}{\begin{enumerate}[start=8]\item Predefined predicates\end{enumerate}}\\
\multicolumn{4}{|p{\dimexpr 16cm + 6\tabcolsep\relax}|}{\begin{enumerate}[start=9]\item Programming style and techniques\end{enumerate}}\\
\multicolumn{4}{|p{\dimexpr 16cm + 6\tabcolsep\relax}|}{\begin{enumerate}[start=10]\item Operations on data structures\end{enumerate}}\\
\multicolumn{4}{|p{\dimexpr 16cm + 6\tabcolsep\relax}|}{\begin{enumerate}[start=11]\item Prolog's relation to Logic: predicate calculus, clausal form, resolution, theorem proving, and Horn clauses\end{enumerate}}\\
\multicolumn{4}{|p{\dimexpr 16cm + 6\tabcolsep\relax}|}{}\\
\hline
\multicolumn{4}{|p{16cm}|}{Basic Bibliography:}\\*
\multicolumn{4}{|p{\dimexpr 16cm + 6\tabcolsep\relax}|}{\begin{itemize}\item BRATKO, Ivan. Prolog programming for artificial intelligence. 3rd ed. Harlow: Pearson Education, 2001.\end{itemize}}\\
\multicolumn{4}{|p{\dimexpr 16cm + 6\tabcolsep\relax}|}{\begin{itemize}\item CASANOVA, Marco Antonio. Programação em lógica e a linguagem PROLOG. São Paulo: E. Blucher, c1987.\end{itemize}}\\
\multicolumn{4}{|p{\dimexpr 16cm + 6\tabcolsep\relax}|}{\begin{itemize}\item CLOCKSIN, W. F. Programming in Prolog. 4th ed. Berlin: Springer-Verlag, c1994.\end{itemize}}\\
\multicolumn{4}{|p{16cm}|}{}\\
\hline
\multicolumn{4}{|p{16cm}|}{Complementary Bibliography:}\\*
\multicolumn{4}{|p{\dimexpr 16cm + 6\tabcolsep\relax}|}{\begin{itemize}\item O'KEEFE, Richard. The Craft of Prolog. MIT Press, 2009. (item inexistente na UFOP)\end{itemize}}\\
\multicolumn{4}{|p{\dimexpr 16cm + 6\tabcolsep\relax}|}{\begin{itemize}\item BRAMER, Max. Logic Programming with Prolog. 2. ed. Springer, 2014. (item inexistente na UFOP)\end{itemize}}\\
\multicolumn{4}{|p{\dimexpr 16cm + 6\tabcolsep\relax}|}{\begin{itemize}\item MAIER, David; WARREN, David S. Computing with logic: logic programming with prolog. Menlo Park, Calif.: Benjamin/Cummings, c1988.\end{itemize}}\\
\multicolumn{4}{|p{\dimexpr 16cm + 6\tabcolsep\relax}|}{\begin{itemize}\item DERANSART, Pierre; ED-DBALI, AbdelAli; CERVONI, Laurent. Prolog: the standard: reference manual. Berlin: Springer-Verlag, 1996.\end{itemize}}\\
\multicolumn{4}{|p{\dimexpr 16cm + 6\tabcolsep\relax}|}{\begin{itemize}\item STERLING, Leon; SHAPIRO, Ehud; foreword by David H. D. Warren. The art of prolog: advanced programming techniques. 2nd ed. Cambridge: The MIT Press, 1994.\end{itemize}}\\
\hline
\end{longtable}
\end{center}

\clearpage
\begin{center}
\begin{tabular}{lccr}
 \multirow{3}{*}{\includegraphics[height=2.7cm]{brasao.png}} &
 \multicolumn{2}{c}{\bfseries FEDERAL UNIVERSITY OF OURO PRETO} &
 \ \ \ \ \multirow{3}{*}{\includegraphics[height=2.7cm]{ufop.png}} \\
 & \multicolumn{2}{c}{\bfseries DEAN'S OFFICE FOR UNDERGRADUATE STUDIES} & \\
 & \multicolumn{2}{c}{} & \\
 & \multicolumn{2}{c}{\Large\bfseries COURSE PROGRAM} & \\
\end{tabular}
\end{center}

\begin{center}
\begin{longtable}{|p{4cm}|p{4cm}|p{4cm}|p{4cm}|}
\hline
\multicolumn{3}{|p{12cm}|}{Course Name in English:} &
\multicolumn{1}{p{4cm}|}{Code:} \\ 
\multicolumn{3}{|p{12cm}|}{\textbf{Web Information Retrieval}} &
\textbf{BCC599}\\ 
\multicolumn{3}{|p{12cm}|}{Course Name in Portuguese:} & \\ 
\multicolumn{3}{|p{12cm}|}{\textbf{Recuperação de Informação na Web}} & \\ 
\hline
\multicolumn{3}{|p{12cm}|}{Department Name and Acronym:} & Academic Unit: \\ 
\multicolumn{3}{|p{12cm}|}{Department of Computer Science (DECOM)} & {ICEB} \\ 
\hline
\multicolumn{4}{|p{16cm}|}{Delivery Mode:
[X] in person \hspace{0.6cm}
[ ] blended \hspace{0.6cm}
[ ] distance learning}\\
\hline
\multicolumn{2}{|p{8cm}|}{Semester workload} &
\multicolumn{2}{p{8cm}|}{Weekly workload}\\
\hline
\multicolumn{1}{|p{4cm}|}{Total} &
\multicolumn{1}{p{4cm}|}{Extension} &
\multicolumn{1}{p{4cm}|}{Theoretical} &
\multicolumn{1}{p{4cm}|}{Practical} \\ 
\multicolumn{1}{|p{4cm}|}{60\,hours} &
\multicolumn{1}{p{4cm}|}{0\;hours} &
\multicolumn{1}{p{4cm}|}{4\;class hours} &
\multicolumn{1}{p{4cm}|}{0\;class hours} \\ 
\hline
\multicolumn{4}{|p{16cm}|}{Course Description:}\\
\multicolumn{4}{|p{16cm}|}{}\\
\multicolumn{4}{|p{\dimexpr 16cm + 6\tabcolsep\relax}|}{Modeling of information retrieval systems. Characterization of the web. Operation of crawling systems. Distributed indexing. Query languages. Query operations. Generation of ranking functions. Evaluation of results.}\\
\multicolumn{4}{|p{16cm}|}{}\\
\hline
\multicolumn{4}{|p{16cm}|}{Course Content:}\\*
\multicolumn{4}{|p{\dimexpr 16cm + 6\tabcolsep\relax}|}{\begin{enumerate}[start=1]\item Introduction\end{enumerate}}\\
\multicolumn{4}{|p{\dimexpr 16cm + 6\tabcolsep\relax}|}{\begin{enumerate}[start=2]\item Boolean retrieval\end{enumerate}}\\
\multicolumn{4}{|p{\dimexpr 16cm + 6\tabcolsep\relax}|}{\begin{enumerate}[start=3]\item Term vocabulary and inverted lists (postings)\end{enumerate}}\\
\multicolumn{4}{|p{\dimexpr 16cm + 6\tabcolsep\relax}|}{\begin{enumerate}[start=4]\item Dictionaries and tolerant retrieval\end{enumerate}}\\
\multicolumn{4}{|p{\dimexpr 16cm + 6\tabcolsep\relax}|}{\begin{enumerate}[start=5]\item Index construction\end{enumerate}}\\
\multicolumn{4}{|p{\dimexpr 16cm + 6\tabcolsep\relax}|}{\begin{enumerate}[start=6]\item Index compression\end{enumerate}}\\
\multicolumn{4}{|p{\dimexpr 16cm + 6\tabcolsep\relax}|}{\begin{enumerate}[start=7]\item Scoring, term weighting, and the vector space model\end{enumerate}}\\
\multicolumn{4}{|p{\dimexpr 16cm + 6\tabcolsep\relax}|}{\begin{enumerate}[start=8]\item Computing scores in a complete search system\end{enumerate}}\\
\multicolumn{4}{|p{\dimexpr 16cm + 6\tabcolsep\relax}|}{\begin{enumerate}[start=9]\item Evaluation in information retrieval\end{enumerate}}\\
\multicolumn{4}{|p{\dimexpr 16cm + 6\tabcolsep\relax}|}{\begin{enumerate}[start=10]\item Relevance feedback and query expansion\end{enumerate}}\\
\multicolumn{4}{|p{\dimexpr 16cm + 6\tabcolsep\relax}|}{\begin{enumerate}[start=11]\item Introduction to Web search\end{enumerate}}\\
\multicolumn{4}{|p{\dimexpr 16cm + 6\tabcolsep\relax}|}{\begin{enumerate}[start=12]\item Web crawling and indexing\end{enumerate}}\\
\multicolumn{4}{|p{\dimexpr 16cm + 6\tabcolsep\relax}|}{\begin{enumerate}[start=13]\item Link analysis\end{enumerate}}\\
\multicolumn{4}{|p{\dimexpr 16cm + 6\tabcolsep\relax}|}{}\\
\hline
\multicolumn{4}{|p{16cm}|}{Basic Bibliography:}\\*
\multicolumn{4}{|p{\dimexpr 16cm + 6\tabcolsep\relax}|}{\begin{itemize}\item MANNING, Christopher D. Introduction to information retrieval. New York: Cambridge University Press, 2008.\end{itemize}}\\
\multicolumn{4}{|p{\dimexpr 16cm + 6\tabcolsep\relax}|}{\begin{itemize}\item LIU, Bing. Web data mining: exploring hyperlinks, contents, and usage data. Berlin.\end{itemize}}\\
\multicolumn{4}{|p{\dimexpr 16cm + 6\tabcolsep\relax}|}{\begin{itemize}\item BAEZA-YATES, Ricardo; RIBEIRO-NETO, Berthier; tradução Leandro Krug Wives, Viviane Pereira Moreira. Recuperação de informação: conceitos e tecnologia das máquinas de busca. 2. ed. Porto Alegre: Bookman, 2013.\end{itemize}}\\
\multicolumn{4}{|p{16cm}|}{}\\
\hline
\multicolumn{4}{|p{16cm}|}{Complementary Bibliography:}\\*
\multicolumn{4}{|p{\dimexpr 16cm + 6\tabcolsep\relax}|}{\begin{itemize}\item GOLDSCHMIDT, Ronaldo; PASSOS, Emmanuel. Data Mining: um guia prático: conceitos, técnicas, ferramentas, orientações e aplicações. Rio de Janeiro: Elsevier.\end{itemize}}\\
\multicolumn{4}{|p{\dimexpr 16cm + 6\tabcolsep\relax}|}{\begin{itemize}\item CHAKRABARTI, Soumen. Mining the Web: discovering knowledge from hypertext data. San Francisco, CA: Morgan Kaufmann Publishers, c2003.\end{itemize}}\\
\multicolumn{4}{|p{\dimexpr 16cm + 6\tabcolsep\relax}|}{\begin{itemize}\item HAN, Jiawei; KAMBER, Micheline; PEI, Jian. Data mining: concepts and techniques. 3. ed. Boston: Elsevier, 2012.\end{itemize}}\\
\multicolumn{4}{|p{\dimexpr 16cm + 6\tabcolsep\relax}|}{\begin{itemize}\item WITTEN, Ian H.; MOFFAT, Alistair; BELL, Timothy C. Managing gigabytes: compressing and indexing documents and images. 2. ed. San Francisco, Calif.: Morgan Kaufmann Publishers, 1999.\end{itemize}}\\
\multicolumn{4}{|p{\dimexpr 16cm + 6\tabcolsep\relax}|}{\begin{itemize}\item HEAPS, H. S. Information retrieval, computational and theoretical aspects. New York.\end{itemize}}\\
\hline
\end{longtable}
\end{center}

\clearpage
\begin{center}
\begin{tabular}{lccr}
 \multirow{3}{*}{\includegraphics[height=2.7cm]{brasao.png}} &
 \multicolumn{2}{c}{\bfseries FEDERAL UNIVERSITY OF OURO PRETO} &
 \ \ \ \ \multirow{3}{*}{\includegraphics[height=2.7cm]{ufop.png}} \\
 & \multicolumn{2}{c}{\bfseries DEAN'S OFFICE FOR UNDERGRADUATE STUDIES} & \\
 & \multicolumn{2}{c}{} & \\
 & \multicolumn{2}{c}{\Large\bfseries COURSE PROGRAM} & \\
\end{tabular}
\end{center}

\begin{center}
\begin{longtable}{|p{4cm}|p{4cm}|p{4cm}|p{4cm}|}
\hline
\multicolumn{3}{|p{12cm}|}{Course Name in English:} &
\multicolumn{1}{p{4cm}|}{Code:} \\ 
\multicolumn{3}{|p{12cm}|}{\textbf{Web Data Management}} &
\textbf{BCC610}\\ 
\multicolumn{3}{|p{12cm}|}{Course Name in Portuguese:} & \\ 
\multicolumn{3}{|p{12cm}|}{\textbf{Gerência de Dados na Web}} & \\ 
\hline
\multicolumn{3}{|p{12cm}|}{Department Name and Acronym:} & Academic Unit: \\ 
\multicolumn{3}{|p{12cm}|}{Department of Computer Science (DECOM)} & {ICEB} \\ 
\hline
\multicolumn{4}{|p{16cm}|}{Delivery Mode:
[X] in person \hspace{0.6cm}
[ ] blended \hspace{0.6cm}
[ ] distance learning}\\
\hline
\multicolumn{2}{|p{8cm}|}{Semester workload} &
\multicolumn{2}{p{8cm}|}{Weekly workload}\\
\hline
\multicolumn{1}{|p{4cm}|}{Total} &
\multicolumn{1}{p{4cm}|}{Extension} &
\multicolumn{1}{p{4cm}|}{Theoretical} &
\multicolumn{1}{p{4cm}|}{Practical} \\ 
\multicolumn{1}{|p{4cm}|}{60\,hours} &
\multicolumn{1}{p{4cm}|}{0\;hours} &
\multicolumn{1}{p{4cm}|}{4\;class hours} &
\multicolumn{1}{p{4cm}|}{0\;class hours} \\ 
\hline
\multicolumn{4}{|p{16cm}|}{Course Description:}\\
\multicolumn{4}{|p{16cm}|}{}\\
\multicolumn{4}{|p{\dimexpr 16cm + 6\tabcolsep\relax}|}{Characterization of Web data; modeling of semi-structured data; query languages; extraction of Web data; integration of Web data; applications.}\\
\multicolumn{4}{|p{16cm}|}{}\\
\hline
\multicolumn{4}{|p{16cm}|}{Course Content:}\\*
\multicolumn{4}{|p{\dimexpr 16cm + 6\tabcolsep\relax}|}{\begin{enumerate}[start=1]\item Web data: characteristics of Web data, Web versus databases, architecture of Web applications\end{enumerate}}\\
\multicolumn{4}{|p{\dimexpr 16cm + 6\tabcolsep\relax}|}{\begin{enumerate}[start=2]\item Modeling semi-structured data and XML: characteristics of semi-structured data, models for representing semi-structured data, XML, relational representation of XML data, publishing relational data as XML, XML structure discovery, tools for XML data management, XML data compression\end{enumerate}}\\
\multicolumn{4}{|p{\dimexpr 16cm + 6\tabcolsep\relax}|}{\begin{enumerate}[start=3]\item Query languages: query languages for the Web, query languages for semi-structured data, XML query and manipulation languages\end{enumerate}}\\
\multicolumn{4}{|p{\dimexpr 16cm + 6\tabcolsep\relax}|}{\begin{enumerate}[start=4]\item Extracting data from Web sources: characterization of the problem, main approaches, tools, evaluation of extraction quality\end{enumerate}}\\
\multicolumn{4}{|p{\dimexpr 16cm + 6\tabcolsep\relax}|}{\begin{enumerate}[start=5]\item Web data integration: characterization of the integration problem, approaches and architectures, data fusion, similarity measures\end{enumerate}}\\
\multicolumn{4}{|p{\dimexpr 16cm + 6\tabcolsep\relax}|}{\begin{enumerate}[start=6]\item Advanced topics: query processing on XML documents, hidden Web data access, entity recognition and extraction, the Semantic Web, the Web of Data, Web page crawling, digital libraries\end{enumerate}}\\
\multicolumn{4}{|p{\dimexpr 16cm + 6\tabcolsep\relax}|}{}\\
\hline
\multicolumn{4}{|p{16cm}|}{Basic Bibliography:}\\*
\multicolumn{4}{|p{\dimexpr 16cm + 6\tabcolsep\relax}|}{\begin{itemize}\item GRAVES, Mark. Projeto de banco de dados com XML. São Paulo: Pearson Education, 2003.\end{itemize}}\\
\multicolumn{4}{|p{\dimexpr 16cm + 6\tabcolsep\relax}|}{\begin{itemize}\item BREITMAN, Karin Koogan. Web semântica: a internet do futuro. Rio de Janeiro: LTC, 2006.\end{itemize}}\\
\multicolumn{4}{|p{\dimexpr 16cm + 6\tabcolsep\relax}|}{\begin{itemize}\item BAEZA-YATES, R. Modern information retrieval. New York: Addison Wesley, 2011.\end{itemize}}\\
\multicolumn{4}{|p{\dimexpr 16cm + 6\tabcolsep\relax}|}{\begin{itemize}\item ABITEBOUL, S. Gerenciando dados na WEB. Rio de Janeiro: Campus, 2000.\end{itemize}}\\
\multicolumn{4}{|p{16cm}|}{}\\
\hline
\multicolumn{4}{|p{16cm}|}{Complementary Bibliography:}\\*
\multicolumn{4}{|p{\dimexpr 16cm + 6\tabcolsep\relax}|}{\begin{itemize}\item ANTONIOU, G. A semantic web primer. Cambridge: MIT Press, 2008.\end{itemize}}\\
\multicolumn{4}{|p{\dimexpr 16cm + 6\tabcolsep\relax}|}{\begin{itemize}\item ELMASRI, Ramez. Sistemas de banco de dados. 6. ed. São Paulo: Pearson.\end{itemize}}\\
\multicolumn{4}{|p{\dimexpr 16cm + 6\tabcolsep\relax}|}{\begin{itemize}\item KANTARDZIC, Mehmed. Data mining: concepts, models, methods, and algorithms. Hoboken, NJ: Wiley-Interscience.\end{itemize}}\\
\multicolumn{4}{|p{\dimexpr 16cm + 6\tabcolsep\relax}|}{\begin{itemize}\item LIU, Bing. Web data mining: exploring hyperlinks, contents, and usage data. Berlin.\end{itemize}}\\
\multicolumn{4}{|p{\dimexpr 16cm + 6\tabcolsep\relax}|}{\begin{itemize}\item GARCIA-MOLINA, Hector; ULLMAN, Jeffrey D.; WIDOM, Jennifer. Database systems: the complete book. 2nd ed. Upper Saddle River: Pearson Prentice Hall, c2009.\end{itemize}}\\
\hline
\end{longtable}
\end{center}

\clearpage
\begin{center}
\begin{tabular}{lccr}
 \multirow{3}{*}{\includegraphics[height=2.7cm]{brasao.png}} &
 \multicolumn{2}{c}{\bfseries FEDERAL UNIVERSITY OF OURO PRETO} &
 \ \ \ \ \multirow{3}{*}{\includegraphics[height=2.7cm]{ufop.png}} \\
 & \multicolumn{2}{c}{\bfseries DEAN'S OFFICE FOR UNDERGRADUATE STUDIES} & \\
 & \multicolumn{2}{c}{} & \\
 & \multicolumn{2}{c}{\Large\bfseries COURSE PROGRAM} & \\
\end{tabular}
\end{center}

\begin{center}
\begin{longtable}{|p{4cm}|p{4cm}|p{4cm}|p{4cm}|}
\hline
\multicolumn{3}{|p{12cm}|}{Course Name in English:} &
\multicolumn{1}{p{4cm}|}{Code:} \\ 
\multicolumn{3}{|p{12cm}|}{\textbf{Evolutionary Computing}} &
\textbf{BCC541}\\ 
\multicolumn{3}{|p{12cm}|}{Course Name in Portuguese:} & \\ 
\multicolumn{3}{|p{12cm}|}{\textbf{Computação Evolutiva}} & \\ 
\hline
\multicolumn{3}{|p{12cm}|}{Department Name and Acronym:} & Academic Unit: \\ 
\multicolumn{3}{|p{12cm}|}{Department of Computer Science (DECOM)} & {ICEB} \\ 
\hline
\multicolumn{4}{|p{16cm}|}{Delivery Mode:
[X] in person \hspace{0.6cm}
[ ] blended \hspace{0.6cm}
[ ] distance learning}\\
\hline
\multicolumn{2}{|p{8cm}|}{Semester workload} &
\multicolumn{2}{p{8cm}|}{Weekly workload}\\
\hline
\multicolumn{1}{|p{4cm}|}{Total} &
\multicolumn{1}{p{4cm}|}{Extension} &
\multicolumn{1}{p{4cm}|}{Theoretical} &
\multicolumn{1}{p{4cm}|}{Practical} \\ 
\multicolumn{1}{|p{4cm}|}{60\,hours} &
\multicolumn{1}{p{4cm}|}{0\;hours} &
\multicolumn{1}{p{4cm}|}{4\;class hours} &
\multicolumn{1}{p{4cm}|}{0\;class hours} \\ 
\hline
\multicolumn{4}{|p{16cm}|}{Course Description:}\\
\multicolumn{4}{|p{16cm}|}{}\\
\multicolumn{4}{|p{\dimexpr 16cm + 6\tabcolsep\relax}|}{Introduction to evolutionary computing; evolution by natural selection; basic concepts and components of evolutionary algorithms; genetic algorithms; evolution strategies; evolutionary programming; genetic programming; classifier systems; hybrid algorithms; constraint handling; multi-objective optimization with evolutionary algorithms.}\\
\multicolumn{4}{|p{16cm}|}{}\\
\hline
\multicolumn{4}{|p{16cm}|}{Course Content:}\\*
\multicolumn{4}{|p{\dimexpr 16cm + 6\tabcolsep\relax}|}{\begin{enumerate}[start=1]\item Introduction to evolutionary computing and its relation to biological models\end{enumerate}}\\
\multicolumn{4}{|p{\dimexpr 16cm + 6\tabcolsep\relax}|}{\begin{enumerate}[start=2]\item Basic components common to evolutionary algorithms\end{enumerate}}\\
\multicolumn{4}{|p{\dimexpr 16cm + 6\tabcolsep\relax}|}{\begin{enumerate}[start=3]\item Categories of evolutionary algorithms\end{enumerate}}\\
\multicolumn{4}{|p{\dimexpr 16cm + 6\tabcolsep\relax}|}{\begin{enumerate}[start=4]\item Evolutionary systems applied to classification systems\end{enumerate}}\\
\multicolumn{4}{|p{\dimexpr 16cm + 6\tabcolsep\relax}|}{\begin{enumerate}[start=5]\item Hybrid algorithms\end{enumerate}}\\
\multicolumn{4}{|p{\dimexpr 16cm + 6\tabcolsep\relax}|}{\begin{enumerate}[start=6]\item Introduction to multi-objective optimization\end{enumerate}}\\
\multicolumn{4}{|p{\dimexpr 16cm + 6\tabcolsep\relax}|}{}\\
\hline
\multicolumn{4}{|p{16cm}|}{Basic Bibliography:}\\*
\multicolumn{4}{|p{\dimexpr 16cm + 6\tabcolsep\relax}|}{\begin{itemize}\item BÄCK, Thomas. Handbook of evolutionary computation. CRC Press, 1997.\end{itemize}}\\
\multicolumn{4}{|p{\dimexpr 16cm + 6\tabcolsep\relax}|}{\begin{itemize}\item David B. Fogel. Evolutionary computation: toward a new philosophy of machine intelligence. 3rd ed. New Jersey: Wiley Interscience, 2006.\end{itemize}}\\
\multicolumn{4}{|p{16cm}|}{}\\
\hline
\multicolumn{4}{|p{16cm}|}{Complementary Bibliography:}\\*
\multicolumn{4}{|p{\dimexpr 16cm + 6\tabcolsep\relax}|}{\begin{itemize}\item BACK, Thomas. Evolutionary algorithms in theory and practice. Oxford Univ. Press, 1996.\end{itemize}}\\
\multicolumn{4}{|p{\dimexpr 16cm + 6\tabcolsep\relax}|}{\begin{itemize}\item DEB, Kalyanmoy. Multi-objective optimization using evolutionary algorithms. John Wiley \& Sons, 2001.\end{itemize}}\\
\multicolumn{4}{|p{\dimexpr 16cm + 6\tabcolsep\relax}|}{\begin{itemize}\item YAO, Xin. Evolutionary computation: Theory and applications. World Scientific, 1999.\end{itemize}}\\
\multicolumn{4}{|p{\dimexpr 16cm + 6\tabcolsep\relax}|}{\begin{itemize}\item JOHNSON, R. P. Automatic generation of neural network architecture using evolutionary computation. World Scientific, 1997.\end{itemize}}\\
\multicolumn{4}{|p{\dimexpr 16cm + 6\tabcolsep\relax}|}{\begin{itemize}\item Kenneth A. De Jong. Evolutionary computation: a unified approach. Cambridge, Mass.: MIT Press, 2006.\end{itemize}}\\
\hline
\end{longtable}
\end{center}

\clearpage
\begin{center}
\begin{tabular}{lccr}
 \multirow{3}{*}{\includegraphics[height=2.7cm]{brasao.png}} &
 \multicolumn{2}{c}{\bfseries FEDERAL UNIVERSITY OF OURO PRETO} &
 \ \ \ \ \multirow{3}{*}{\includegraphics[height=2.7cm]{ufop.png}} \\
 & \multicolumn{2}{c}{\bfseries DEAN'S OFFICE FOR UNDERGRADUATE STUDIES} & \\
 & \multicolumn{2}{c}{} & \\
 & \multicolumn{2}{c}{\Large\bfseries COURSE PROGRAM} & \\
\end{tabular}
\end{center}

\begin{center}
\begin{longtable}{|p{4cm}|p{4cm}|p{4cm}|p{4cm}|}
\hline
\multicolumn{3}{|p{12cm}|}{Course Name in English:} &
\multicolumn{1}{p{4cm}|}{Code:} \\ 
\multicolumn{3}{|p{12cm}|}{\textbf{Network Optimization}} &
\textbf{BCC543}\\ 
\multicolumn{3}{|p{12cm}|}{Course Name in Portuguese:} & \\ 
\multicolumn{3}{|p{12cm}|}{\textbf{Otimização em Redes}} & \\ 
\hline
\multicolumn{3}{|p{12cm}|}{Department Name and Acronym:} & Academic Unit: \\ 
\multicolumn{3}{|p{12cm}|}{Department of Computer Science (DECOM)} & {ICEB} \\ 
\hline
\multicolumn{4}{|p{16cm}|}{Delivery Mode:
[X] in person \hspace{0.6cm}
[ ] blended \hspace{0.6cm}
[ ] distance learning}\\
\hline
\multicolumn{2}{|p{8cm}|}{Semester workload} &
\multicolumn{2}{p{8cm}|}{Weekly workload}\\
\hline
\multicolumn{1}{|p{4cm}|}{Total} &
\multicolumn{1}{p{4cm}|}{Extension} &
\multicolumn{1}{p{4cm}|}{Theoretical} &
\multicolumn{1}{p{4cm}|}{Practical} \\ 
\multicolumn{1}{|p{4cm}|}{60\,hours} &
\multicolumn{1}{p{4cm}|}{0\;hours} &
\multicolumn{1}{p{4cm}|}{4\;class hours} &
\multicolumn{1}{p{4cm}|}{0\;class hours} \\ 
\hline
\multicolumn{4}{|p{16cm}|}{Course Description:}\\
\multicolumn{4}{|p{16cm}|}{}\\
\multicolumn{4}{|p{\dimexpr 16cm + 6\tabcolsep\relax}|}{Basic concepts of graphs; network flow models; shortest path, maximum flow, and minimum-cost flow algorithms; applications and implementation of specialized algorithms.}\\
\multicolumn{4}{|p{16cm}|}{}\\
\hline
\multicolumn{4}{|p{16cm}|}{Course Content:}\\*
\multicolumn{4}{|p{\dimexpr 16cm + 6\tabcolsep\relax}|}{\begin{enumerate}[start=1]\item Basic graph concepts: notation and definitions; network representation; examples of network models\end{enumerate}}\\
\multicolumn{4}{|p{\dimexpr 16cm + 6\tabcolsep\relax}|}{\begin{enumerate}[start=2]\item Shortest path problems: introduction; applications; solution algorithms and implementations\end{enumerate}}\\
\multicolumn{4}{|p{\dimexpr 16cm + 6\tabcolsep\relax}|}{\begin{enumerate}[start=3]\item Maximum flow problem: introduction; applications; flows and cuts; max-flow min-cut theorem; solution algorithms and implementations\end{enumerate}}\\
\multicolumn{4}{|p{\dimexpr 16cm + 6\tabcolsep\relax}|}{\begin{enumerate}[start=4]\item Minimum-cost flow problems: introduction; applications; solution algorithms and implementations\end{enumerate}}\\
\multicolumn{4}{|p{\dimexpr 16cm + 6\tabcolsep\relax}|}{\begin{enumerate}[start=5]\item Applications: assignment problems; matching problems; routing problems; sequencing problems\end{enumerate}}\\
\multicolumn{4}{|p{\dimexpr 16cm + 6\tabcolsep\relax}|}{}\\
\hline
\multicolumn{4}{|p{16cm}|}{Basic Bibliography:}\\*
\multicolumn{4}{|p{\dimexpr 16cm + 6\tabcolsep\relax}|}{\begin{itemize}\item Marco Goldbarg; Elizabeth Goldbarg. Grafos: conceitos, algoritmos e aplicações. Rio de Janeiro: Elsevier, 2012.\end{itemize}}\\
\multicolumn{4}{|p{\dimexpr 16cm + 6\tabcolsep\relax}|}{\begin{itemize}\item Chvatal, Vasek. Linear programming. New York: W. H. Freeman, 1983.\end{itemize}}\\
\multicolumn{4}{|p{\dimexpr 16cm + 6\tabcolsep\relax}|}{\begin{itemize}\item Frazer, J. Ronald. Applied linear programming. Englewood Cliffs, N.J.: Prentice-Hall, 1968.\end{itemize}}\\
\multicolumn{4}{|p{\dimexpr 16cm + 6\tabcolsep\relax}|}{\begin{itemize}\item Ravindra K. Ahuja, Thomas L. Magnanti, James B. Orlin. Network flows: theory, algorithms and applications. Englewood Cliffs: Prentice-Hall, 1993.\end{itemize}}\\
\multicolumn{4}{|p{\dimexpr 16cm + 6\tabcolsep\relax}|}{\begin{itemize}\item Bazaraa, M. S. Linear programming and network flows. 2. ed. New York: J. Wiley, 1990.\end{itemize}}\\
\multicolumn{4}{|p{16cm}|}{}\\
\hline
\multicolumn{4}{|p{16cm}|}{Complementary Bibliography:}\\*
\multicolumn{4}{|p{\dimexpr 16cm + 6\tabcolsep\relax}|}{\begin{itemize}\item Hu, T. C. Integer programming and network flows. Reading: Addison Wesley, 1970.\end{itemize}}\\
\multicolumn{4}{|p{\dimexpr 16cm + 6\tabcolsep\relax}|}{\begin{itemize}\item Marcos Arenales, Vinicius Armentano, Reinaldo Morabito, Horacio Yanasse. Pesquisa Operacional: para cursos de engenharia. Rio de Janeiro: Elsevier, 2007.\end{itemize}}\\
\multicolumn{4}{|p{\dimexpr 16cm + 6\tabcolsep\relax}|}{\begin{itemize}\item Wolsey, Laurence A. Integer programming. New York: John Wiley \& Sons, 1998.\end{itemize}}\\
\multicolumn{4}{|p{\dimexpr 16cm + 6\tabcolsep\relax}|}{\begin{itemize}\item Marco Cesar Goldbarg, Henrique Pacca L. Luna. Otimização combinatória e programação linear: modelos e algoritmos. 2. ed. rev. e atual. Rio de Janeiro: Campus, 2005.\end{itemize}}\\
\hline
\end{longtable}
\end{center}

\clearpage
\begin{center}
\begin{tabular}{lccr}
 \multirow{3}{*}{\includegraphics[height=2.7cm]{brasao.png}} &
 \multicolumn{2}{c}{\bfseries FEDERAL UNIVERSITY OF OURO PRETO} &
 \ \ \ \ \multirow{3}{*}{\includegraphics[height=2.7cm]{ufop.png}} \\
 & \multicolumn{2}{c}{\bfseries DEAN'S OFFICE FOR UNDERGRADUATE STUDIES} & \\
 & \multicolumn{2}{c}{} & \\
 & \multicolumn{2}{c}{\Large\bfseries COURSE PROGRAM} & \\
\end{tabular}
\end{center}

\begin{center}
\begin{longtable}{|p{4cm}|p{4cm}|p{4cm}|p{4cm}|}
\hline
\multicolumn{3}{|p{12cm}|}{Course Name in English:} &
\multicolumn{1}{p{4cm}|}{Code:} \\ 
\multicolumn{3}{|p{12cm}|}{\textbf{Linear and Integer Optimization}} &
\textbf{BCC544}\\ 
\multicolumn{3}{|p{12cm}|}{Course Name in Portuguese:} & \\ 
\multicolumn{3}{|p{12cm}|}{\textbf{Otimização Linear e Inteira}} & \\ 
\hline
\multicolumn{3}{|p{12cm}|}{Department Name and Acronym:} & Academic Unit: \\ 
\multicolumn{3}{|p{12cm}|}{Department of Computer Science (DECOM)} & {ICEB} \\ 
\hline
\multicolumn{4}{|p{16cm}|}{Delivery Mode:
[X] in person \hspace{0.6cm}
[ ] blended \hspace{0.6cm}
[ ] distance learning}\\
\hline
\multicolumn{2}{|p{8cm}|}{Semester workload} &
\multicolumn{2}{p{8cm}|}{Weekly workload}\\
\hline
\multicolumn{1}{|p{4cm}|}{Total} &
\multicolumn{1}{p{4cm}|}{Extension} &
\multicolumn{1}{p{4cm}|}{Theoretical} &
\multicolumn{1}{p{4cm}|}{Practical} \\ 
\multicolumn{1}{|p{4cm}|}{60\,hours} &
\multicolumn{1}{p{4cm}|}{0\;hours} &
\multicolumn{1}{p{4cm}|}{4\;class hours} &
\multicolumn{1}{p{4cm}|}{0\;class hours} \\ 
\hline
\multicolumn{4}{|p{16cm}|}{Course Description:}\\
\multicolumn{4}{|p{16cm}|}{}\\
\multicolumn{4}{|p{\dimexpr 16cm + 6\tabcolsep\relax}|}{Linear programming modeling. Simplex algorithm. Duality. Sensitivity analysis. Column generation. Dantzig-Wolfe decomposition methods. Integer programming modeling. Implicit enumeration. Cutting planes. Bounds and relaxations.}\\
\multicolumn{4}{|p{16cm}|}{}\\
\hline
\multicolumn{4}{|p{16cm}|}{Course Content:}\\*
\multicolumn{4}{|p{\dimexpr 16cm + 6\tabcolsep\relax}|}{\begin{enumerate}[start=1]\item Introduction to Linear Programming (History)\end{enumerate}}\\
\multicolumn{4}{|p{\dimexpr 16cm + 6\tabcolsep\relax}|}{\begin{enumerate}[start=2]\item Modeling simple problems\end{enumerate}}\\
\multicolumn{4}{|p{\dimexpr 16cm + 6\tabcolsep\relax}|}{\begin{enumerate}[start=3]\item Graphical method\end{enumerate}}\\
\multicolumn{4}{|p{\dimexpr 16cm + 6\tabcolsep\relax}|}{\begin{enumerate}[start=4]\item Problem modeling and example applications\end{enumerate}}\\
\multicolumn{4}{|p{\dimexpr 16cm + 6\tabcolsep\relax}|}{\begin{enumerate}[start=5]\item MathProg modeling language\end{enumerate}}\\
\multicolumn{4}{|p{\dimexpr 16cm + 6\tabcolsep\relax}|}{\begin{enumerate}[start=6]\item Simplex Algorithm\end{enumerate}}\\
\multicolumn{4}{|p{\dimexpr 16cm + 6\tabcolsep\relax}|}{\begin{enumerate}[start=7]\item Duality\end{enumerate}}\\
\multicolumn{4}{|p{\dimexpr 16cm + 6\tabcolsep\relax}|}{\begin{enumerate}[start=8]\item Sensitivity analysis\end{enumerate}}\\
\multicolumn{4}{|p{\dimexpr 16cm + 6\tabcolsep\relax}|}{\begin{enumerate}[start=9]\item Column generation\end{enumerate}}\\
\multicolumn{4}{|p{\dimexpr 16cm + 6\tabcolsep\relax}|}{\begin{enumerate}[start=10]\item Dantzig-Wolfe decomposition\end{enumerate}}\\
\multicolumn{4}{|p{\dimexpr 16cm + 6\tabcolsep\relax}|}{\begin{enumerate}[start=11]\item Integer Programming modeling\end{enumerate}}\\
\multicolumn{4}{|p{\dimexpr 16cm + 6\tabcolsep\relax}|}{\begin{enumerate}[start=12]\item Evaluating Formulation Quality\end{enumerate}}\\
\multicolumn{4}{|p{\dimexpr 16cm + 6\tabcolsep\relax}|}{\begin{enumerate}[start=13]\item Generation of valid inequalities and automatic reformulation\end{enumerate}}\\
\multicolumn{4}{|p{\dimexpr 16cm + 6\tabcolsep\relax}|}{\begin{enumerate}[start=14]\item Techniques for generating strong dual and primal bounds\end{enumerate}}\\
\multicolumn{4}{|p{\dimexpr 16cm + 6\tabcolsep\relax}|}{\begin{enumerate}[start=15]\item Column generation in Integer Programming\end{enumerate}}\\
\multicolumn{4}{|p{\dimexpr 16cm + 6\tabcolsep\relax}|}{}\\
\hline
\multicolumn{4}{|p{16cm}|}{Basic Bibliography:}\\*
\multicolumn{4}{|p{\dimexpr 16cm + 6\tabcolsep\relax}|}{\begin{itemize}\item Maculan, Nelson. Otimização linear. Brasília: Ed. UnB, 2006.\end{itemize}}\\
\multicolumn{4}{|p{\dimexpr 16cm + 6\tabcolsep\relax}|}{\begin{itemize}\item Chvatal, Vasek. Linear programming. New York: W. H. Freeman, 1983.\end{itemize}}\\
\multicolumn{4}{|p{\dimexpr 16cm + 6\tabcolsep\relax}|}{\begin{itemize}\item Wolsey, Laurence A. Integer programming. New York: John Wiley \& Sons, 1998.\end{itemize}}\\
\multicolumn{4}{|p{16cm}|}{}\\
\hline
\multicolumn{4}{|p{16cm}|}{Complementary Bibliography:}\\*
\multicolumn{4}{|p{\dimexpr 16cm + 6\tabcolsep\relax}|}{\begin{itemize}\item Edited by John K. Karlof. Integer programming: theory and practice. Boca Raton, Fla.\end{itemize}}\\
\multicolumn{4}{|p{\dimexpr 16cm + 6\tabcolsep\relax}|}{\begin{itemize}\item Dasgupta, Sanjoy. Algoritmos. São Paulo: McGraw-Hill, 2009.\end{itemize}}\\
\multicolumn{4}{|p{\dimexpr 16cm + 6\tabcolsep\relax}|}{\begin{itemize}\item Michael Jünger... [et al.], editors. 50 years of integer programming, 1958-2008: the early years and state-of-the-art surveys. Heidelberg: Springer, 2010.\end{itemize}}\\
\multicolumn{4}{|p{\dimexpr 16cm + 6\tabcolsep\relax}|}{\begin{itemize}\item Lee, Jon. A first course in combinatorial optimization. Cambridge, UK.\end{itemize}}\\
\multicolumn{4}{|p{\dimexpr 16cm + 6\tabcolsep\relax}|}{\begin{itemize}\item Marco Cesar Goldbarg, Henrique Pacca L. Luna. Otimização combinatória e programação linear: modelos e algoritmos. Rio de Janeiro: Campus, 2000.\end{itemize}}\\
\multicolumn{4}{|p{\dimexpr 16cm + 6\tabcolsep\relax}|}{\begin{itemize}\item Dantzig, George B. Linear programming and extensions. Princeton, N.J.: Princeton University Press, 1963.\end{itemize}}\\
\hline
\end{longtable}
\end{center}

\clearpage
\begin{center}
\begin{tabular}{lccr}
 \multirow{3}{*}{\includegraphics[height=2.7cm]{brasao.png}} &
 \multicolumn{2}{c}{\bfseries FEDERAL UNIVERSITY OF OURO PRETO} &
 \ \ \ \ \multirow{3}{*}{\includegraphics[height=2.7cm]{ufop.png}} \\
 & \multicolumn{2}{c}{\bfseries DEAN'S OFFICE FOR UNDERGRADUATE STUDIES} & \\
 & \multicolumn{2}{c}{} & \\
 & \multicolumn{2}{c}{\Large\bfseries COURSE PROGRAM} & \\
\end{tabular}
\end{center}

\begin{center}
\begin{longtable}{|p{4cm}|p{4cm}|p{4cm}|p{4cm}|}
\hline
\multicolumn{3}{|p{12cm}|}{Course Name in English:} &
\multicolumn{1}{p{4cm}|}{Code:} \\ 
\multicolumn{3}{|p{12cm}|}{\textbf{Multi-Objective Optimization Techniques}} &
\textbf{BCC545}\\ 
\multicolumn{3}{|p{12cm}|}{Course Name in Portuguese:} & \\ 
\multicolumn{3}{|p{12cm}|}{\textbf{Técnicas de Otimização Multiobjetivo}} & \\ 
\hline
\multicolumn{3}{|p{12cm}|}{Department Name and Acronym:} & Academic Unit: \\ 
\multicolumn{3}{|p{12cm}|}{Department of Computer Science (DECOM)} & {ICEB} \\ 
\hline
\multicolumn{4}{|p{16cm}|}{Delivery Mode:
[X] in person \hspace{0.6cm}
[ ] blended \hspace{0.6cm}
[ ] distance learning}\\
\hline
\multicolumn{2}{|p{8cm}|}{Semester workload} &
\multicolumn{2}{p{8cm}|}{Weekly workload}\\
\hline
\multicolumn{1}{|p{4cm}|}{Total} &
\multicolumn{1}{p{4cm}|}{Extension} &
\multicolumn{1}{p{4cm}|}{Theoretical} &
\multicolumn{1}{p{4cm}|}{Practical} \\ 
\multicolumn{1}{|p{4cm}|}{60\,hours} &
\multicolumn{1}{p{4cm}|}{0\;hours} &
\multicolumn{1}{p{4cm}|}{4\;class hours} &
\multicolumn{1}{p{4cm}|}{0\;class hours} \\ 
\hline
\multicolumn{4}{|p{16cm}|}{Course Description:}\\
\multicolumn{4}{|p{16cm}|}{}\\
\multicolumn{4}{|p{\dimexpr 16cm + 6\tabcolsep\relax}|}{Mathematical foundations; basic concepts of nonlinear optimization; optimality conditions; unconstrained and constrained minimization methods; objective functionals and objective vectors; ordered sets; Pareto solutions: analytical characterization; generation of Pareto solutions; preference indication; algorithms for multi-objective optimization; multi-objective optimization with evolutionary algorithms; decision making.}\\
\multicolumn{4}{|p{16cm}|}{}\\
\hline
\multicolumn{4}{|p{16cm}|}{Course Content:}\\*
\multicolumn{4}{|p{\dimexpr 16cm + 6\tabcolsep\relax}|}{\begin{enumerate}[start=1]\item Course plan presentation and introduction: The vector optimization problem; Notation\end{enumerate}}\\
\multicolumn{4}{|p{\dimexpr 16cm + 6\tabcolsep\relax}|}{\begin{enumerate}[start=2]\item Ordering of solutions: Pareto-optimal set; locally Pareto-optimal set; utopian solution\end{enumerate}}\\
\multicolumn{4}{|p{\dimexpr 16cm + 6\tabcolsep\relax}|}{\begin{enumerate}[start=3]\item The problem of determining efficient solutions: Kuhn-Tucker conditions for efficiency\end{enumerate}}\\
\multicolumn{4}{|p{\dimexpr 16cm + 6\tabcolsep\relax}|}{\begin{enumerate}[start=4]\item Generation of efficient solutions: Weighted-sum approach; Epsilon-constraint approach; goal-programming approach; efficiency test\end{enumerate}}\\
\multicolumn{4}{|p{\dimexpr 16cm + 6\tabcolsep\relax}|}{\begin{enumerate}[start=5]\item Structure of the Pareto-optimal set\end{enumerate}}\\
\multicolumn{4}{|p{\dimexpr 16cm + 6\tabcolsep\relax}|}{\begin{enumerate}[start=6]\item Multi-objective Analysis: Consistency; ordering and dominance; extension\end{enumerate}}\\
\multicolumn{4}{|p{\dimexpr 16cm + 6\tabcolsep\relax}|}{\begin{enumerate}[start=7]\item Multi-objective Decision and Synthesis\end{enumerate}}\\
\multicolumn{4}{|p{\dimexpr 16cm + 6\tabcolsep\relax}|}{\begin{enumerate}[start=8]\item Multi-objective Genetic Algorithms: Construction of Multi-objective Genetic Algorithms\end{enumerate}}\\
\multicolumn{4}{|p{\dimexpr 16cm + 6\tabcolsep\relax}|}{}\\
\hline
\multicolumn{4}{|p{16cm}|}{Basic Bibliography:}\\*
\multicolumn{4}{|p{\dimexpr 16cm + 6\tabcolsep\relax}|}{\begin{itemize}\item DEB, Kalyanmoy. Multi-Objective Optimization using Evolutionary Algorithms. John Wiley \& Sons, 2001.\end{itemize}}\\
\multicolumn{4}{|p{\dimexpr 16cm + 6\tabcolsep\relax}|}{\begin{itemize}\item COELLO COELLO, Carlos. Evolutionary Algorithms for Solving Multi-Objective Problems. 2. ed. Springer, 2007.\end{itemize}}\\
\multicolumn{4}{|p{\dimexpr 16cm + 6\tabcolsep\relax}|}{\begin{itemize}\item Vira Chankong and Yacov Y. Haimes. Multiobjective decision making: theory and methodology. New York: North-Holland, 1983.\end{itemize}}\\
\multicolumn{4}{|p{\dimexpr 16cm + 6\tabcolsep\relax}|}{\begin{itemize}\item Yann Collette, Patrick Siarry. Multiobjective optimization: principles and case studies. Berlin.\end{itemize}}\\
\multicolumn{4}{|p{\dimexpr 16cm + 6\tabcolsep\relax}|}{\begin{itemize}\item Mokhtar S. Bazaraa, Hanif D. Sherali, C. M. Shetty. Nonlinear programming: theory and algorithms. 3rd ed. Hoboken, N.J.: Wiley-Interscience, 2006.\end{itemize}}\\
\multicolumn{4}{|p{16cm}|}{}\\
\hline
\multicolumn{4}{|p{16cm}|}{Complementary Bibliography:}\\*
\multicolumn{4}{|p{\dimexpr 16cm + 6\tabcolsep\relax}|}{\begin{itemize}\item SAWARAGI, Yoshikazu. Theory of Multiobjective Optimization. Elsevier, 1985.\end{itemize}}\\
\multicolumn{4}{|p{\dimexpr 16cm + 6\tabcolsep\relax}|}{\begin{itemize}\item ABRAHAM, Ajith. Evolutionary Multiobjective Optimization: Theoretical Advances and Applications. Springer, 2005.\end{itemize}}\\
\multicolumn{4}{|p{\dimexpr 16cm + 6\tabcolsep\relax}|}{\begin{itemize}\item EICHFELDER, Gabriele. Adaptive Scalarization Methods in Multiobjective Optimization. Springer, 2018.\end{itemize}}\\
\multicolumn{4}{|p{\dimexpr 16cm + 6\tabcolsep\relax}|}{\begin{itemize}\item MIETTINEN, Kaisa. Nonlinear Multiobjective Optimization. Springer, 1999.\end{itemize}}\\
\multicolumn{4}{|p{\dimexpr 16cm + 6\tabcolsep\relax}|}{\begin{itemize}\item TANG, Kit. Multiobjective Optimization Methodology: A Jumping Gene Approach. 1. ed. CRC Press, 2018.\end{itemize}}\\
\hline
\end{longtable}
\end{center}

\clearpage
\begin{center}
\begin{tabular}{lccr}
 \multirow{3}{*}{\includegraphics[height=2.7cm]{brasao.png}} &
 \multicolumn{2}{c}{\bfseries FEDERAL UNIVERSITY OF OURO PRETO} &
 \ \ \ \ \multirow{3}{*}{\includegraphics[height=2.7cm]{ufop.png}} \\
 & \multicolumn{2}{c}{\bfseries DEAN'S OFFICE FOR UNDERGRADUATE STUDIES} & \\
 & \multicolumn{2}{c}{} & \\
 & \multicolumn{2}{c}{\Large\bfseries COURSE PROGRAM} & \\
\end{tabular}
\end{center}

\begin{center}
\begin{longtable}{|p{4cm}|p{4cm}|p{4cm}|p{4cm}|}
\hline
\multicolumn{3}{|p{12cm}|}{Course Name in English:} &
\multicolumn{1}{p{4cm}|}{Code:} \\ 
\multicolumn{3}{|p{12cm}|}{\textbf{Metaheuristics Techniques in Combinatorial Optimization}} &
\textbf{BCC546}\\ 
\multicolumn{3}{|p{12cm}|}{Course Name in Portuguese:} & \\ 
\multicolumn{3}{|p{12cm}|}{\textbf{Técnicas Metaheurísticas para Otimização Combinatória}} & \\ 
\hline
\multicolumn{3}{|p{12cm}|}{Department Name and Acronym:} & Academic Unit: \\ 
\multicolumn{3}{|p{12cm}|}{Department of Computer Science (DECOM)} & {ICEB} \\ 
\hline
\multicolumn{4}{|p{16cm}|}{Delivery Mode:
[X] in person \hspace{0.6cm}
[ ] blended \hspace{0.6cm}
[ ] distance learning}\\
\hline
\multicolumn{2}{|p{8cm}|}{Semester workload} &
\multicolumn{2}{p{8cm}|}{Weekly workload}\\
\hline
\multicolumn{1}{|p{4cm}|}{Total} &
\multicolumn{1}{p{4cm}|}{Extension} &
\multicolumn{1}{p{4cm}|}{Theoretical} &
\multicolumn{1}{p{4cm}|}{Practical} \\ 
\multicolumn{1}{|p{4cm}|}{60\,hours} &
\multicolumn{1}{p{4cm}|}{0\;hours} &
\multicolumn{1}{p{4cm}|}{4\;class hours} &
\multicolumn{1}{p{4cm}|}{0\;class hours} \\ 
\hline
\multicolumn{4}{|p{16cm}|}{Course Description:}\\
\multicolumn{4}{|p{16cm}|}{}\\
\multicolumn{4}{|p{\dimexpr 16cm + 6\tabcolsep\relax}|}{Combinatorial optimization problems. Classical heuristics. Local search. Multi-start methods. Trajectory-based metaheuristics. Population-based metaheuristics. Nature-inspired methods. Applications.}\\
\multicolumn{4}{|p{16cm}|}{}\\
\hline
\multicolumn{4}{|p{16cm}|}{Course Content:}\\*
\multicolumn{4}{|p{\dimexpr 16cm + 6\tabcolsep\relax}|}{\begin{enumerate}[start=1]\item Introduction: combinatorial optimization problems\end{enumerate}}\\
\multicolumn{4}{|p{\dimexpr 16cm + 6\tabcolsep\relax}|}{\begin{enumerate}[start=2]\item Greedy algorithms\end{enumerate}}\\
\multicolumn{4}{|p{\dimexpr 16cm + 6\tabcolsep\relax}|}{\begin{enumerate}[start=3]\item Local search\end{enumerate}}\\
\multicolumn{4}{|p{\dimexpr 16cm + 6\tabcolsep\relax}|}{\begin{enumerate}[start=4]\item Large neighborhood search\end{enumerate}}\\
\multicolumn{4}{|p{\dimexpr 16cm + 6\tabcolsep\relax}|}{\begin{enumerate}[start=5]\item Path relinking\end{enumerate}}\\
\multicolumn{4}{|p{\dimexpr 16cm + 6\tabcolsep\relax}|}{\begin{enumerate}[start=6]\item Simulated Annealing\end{enumerate}}\\
\multicolumn{4}{|p{\dimexpr 16cm + 6\tabcolsep\relax}|}{\begin{enumerate}[start=7]\item Tabu Search\end{enumerate}}\\
\multicolumn{4}{|p{\dimexpr 16cm + 6\tabcolsep\relax}|}{\begin{enumerate}[start=8]\item Iterated local search\end{enumerate}}\\
\multicolumn{4}{|p{\dimexpr 16cm + 6\tabcolsep\relax}|}{\begin{enumerate}[start=9]\item Variable neighborhood search\end{enumerate}}\\
\multicolumn{4}{|p{\dimexpr 16cm + 6\tabcolsep\relax}|}{\begin{enumerate}[start=10]\item Multi-start procedures\end{enumerate}}\\
\multicolumn{4}{|p{\dimexpr 16cm + 6\tabcolsep\relax}|}{\begin{enumerate}[start=11]\item Nature-inspired algorithms\end{enumerate}}\\
\multicolumn{4}{|p{\dimexpr 16cm + 6\tabcolsep\relax}|}{\begin{enumerate}[start=12]\item Evolutionary algorithms\end{enumerate}}\\
\multicolumn{4}{|p{\dimexpr 16cm + 6\tabcolsep\relax}|}{\begin{enumerate}[start=13]\item Scatter Search\end{enumerate}}\\
\multicolumn{4}{|p{\dimexpr 16cm + 6\tabcolsep\relax}|}{\begin{enumerate}[start=14]\item Applications in routing\end{enumerate}}\\
\multicolumn{4}{|p{\dimexpr 16cm + 6\tabcolsep\relax}|}{\begin{enumerate}[start=15]\item Applications in sequencing\end{enumerate}}\\
\multicolumn{4}{|p{\dimexpr 16cm + 6\tabcolsep\relax}|}{\begin{enumerate}[start=16]\item Applications in timetabling\end{enumerate}}\\
\multicolumn{4}{|p{\dimexpr 16cm + 6\tabcolsep\relax}|}{\begin{enumerate}[start=17]\item Applications in network design\end{enumerate}}\\
\multicolumn{4}{|p{\dimexpr 16cm + 6\tabcolsep\relax}|}{}\\
\hline
\multicolumn{4}{|p{16cm}|}{Basic Bibliography:}\\*
\multicolumn{4}{|p{\dimexpr 16cm + 6\tabcolsep\relax}|}{\begin{itemize}\item Handbook of approximation algorithms and metaheuristics. New York: Chapman \& Hall/CRC, 2007.\end{itemize}}\\
\multicolumn{4}{|p{\dimexpr 16cm + 6\tabcolsep\relax}|}{\begin{itemize}\item Voss, Stefan. Introduction to computational optimization models for production planning in a supply chain. 2nd ed. Berlin.\end{itemize}}\\
\multicolumn{4}{|p{\dimexpr 16cm + 6\tabcolsep\relax}|}{\begin{itemize}\item El-Ghazali Talbi. Metaheuristics: from design to implementation. Hoboken, N.J.: John Wiley \& Sons, 2009.\end{itemize}}\\
\multicolumn{4}{|p{\dimexpr 16cm + 6\tabcolsep\relax}|}{\begin{itemize}\item Zbigniew Michalewicz, David B. Fogel. How to solve it: modern heuristics. 2nd rev. and extended ed. Berlin.\end{itemize}}\\
\multicolumn{4}{|p{\dimexpr 16cm + 6\tabcolsep\relax}|}{\begin{itemize}\item Mauricio G.C. Resende, Jorge Pinho de Sousa, with assistance from Ana Viana (eds.). Metaheuristics: computer decision-making. Boston: Kluwer Academic Publishers, 2004.\end{itemize}}\\
\multicolumn{4}{|p{16cm}|}{}\\
\hline
\multicolumn{4}{|p{16cm}|}{Complementary Bibliography:}\\*
\multicolumn{4}{|p{\dimexpr 16cm + 6\tabcolsep\relax}|}{\begin{itemize}\item Local search in combinatorial optimization. New Jersey: Princeton University Press, 2003.\end{itemize}}\\
\multicolumn{4}{|p{\dimexpr 16cm + 6\tabcolsep\relax}|}{\begin{itemize}\item J. Dréo... [et al.]. Metaheuristics for hard optimization: methods and case studies. Berlin: Springer, 2006.\end{itemize}}\\
\multicolumn{4}{|p{\dimexpr 16cm + 6\tabcolsep\relax}|}{\begin{itemize}\item Kenneth V. Price, Rainer M. Storn, Jouni A. Lampinen. Differential evolution: a practical approach to global optimization. Berlin.\end{itemize}}\\
\multicolumn{4}{|p{\dimexpr 16cm + 6\tabcolsep\relax}|}{\begin{itemize}\item Dorigo, Marco. Ant colony optimization. Cambridge: The MIT Press, 2004.\end{itemize}}\\
\multicolumn{4}{|p{\dimexpr 16cm + 6\tabcolsep\relax}|}{\begin{itemize}\item Edited by Vittorio Maniezzo, Thomas Stutzle, Stefan Voss. Matheuristics: hybridizing metaheuristics and mathematical programming. New York: Springer, 2009.\end{itemize}}\\
\hline
\end{longtable}
\end{center}

\clearpage
\begin{center}
\begin{tabular}{lccr}
 \multirow{3}{*}{\includegraphics[height=2.7cm]{brasao.png}} &
 \multicolumn{2}{c}{\bfseries FEDERAL UNIVERSITY OF OURO PRETO} &
 \ \ \ \ \multirow{3}{*}{\includegraphics[height=2.7cm]{ufop.png}} \\
 & \multicolumn{2}{c}{\bfseries DEAN'S OFFICE FOR UNDERGRADUATE STUDIES} & \\
 & \multicolumn{2}{c}{} & \\
 & \multicolumn{2}{c}{\Large\bfseries COURSE PROGRAM} & \\
\end{tabular}
\end{center}

\begin{center}
\begin{longtable}{|p{4cm}|p{4cm}|p{4cm}|p{4cm}|}
\hline
\multicolumn{3}{|p{12cm}|}{Course Name in English:} &
\multicolumn{1}{p{4cm}|}{Code:} \\ 
\multicolumn{3}{|p{12cm}|}{\textbf{Software Project Management}} &
\textbf{BCC547}\\ 
\multicolumn{3}{|p{12cm}|}{Course Name in Portuguese:} & \\ 
\multicolumn{3}{|p{12cm}|}{\textbf{Gerência de Projetos de Software}} & \\ 
\hline
\multicolumn{3}{|p{12cm}|}{Department Name and Acronym:} & Academic Unit: \\ 
\multicolumn{3}{|p{12cm}|}{Department of Computer Science (DECOM)} & {ICEB} \\ 
\hline
\multicolumn{4}{|p{16cm}|}{Delivery Mode:
[X] in person \hspace{0.6cm}
[ ] blended \hspace{0.6cm}
[ ] distance learning}\\
\hline
\multicolumn{2}{|p{8cm}|}{Semester workload} &
\multicolumn{2}{p{8cm}|}{Weekly workload}\\
\hline
\multicolumn{1}{|p{4cm}|}{Total} &
\multicolumn{1}{p{4cm}|}{Extension} &
\multicolumn{1}{p{4cm}|}{Theoretical} &
\multicolumn{1}{p{4cm}|}{Practical} \\ 
\multicolumn{1}{|p{4cm}|}{60\,hours} &
\multicolumn{1}{p{4cm}|}{0\;hours} &
\multicolumn{1}{p{4cm}|}{4\;class hours} &
\multicolumn{1}{p{4cm}|}{0\;class hours} \\ 
\hline
\multicolumn{4}{|p{16cm}|}{Course Description:}\\
\multicolumn{4}{|p{16cm}|}{}\\
\multicolumn{4}{|p{\dimexpr 16cm + 6\tabcolsep\relax}|}{Software project management, software process, software life cycle, project management according to PMI-PMBOK, agile project management (SCRUM), management tools.}\\
\multicolumn{4}{|p{16cm}|}{}\\
\hline
\multicolumn{4}{|p{16cm}|}{Course Content:}\\*
\multicolumn{4}{|p{\dimexpr 16cm + 6\tabcolsep\relax}|}{\begin{enumerate}[start=1]\item Software Project Management: Project Management; Software Project Management; Software Engineering and Project Management; project management process; traditional and agile management\end{enumerate}}\\
\multicolumn{4}{|p{\dimexpr 16cm + 6\tabcolsep\relax}|}{\begin{enumerate}[start=2]\item Software Process: Process; software development processes; life cycle\end{enumerate}}\\
\multicolumn{4}{|p{\dimexpr 16cm + 6\tabcolsep\relax}|}{\begin{enumerate}[start=3]\item Software: What is software; software as a product\end{enumerate}}\\
\multicolumn{4}{|p{\dimexpr 16cm + 6\tabcolsep\relax}|}{\begin{enumerate}[start=4]\item PMI: What is PMI; process groups; knowledge areas; projects; integration management; scope management; time management; cost management; quality management; human resources management; communications management; risk management; procurement management; stakeholder management\end{enumerate}}\\
\multicolumn{4}{|p{\dimexpr 16cm + 6\tabcolsep\relax}|}{\begin{enumerate}[start=5]\item SCRUM: Definition; SCRUM model; roles; artifacts; meetings\end{enumerate}}\\
\multicolumn{4}{|p{\dimexpr 16cm + 6\tabcolsep\relax}|}{\begin{enumerate}[start=6]\item Tools: Tools supporting the various activities and practices of project management in traditional and agile approaches\end{enumerate}}\\
\multicolumn{4}{|p{\dimexpr 16cm + 6\tabcolsep\relax}|}{}\\
\hline
\multicolumn{4}{|p{16cm}|}{Basic Bibliography:}\\*
\multicolumn{4}{|p{\dimexpr 16cm + 6\tabcolsep\relax}|}{\begin{itemize}\item Martins, José Carlos Cordeiro. Gerenciando projetos de desenvolvimento de software com PMI, RUP e UML. 4. ed. Rio de Janeiro: Brasport, 2007.\end{itemize}}\\
\multicolumn{4}{|p{\dimexpr 16cm + 6\tabcolsep\relax}|}{\begin{itemize}\item Schwaber, Ken. Agile project management with Scrum. Redmond, Wash.: Microsoft Press, 2004.\end{itemize}}\\
\multicolumn{4}{|p{\dimexpr 16cm + 6\tabcolsep\relax}|}{\begin{itemize}\item Kim Heldman; tradução Luciana do Amaral Teixeira; revisão técnica Martius Vicente Rodriguez Y Rodriguez. Gerência de projetos = PMP Project management professional: guia para o exame oficial do PMI. 5. ed. rev. e atual. Rio de Janeiro: Elsevier, 2009.\end{itemize}}\\
\multicolumn{4}{|p{16cm}|}{}\\
\hline
\multicolumn{4}{|p{16cm}|}{Complementary Bibliography:}\\*
\multicolumn{4}{|p{\dimexpr 16cm + 6\tabcolsep\relax}|}{\begin{itemize}\item Sommerville, Ian. Engenharia de software. 9. ed. São Paulo: Addison Wesley, 2011.\end{itemize}}\\
\multicolumn{4}{|p{\dimexpr 16cm + 6\tabcolsep\relax}|}{\begin{itemize}\item Roger S. Pressman; tradução Ariovaldo Griesi, Mario Moro Fecchio; revisão técnica Reginaldo Arakaki, Julio Arakaki, Renato Manzan de Andrade. Engenharia de software: uma abordagem profissional. 7. ed. Porto Alegre, RS: AMGH, 2011.\end{itemize}}\\
\multicolumn{4}{|p{\dimexpr 16cm + 6\tabcolsep\relax}|}{\begin{itemize}\item Daniel Capaldo Amaral... [et al.]. Gerenciamento Ágil de projetos: aplicação em produtos inovadores. São Paulo: Saraiva, 2011.\end{itemize}}\\
\multicolumn{4}{|p{\dimexpr 16cm + 6\tabcolsep\relax}|}{\begin{itemize}\item Gattoni, Roberto Luís Capuruço. Gestão do conhecimento aplicada à prática da gerência de projetos. Belo Horizonte, MG: FUMEC-FACE.\end{itemize}}\\
\multicolumn{4}{|p{\dimexpr 16cm + 6\tabcolsep\relax}|}{\begin{itemize}\item Joseph Phillips; consultoria editorial Lorenzo Ridolfi, Sérgio Colcher; tradução Ana Beatriz Tavares, Daniela Lacerda; revisão técnica Andrea Hemerly. Gerência de projetos de tecnologia da informação: no caminho certo, do início ao fim. Rio de Janeiro: Campus.\end{itemize}}\\
\hline
\end{longtable}
\end{center}

\clearpage
\begin{center}
\begin{tabular}{lccr}
 \multirow{3}{*}{\includegraphics[height=2.7cm]{brasao.png}} &
 \multicolumn{2}{c}{\bfseries FEDERAL UNIVERSITY OF OURO PRETO} &
 \ \ \ \ \multirow{3}{*}{\includegraphics[height=2.7cm]{ufop.png}} \\
 & \multicolumn{2}{c}{\bfseries DEAN'S OFFICE FOR UNDERGRADUATE STUDIES} & \\
 & \multicolumn{2}{c}{} & \\
 & \multicolumn{2}{c}{\Large\bfseries COURSE PROGRAM} & \\
\end{tabular}
\end{center}

\begin{center}
\begin{longtable}{|p{4cm}|p{4cm}|p{4cm}|p{4cm}|}
\hline
\multicolumn{3}{|p{12cm}|}{Course Name in English:} &
\multicolumn{1}{p{4cm}|}{Code:} \\ 
\multicolumn{3}{|p{12cm}|}{\textbf{Software Quality}} &
\textbf{BCC553}\\ 
\multicolumn{3}{|p{12cm}|}{Course Name in Portuguese:} & \\ 
\multicolumn{3}{|p{12cm}|}{\textbf{Qualidade de Software}} & \\ 
\hline
\multicolumn{3}{|p{12cm}|}{Department Name and Acronym:} & Academic Unit: \\ 
\multicolumn{3}{|p{12cm}|}{Department of Computer Science (DECOM)} & {ICEB} \\ 
\hline
\multicolumn{4}{|p{16cm}|}{Delivery Mode:
[X] in person \hspace{0.6cm}
[ ] blended \hspace{0.6cm}
[ ] distance learning}\\
\hline
\multicolumn{2}{|p{8cm}|}{Semester workload} &
\multicolumn{2}{p{8cm}|}{Weekly workload}\\
\hline
\multicolumn{1}{|p{4cm}|}{Total} &
\multicolumn{1}{p{4cm}|}{Extension} &
\multicolumn{1}{p{4cm}|}{Theoretical} &
\multicolumn{1}{p{4cm}|}{Practical} \\ 
\multicolumn{1}{|p{4cm}|}{60\,hours} &
\multicolumn{1}{p{4cm}|}{0\;hours} &
\multicolumn{1}{p{4cm}|}{4\;class hours} &
\multicolumn{1}{p{4cm}|}{0\;class hours} \\ 
\hline
\multicolumn{4}{|p{16cm}|}{Course Description:}\\
\multicolumn{4}{|p{16cm}|}{}\\
\multicolumn{4}{|p{\dimexpr 16cm + 6\tabcolsep\relax}|}{What is quality. Software processes. Process quality. Product quality. MPS-BR (Software and service development). CMMI. Metrics. Testing. Verification. Validation. Auditing.}\\
\multicolumn{4}{|p{16cm}|}{}\\
\hline
\multicolumn{4}{|p{16cm}|}{Course Content:}\\*
\multicolumn{4}{|p{\dimexpr 16cm + 6\tabcolsep\relax}|}{\begin{enumerate}[start=1]\item What is Quality: Quality; Software Quality; Process Quality\end{enumerate}}\\
\multicolumn{4}{|p{\dimexpr 16cm + 6\tabcolsep\relax}|}{\begin{enumerate}[start=2]\item Process Quality: What is a process and a software process; process maturity; norms and standards; process improvement\end{enumerate}}\\
\multicolumn{4}{|p{\dimexpr 16cm + 6\tabcolsep\relax}|}{\begin{enumerate}[start=3]\item Product Quality: What is a product and a software product; norms and standards\end{enumerate}}\\
\multicolumn{4}{|p{\dimexpr 16cm + 6\tabcolsep\relax}|}{\begin{enumerate}[start=4]\item MPS-BR: Software Development; Service\end{enumerate}}\\
\multicolumn{4}{|p{\dimexpr 16cm + 6\tabcolsep\relax}|}{\begin{enumerate}[start=5]\item CMMI\end{enumerate}}\\
\multicolumn{4}{|p{\dimexpr 16cm + 6\tabcolsep\relax}|}{\begin{enumerate}[start=6]\item Quality monitoring and control: Metrics; Testing; Verification; Validation; Auditing\end{enumerate}}\\
\multicolumn{4}{|p{\dimexpr 16cm + 6\tabcolsep\relax}|}{}\\
\hline
\multicolumn{4}{|p{16cm}|}{Basic Bibliography:}\\*
\multicolumn{4}{|p{\dimexpr 16cm + 6\tabcolsep\relax}|}{\begin{itemize}\item Mary Beth Chrissis; Mike Konrad; Sandy Shrum. CMMI: guidelines for process integration and product improvement. 2. ed. Boston: Addison Wesley, 2007.\end{itemize}}\\
\multicolumn{4}{|p{\dimexpr 16cm + 6\tabcolsep\relax}|}{\begin{itemize}\item Sommerville, Ian. Engenharia de software. 9. ed. São Paulo: Addison Wesley, 2011.\end{itemize}}\\
\multicolumn{4}{|p{16cm}|}{}\\
\hline
\multicolumn{4}{|p{16cm}|}{Complementary Bibliography:}\\*
\multicolumn{4}{|p{\dimexpr 16cm + 6\tabcolsep\relax}|}{\begin{itemize}\item Bartié, Alexandre. Garantia da qualidade de software. Rio de Janeiro: Campus, 2002.\end{itemize}}\\
\multicolumn{4}{|p{\dimexpr 16cm + 6\tabcolsep\relax}|}{\begin{itemize}\item Edson Junio Rodrigues Lobo. Curso de engenharia de software: métodos e processos para garantir a qualidade no desenvolvimento de softwares. São Paulo: Digeratis Books, 2007.\end{itemize}}\\
\multicolumn{4}{|p{\dimexpr 16cm + 6\tabcolsep\relax}|}{\begin{itemize}\item Roger S. Pressman; tradução Ariovaldo Griesi, Mario Moro Fecchio; revisão técnica Reginaldo Arakaki, Julio Arakaki, Renato Manzan de Andrade. Engenharia de software: uma abordagem profissional. 7. ed. Porto Alegre, RS: AMGH, 2011.\end{itemize}}\\
\multicolumn{4}{|p{\dimexpr 16cm + 6\tabcolsep\relax}|}{\begin{itemize}\item Handbook of software quality assurance. 4th ed. Boston: Artech House, 2008.\end{itemize}}\\
\multicolumn{4}{|p{\dimexpr 16cm + 6\tabcolsep\relax}|}{\begin{itemize}\item Dennis M. Ahern, Aaron Clouse, Richard Turner. CMMI distilled: a practical introduction to integrated process improvement. 3rd ed. Upper Saddle River, NJ: Addison Wesley, 2008.\end{itemize}}\\
\hline
\end{longtable}
\end{center}

\clearpage
\begin{center}
\begin{tabular}{lccr}
 \multirow{3}{*}{\includegraphics[height=2.7cm]{brasao.png}} &
 \multicolumn{2}{c}{\bfseries FEDERAL UNIVERSITY OF OURO PRETO} &
 \ \ \ \ \multirow{3}{*}{\includegraphics[height=2.7cm]{ufop.png}} \\
 & \multicolumn{2}{c}{\bfseries DEAN'S OFFICE FOR UNDERGRADUATE STUDIES} & \\
 & \multicolumn{2}{c}{} & \\
 & \multicolumn{2}{c}{\Large\bfseries COURSE PROGRAM} & \\
\end{tabular}
\end{center}

\begin{center}
\begin{longtable}{|p{4cm}|p{4cm}|p{4cm}|p{4cm}|}
\hline
\multicolumn{3}{|p{12cm}|}{Course Name in English:} &
\multicolumn{1}{p{4cm}|}{Code:} \\ 
\multicolumn{3}{|p{12cm}|}{\textbf{Interaction Design}} &
\textbf{BCC555}\\ 
\multicolumn{3}{|p{12cm}|}{Course Name in Portuguese:} & \\ 
\multicolumn{3}{|p{12cm}|}{\textbf{Design de Interação}} & \\ 
\hline
\multicolumn{3}{|p{12cm}|}{Department Name and Acronym:} & Academic Unit: \\ 
\multicolumn{3}{|p{12cm}|}{Department of Computer Science (DECOM)} & {ICEB} \\ 
\hline
\multicolumn{4}{|p{16cm}|}{Delivery Mode:
[X] in person \hspace{0.6cm}
[ ] blended \hspace{0.6cm}
[ ] distance learning}\\
\hline
\multicolumn{2}{|p{8cm}|}{Semester workload} &
\multicolumn{2}{p{8cm}|}{Weekly workload}\\
\hline
\multicolumn{1}{|p{4cm}|}{Total} &
\multicolumn{1}{p{4cm}|}{Extension} &
\multicolumn{1}{p{4cm}|}{Theoretical} &
\multicolumn{1}{p{4cm}|}{Practical} \\ 
\multicolumn{1}{|p{4cm}|}{60\,hours} &
\multicolumn{1}{p{4cm}|}{0\;hours} &
\multicolumn{1}{p{4cm}|}{4\;class hours} &
\multicolumn{1}{p{4cm}|}{0\;class hours} \\ 
\hline
\multicolumn{4}{|p{16cm}|}{Course Description:}\\
\multicolumn{4}{|p{16cm}|}{}\\
\multicolumn{4}{|p{\dimexpr 16cm + 6\tabcolsep\relax}|}{The Interaction Design process; requirements gathering; task modeling and interaction modeling; prototyping and construction; notions of Visual Design.}\\
\multicolumn{4}{|p{16cm}|}{}\\
\hline
\multicolumn{4}{|p{16cm}|}{Course Content:}\\*
\multicolumn{4}{|p{\dimexpr 16cm + 6\tabcolsep\relax}|}{\begin{enumerate}[start=1]\item The Interaction Design Process: Examples of good and bad designs; What is Interaction Design; Usability goals and user experience goals; Conceptual Models; Interface Metaphors; Interaction Paradigms; Life cycle models for Interaction Design; Life cycle models in Software Engineering; Integration of Software Engineering and HCI Models\end{enumerate}}\\
\multicolumn{4}{|p{\dimexpr 16cm + 6\tabcolsep\relax}|}{\begin{enumerate}[start=2]\item Requirements Gathering: Functional and Non-Functional Requirements; Questionnaires; Interviews; Focus Groups; Card Sorting; Brainstorming; Contextual Inquiry\end{enumerate}}\\
\multicolumn{4}{|p{\dimexpr 16cm + 6\tabcolsep\relax}|}{\begin{enumerate}[start=3]\item Task and Interaction Modeling: Abstract representations of interface/interaction; user profiles; personas and goals; storyboards; analysis scenarios; interaction scenarios; use cases; hierarchical task analysis; interaction diagrams; MoLIC\end{enumerate}}\\
\multicolumn{4}{|p{\dimexpr 16cm + 6\tabcolsep\relax}|}{\begin{enumerate}[start=4]\item Prototyping and Construction: Paper prototypes; low-fidelity prototypes; high-fidelity prototypes; prototyping tools; Conceptual Design: from requirements to the first design; Physical Design\end{enumerate}}\\
\multicolumn{4}{|p{\dimexpr 16cm + 6\tabcolsep\relax}|}{\begin{enumerate}[start=5]\item Notions of Visual Design: Icon design; layout design; color design; typography design\end{enumerate}}\\
\multicolumn{4}{|p{\dimexpr 16cm + 6\tabcolsep\relax}|}{}\\
\hline
\multicolumn{4}{|p{16cm}|}{Basic Bibliography:}\\*
\multicolumn{4}{|p{\dimexpr 16cm + 6\tabcolsep\relax}|}{\begin{itemize}\item Jennifer Preece, Yvonne Rogers e Helen Sharp; tradução de Viviane Possamai. Design de interação: além da interação homem-computador. Porto Alegre, RS: Bookman, 2005.\end{itemize}}\\
\multicolumn{4}{|p{\dimexpr 16cm + 6\tabcolsep\relax}|}{\begin{itemize}\item Robin Williams; tradução: Laura Karin Gillon. Design para quem não é designer: noções básicas de planejamento visual. 3. ed. São Paulo: Callis, 2009.\end{itemize}}\\
\multicolumn{4}{|p{\dimexpr 16cm + 6\tabcolsep\relax}|}{\begin{itemize}\item Ben Shneiderman \& Catherine Plaisant. Designing the user interface: strategies for effective human-computer interaction. 5th ed. Boston: Addison Wesley, 2010.\end{itemize}}\\
\multicolumn{4}{|p{16cm}|}{}\\
\hline
\multicolumn{4}{|p{16cm}|}{Complementary Bibliography:}\\*
\multicolumn{4}{|p{\dimexpr 16cm + 6\tabcolsep\relax}|}{\begin{itemize}\item Barbosa, Simone Diniz Junqueira. Interação humano-computador. Rio de Janeiro: Campus, 2010.\end{itemize}}\\
\multicolumn{4}{|p{\dimexpr 16cm + 6\tabcolsep\relax}|}{\begin{itemize}\item Benyon, David. Interação humano-computador. 2. ed. São Paulo: Pearson, 2011.\end{itemize}}\\
\multicolumn{4}{|p{\dimexpr 16cm + 6\tabcolsep\relax}|}{\begin{itemize}\item Research methods for human-computer interaction. Cambridge, UK.\end{itemize}}\\
\multicolumn{4}{|p{\dimexpr 16cm + 6\tabcolsep\relax}|}{\begin{itemize}\item Edited by Julie A. Jacko. The human-computer interaction handbook: fundamentals, evolving technologies, and emerging applications. 3rd ed. Boca Raton, FL: CRC Press, 2012.\end{itemize}}\\
\multicolumn{4}{|p{\dimexpr 16cm + 6\tabcolsep\relax}|}{\begin{itemize}\item Souza, Clarisse da Silva Vieira Camelo de. The semiotic engineering of human-computer interaction. Cambridge, MA: MIT Press, 2005.\end{itemize}}\\
\hline
\end{longtable}
\end{center}

\clearpage
\begin{center}
\begin{tabular}{lccr}
 \multirow{3}{*}{\includegraphics[height=2.7cm]{brasao.png}} &
 \multicolumn{2}{c}{\bfseries FEDERAL UNIVERSITY OF OURO PRETO} &
 \ \ \ \ \multirow{3}{*}{\includegraphics[height=2.7cm]{ufop.png}} \\
 & \multicolumn{2}{c}{\bfseries DEAN'S OFFICE FOR UNDERGRADUATE STUDIES} & \\
 & \multicolumn{2}{c}{} & \\
 & \multicolumn{2}{c}{\Large\bfseries COURSE PROGRAM} & \\
\end{tabular}
\end{center}

\begin{center}
\begin{longtable}{|p{4cm}|p{4cm}|p{4cm}|p{4cm}|}
\hline
\multicolumn{3}{|p{12cm}|}{Course Name in English:} &
\multicolumn{1}{p{4cm}|}{Code:} \\ 
\multicolumn{3}{|p{12cm}|}{\textbf{UI and UX for Artificial Intelligence}} &
\textbf{BCC931}\\ 
\multicolumn{3}{|p{12cm}|}{Course Name in Portuguese:} & \\ 
\multicolumn{3}{|p{12cm}|}{\textbf{UI e UX para Inteligência Artificial}} & \\ 
\hline
\multicolumn{3}{|p{12cm}|}{Department Name and Acronym:} & Academic Unit: \\ 
\multicolumn{3}{|p{12cm}|}{Department of Computer Science (DECOM)} & {ICEB} \\ 
\hline
\multicolumn{4}{|p{16cm}|}{Delivery Mode:
[X] in person \hspace{0.6cm}
[ ] blended \hspace{0.6cm}
[ ] distance learning}\\
\hline
\multicolumn{2}{|p{8cm}|}{Semester workload} &
\multicolumn{2}{p{8cm}|}{Weekly workload}\\
\hline
\multicolumn{1}{|p{4cm}|}{Total} &
\multicolumn{1}{p{4cm}|}{Extension} &
\multicolumn{1}{p{4cm}|}{Theoretical} &
\multicolumn{1}{p{4cm}|}{Practical} \\ 
\multicolumn{1}{|p{4cm}|}{60\,hours} &
\multicolumn{1}{p{4cm}|}{0\;hours} &
\multicolumn{1}{p{4cm}|}{4\;class hours} &
\multicolumn{1}{p{4cm}|}{0\;class hours} \\ 
\hline
\multicolumn{4}{|p{16cm}|}{Course Description:}\\
\multicolumn{4}{|p{16cm}|}{}\\
\multicolumn{4}{|p{\dimexpr 16cm + 6\tabcolsep\relax}|}{Study of user interface (UI) and user experience (UX) design principles applied to Artificial Intelligence systems. Focus on usability, accessibility, human-machine interaction, and the evaluation of intelligent interfaces.}\\
\multicolumn{4}{|p{16cm}|}{}\\
\hline
\multicolumn{4}{|p{16cm}|}{Course Content:}\\*
\multicolumn{4}{|p{\dimexpr 16cm + 6\tabcolsep\relax}|}{\begin{enumerate}[start=1]\item Fundamentals of UI and UX: concepts, history, and importance\end{enumerate}}\\
\multicolumn{4}{|p{\dimexpr 16cm + 6\tabcolsep\relax}|}{\begin{enumerate}[start=2]\item Principles of user-centered design\end{enumerate}}\\
\multicolumn{4}{|p{\dimexpr 16cm + 6\tabcolsep\relax}|}{\begin{enumerate}[start=3]\item UX specifics for AI-based systems: transparency, trust, and user control\end{enumerate}}\\
\multicolumn{4}{|p{\dimexpr 16cm + 6\tabcolsep\relax}|}{\begin{enumerate}[start=4]\item Prototyping and wireframing techniques for intelligent interfaces\end{enumerate}}\\
\multicolumn{4}{|p{\dimexpr 16cm + 6\tabcolsep\relax}|}{\begin{enumerate}[start=5]\item Human-AI interaction: conversation design, chatbots, and virtual assistants\end{enumerate}}\\
\multicolumn{4}{|p{\dimexpr 16cm + 6\tabcolsep\relax}|}{\begin{enumerate}[start=6]\item Usability evaluation: qualitative and quantitative methods\end{enumerate}}\\
\multicolumn{4}{|p{\dimexpr 16cm + 6\tabcolsep\relax}|}{\begin{enumerate}[start=7]\item Digital accessibility and inclusion in intelligent systems\end{enumerate}}\\
\multicolumn{4}{|p{\dimexpr 16cm + 6\tabcolsep\relax}|}{\begin{enumerate}[start=8]\item Ethical challenges in UI/UX for AI: bias, privacy, and explainability\end{enumerate}}\\
\multicolumn{4}{|p{\dimexpr 16cm + 6\tabcolsep\relax}|}{\begin{enumerate}[start=9]\item Tools and frameworks for UI/UX development in AI\end{enumerate}}\\
\multicolumn{4}{|p{\dimexpr 16cm + 6\tabcolsep\relax}|}{\begin{enumerate}[start=10]\item Case studies: real-world applications in AI and intelligent interfaces\end{enumerate}}\\
\multicolumn{4}{|p{\dimexpr 16cm + 6\tabcolsep\relax}|}{}\\
\hline
\multicolumn{4}{|p{16cm}|}{Basic Bibliography:}\\*
\multicolumn{4}{|p{\dimexpr 16cm + 6\tabcolsep\relax}|}{\begin{itemize}\item NORMAN, Donald A. The Design of Everyday Things. 2. ed. Basic Books, 2013.\end{itemize}}\\
\multicolumn{4}{|p{\dimexpr 16cm + 6\tabcolsep\relax}|}{\begin{itemize}\item COOPER, Alan; REIMANN, Robert; CRONIN, David. About Face: The Essentials of Interaction Design. 4. ed. Wiley, 2014.\end{itemize}}\\
\multicolumn{4}{|p{\dimexpr 16cm + 6\tabcolsep\relax}|}{\begin{itemize}\item SHNEIDERMAN, Ben; PLAISANT, Catherine; COHEN, Maxine S.; JACOBS, Steven; ELMQVIST, Niklas; DIAKOPOULOS, Nicholas. Designing the User Interface: Strategies for Effective Human-Computer Interaction. 6. ed. Boston: Pearson, 2016.\end{itemize}}\\
\multicolumn{4}{|p{16cm}|}{}\\
\hline
\multicolumn{4}{|p{16cm}|}{Complementary Bibliography:}\\*
\multicolumn{4}{|p{\dimexpr 16cm + 6\tabcolsep\relax}|}{\begin{itemize}\item KORE, Akshay. Designing Human-Centric AI Experiences: Applied UX Design for Artificial Intelligence. New York: Apress, 2022.\end{itemize}}\\
\multicolumn{4}{|p{\dimexpr 16cm + 6\tabcolsep\relax}|}{\begin{itemize}\item MACFADYEN, Louise. Designing AI Interfaces: Design Principles for Creative and Autonomous AI. Sebastopol, CA: O'Reilly Media, 2026.\end{itemize}}\\
\multicolumn{4}{|p{\dimexpr 16cm + 6\tabcolsep\relax}|}{\begin{itemize}\item SHEVAT, Amir. Designing Bots: Creating Conversational Experiences. Sebastopol, CA: O'Reilly Media, 2017.\end{itemize}}\\
\multicolumn{4}{|p{\dimexpr 16cm + 6\tabcolsep\relax}|}{\begin{itemize}\item LEVY, Jaime. UX Strategy: Product Strategy Techniques for Devising Innovative Digital Solutions. 2. ed. Sebastopol, CA: O'Reilly Media, 2021.\end{itemize}}\\
\hline
\end{longtable}
\end{center}

\clearpage
\begin{center}
\begin{tabular}{lccr}
 \multirow{3}{*}{\includegraphics[height=2.7cm]{brasao.png}} &
 \multicolumn{2}{c}{\bfseries FEDERAL UNIVERSITY OF OURO PRETO} &
 \ \ \ \ \multirow{3}{*}{\includegraphics[height=2.7cm]{ufop.png}} \\
 & \multicolumn{2}{c}{\bfseries DEAN'S OFFICE FOR UNDERGRADUATE STUDIES} & \\
 & \multicolumn{2}{c}{} & \\
 & \multicolumn{2}{c}{\Large\bfseries COURSE PROGRAM} & \\
\end{tabular}
\end{center}

\begin{center}
\begin{longtable}{|p{4cm}|p{4cm}|p{4cm}|p{4cm}|}
\hline
\multicolumn{3}{|p{12cm}|}{Course Name in English:} &
\multicolumn{1}{p{4cm}|}{Code:} \\ 
\multicolumn{3}{|p{12cm}|}{\textbf{Audio Processing}} &
\textbf{BCC932}\\ 
\multicolumn{3}{|p{12cm}|}{Course Name in Portuguese:} & \\ 
\multicolumn{3}{|p{12cm}|}{\textbf{Processamento de Áudio}} & \\ 
\hline
\multicolumn{3}{|p{12cm}|}{Department Name and Acronym:} & Academic Unit: \\ 
\multicolumn{3}{|p{12cm}|}{Department of Computer Science (DECOM)} & {ICEB} \\ 
\hline
\multicolumn{4}{|p{16cm}|}{Delivery Mode:
[X] in person \hspace{0.6cm}
[ ] blended \hspace{0.6cm}
[ ] distance learning}\\
\hline
\multicolumn{2}{|p{8cm}|}{Semester workload} &
\multicolumn{2}{p{8cm}|}{Weekly workload}\\
\hline
\multicolumn{1}{|p{4cm}|}{Total} &
\multicolumn{1}{p{4cm}|}{Extension} &
\multicolumn{1}{p{4cm}|}{Theoretical} &
\multicolumn{1}{p{4cm}|}{Practical} \\ 
\multicolumn{1}{|p{4cm}|}{60\,hours} &
\multicolumn{1}{p{4cm}|}{0\;hours} &
\multicolumn{1}{p{4cm}|}{4\;class hours} &
\multicolumn{1}{p{4cm}|}{0\;class hours} \\ 
\hline
\multicolumn{4}{|p{16cm}|}{Course Description:}\\
\multicolumn{4}{|p{16cm}|}{}\\
\multicolumn{4}{|p{\dimexpr 16cm + 6\tabcolsep\relax}|}{Introduction to fundamental concepts and techniques for digital audio signal processing. Analysis, transformation, filtering, and synthesis of audio applied to recognition, sound synthesis, and AI systems.}\\
\multicolumn{4}{|p{16cm}|}{}\\
\hline
\multicolumn{4}{|p{16cm}|}{Course Content:}\\*
\multicolumn{4}{|p{\dimexpr 16cm + 6\tabcolsep\relax}|}{\begin{enumerate}[start=1]\item Traditional acoustic models, Gaussian mixtures, and hidden Markov chains\end{enumerate}}\\
\multicolumn{4}{|p{\dimexpr 16cm + 6\tabcolsep\relax}|}{\begin{enumerate}[start=2]\item Signal representations for audio and speech\end{enumerate}}\\
\multicolumn{4}{|p{\dimexpr 16cm + 6\tabcolsep\relax}|}{\begin{enumerate}[start=3]\item Neural network architectures for speech recognition\end{enumerate}}\\
\multicolumn{4}{|p{\dimexpr 16cm + 6\tabcolsep\relax}|}{\begin{enumerate}[start=4]\item Encoder, vocoder models, and architectures for speech synthesis\end{enumerate}}\\
\multicolumn{4}{|p{\dimexpr 16cm + 6\tabcolsep\relax}|}{\begin{enumerate}[start=5]\item Architectures for music applications\end{enumerate}}\\
\multicolumn{4}{|p{\dimexpr 16cm + 6\tabcolsep\relax}|}{}\\
\hline
\multicolumn{4}{|p{16cm}|}{Basic Bibliography:}\\*
\multicolumn{4}{|p{\dimexpr 16cm + 6\tabcolsep\relax}|}{\begin{itemize}\item KAMATH, U. L.; WHITAKER, J. Deep learning for NLP and speech recognition. Springer Nature, 2019.\end{itemize}}\\
\multicolumn{4}{|p{\dimexpr 16cm + 6\tabcolsep\relax}|}{\begin{itemize}\item DONG, Y.; DENG, Y. Automatic speech recognition. Springer London Limited, 2016.\end{itemize}}\\
\multicolumn{4}{|p{\dimexpr 16cm + 6\tabcolsep\relax}|}{\begin{itemize}\item GOPI, E. S. Digital speech processing using Matlab. New Delhi: Springer, 2014.\end{itemize}}\\
\multicolumn{4}{|p{16cm}|}{}\\
\hline
\multicolumn{4}{|p{16cm}|}{Complementary Bibliography:}\\*
\multicolumn{4}{|p{\dimexpr 16cm + 6\tabcolsep\relax}|}{\begin{itemize}\item CAMASTRA, F.; VINCIARELLI, A. Machine learning for audio, image and video analysis: theory and applications. Springer, 2015.\end{itemize}}\\
\multicolumn{4}{|p{\dimexpr 16cm + 6\tabcolsep\relax}|}{\begin{itemize}\item STEVENS, E.; ANTIGA, L. Deep Learning with PyTorch. Manning Publications, 2020.\end{itemize}}\\
\multicolumn{4}{|p{\dimexpr 16cm + 6\tabcolsep\relax}|}{\begin{itemize}\item GALEONE, P. Hands-On Neural Networks with TensorFlow 2.0: Understand TensorFlow, from static graph to eager execution, and design neural networks. Packt Publishing, 2019.\end{itemize}}\\
\hline
\end{longtable}
\end{center}

\clearpage
\begin{center}
\begin{tabular}{lccr}
 \multirow{3}{*}{\includegraphics[height=2.7cm]{brasao.png}} &
 \multicolumn{2}{c}{\bfseries FEDERAL UNIVERSITY OF OURO PRETO} &
 \ \ \ \ \multirow{3}{*}{\includegraphics[height=2.7cm]{ufop.png}} \\
 & \multicolumn{2}{c}{\bfseries DEAN'S OFFICE FOR UNDERGRADUATE STUDIES} & \\
 & \multicolumn{2}{c}{} & \\
 & \multicolumn{2}{c}{\Large\bfseries COURSE PROGRAM} & \\
\end{tabular}
\end{center}

\begin{center}
\begin{longtable}{|p{4cm}|p{4cm}|p{4cm}|p{4cm}|}
\hline
\multicolumn{3}{|p{12cm}|}{Course Name in English:} &
\multicolumn{1}{p{4cm}|}{Code:} \\ 
\multicolumn{3}{|p{12cm}|}{\textbf{Intelligent Robotics}} &
\textbf{BCC306}\\ 
\multicolumn{3}{|p{12cm}|}{Course Name in Portuguese:} & \\ 
\multicolumn{3}{|p{12cm}|}{\textbf{Robótica Inteligente}} & \\ 
\hline
\multicolumn{3}{|p{12cm}|}{Department Name and Acronym:} & Academic Unit: \\ 
\multicolumn{3}{|p{12cm}|}{Department of Computer Science (DECOM)} & {ICEB} \\ 
\hline
\multicolumn{4}{|p{16cm}|}{Delivery Mode:
[X] in person \hspace{0.6cm}
[ ] blended \hspace{0.6cm}
[ ] distance learning}\\
\hline
\multicolumn{2}{|p{8cm}|}{Semester workload} &
\multicolumn{2}{p{8cm}|}{Weekly workload}\\
\hline
\multicolumn{1}{|p{4cm}|}{Total} &
\multicolumn{1}{p{4cm}|}{Extension} &
\multicolumn{1}{p{4cm}|}{Theoretical} &
\multicolumn{1}{p{4cm}|}{Practical} \\ 
\multicolumn{1}{|p{4cm}|}{60\,hours} &
\multicolumn{1}{p{4cm}|}{0\;hours} &
\multicolumn{1}{p{4cm}|}{4\;class hours} &
\multicolumn{1}{p{4cm}|}{0\;class hours} \\ 
\hline
\multicolumn{4}{|p{16cm}|}{Course Description:}\\
\multicolumn{4}{|p{16cm}|}{}\\
\multicolumn{4}{|p{\dimexpr 16cm + 6\tabcolsep\relax}|}{Introduction to intelligent robotics; sensors, actuators, and embedded systems in robotics; computer vision in robotics; robot control; intelligent techniques for route planning, navigation, and localization of mobile robots; machine learning in robotics; social robotics and human-robot interaction.}\\
\multicolumn{4}{|p{16cm}|}{}\\
\hline
\multicolumn{4}{|p{16cm}|}{Course Content:}\\*
\multicolumn{4}{|p{\dimexpr 16cm + 6\tabcolsep\relax}|}{\begin{enumerate}[start=1]\item Introduction to intelligent robotics\end{enumerate}}\\
\multicolumn{4}{|p{\dimexpr 16cm + 6\tabcolsep\relax}|}{\begin{enumerate}[start=2]\item Sensors, actuators, and embedded systems in robotics\end{enumerate}}\\
\multicolumn{4}{|p{\dimexpr 16cm + 6\tabcolsep\relax}|}{\begin{enumerate}[start=3]\item Computer vision in robotics\end{enumerate}}\\
\multicolumn{4}{|p{\dimexpr 16cm + 6\tabcolsep\relax}|}{\begin{enumerate}[start=4]\item Robot control\end{enumerate}}\\
\multicolumn{4}{|p{\dimexpr 16cm + 6\tabcolsep\relax}|}{\begin{enumerate}[start=5]\item Intelligent techniques for route planning, navigation, and localization of mobile robots\end{enumerate}}\\
\multicolumn{4}{|p{\dimexpr 16cm + 6\tabcolsep\relax}|}{\begin{enumerate}[start=6]\item Machine learning in robotics\end{enumerate}}\\
\multicolumn{4}{|p{\dimexpr 16cm + 6\tabcolsep\relax}|}{\begin{enumerate}[start=7]\item Social robotics and human-robot interaction\end{enumerate}}\\
\multicolumn{4}{|p{\dimexpr 16cm + 6\tabcolsep\relax}|}{}\\
\hline
\multicolumn{4}{|p{16cm}|}{Basic Bibliography:}\\*
\multicolumn{4}{|p{\dimexpr 16cm + 6\tabcolsep\relax}|}{\begin{itemize}\item Craig, John J. Robótica. 3. ed. São Paulo: Pearson Education do Brasil, 2013.\end{itemize}}\\
\multicolumn{4}{|p{\dimexpr 16cm + 6\tabcolsep\relax}|}{\begin{itemize}\item Spong, Mark W.; Vidyasagar, M. Robot dynamics and control. New York: J. Wiley, 1989.\end{itemize}}\\
\multicolumn{4}{|p{\dimexpr 16cm + 6\tabcolsep\relax}|}{\begin{itemize}\item Russell, S. J.; Norvig, P. Inteligência Artificial: uma abordagem moderna. 4. ed. Rio de Janeiro: LTC, 2022.\end{itemize}}\\
\multicolumn{4}{|p{16cm}|}{}\\
\hline
\multicolumn{4}{|p{16cm}|}{Complementary Bibliography:}\\*
\multicolumn{4}{|p{\dimexpr 16cm + 6\tabcolsep\relax}|}{\begin{itemize}\item Romero, R. A. F.; Prestes, E. S. J.; Osório, F. S.; Wolf, D. F. (org.). Robótica móvel. 1. ed. Rio de Janeiro: LTC, 2014.\end{itemize}}\\
\multicolumn{4}{|p{\dimexpr 16cm + 6\tabcolsep\relax}|}{\begin{itemize}\item Romero, R. A. F. (org.). Robótica educacional. 1. ed. São Paulo: Edusp, 2024.\end{itemize}}\\
\multicolumn{4}{|p{\dimexpr 16cm + 6\tabcolsep\relax}|}{\begin{itemize}\item Sutton, R. S.; Barto, A. G. Reinforcement learning: an introduction. 2. ed. Cambridge, MA: MIT Press, 2018.\end{itemize}}\\
\multicolumn{4}{|p{\dimexpr 16cm + 6\tabcolsep\relax}|}{\begin{itemize}\item Elgendy, M. Deep learning for vision systems. 1. ed. Manning, 2020.\end{itemize}}\\
\multicolumn{4}{|p{\dimexpr 16cm + 6\tabcolsep\relax}|}{\begin{itemize}\item Siegwart, R.; Nourbakhsh, I. R.; Scaramuzza, D. Introduction to autonomous mobile robots. 2. ed. Cambridge, MA: MIT Press, 2011.\end{itemize}}\\
\hline
\end{longtable}
\end{center}

\clearpage
\begin{center}
\begin{tabular}{lccr}
 \multirow{3}{*}{\includegraphics[height=2.7cm]{brasao.png}} &
 \multicolumn{2}{c}{\bfseries FEDERAL UNIVERSITY OF OURO PRETO} &
 \ \ \ \ \multirow{3}{*}{\includegraphics[height=2.7cm]{ufop.png}} \\
 & \multicolumn{2}{c}{\bfseries DEAN'S OFFICE FOR UNDERGRADUATE STUDIES} & \\
 & \multicolumn{2}{c}{} & \\
 & \multicolumn{2}{c}{\Large\bfseries COURSE PROGRAM} & \\
\end{tabular}
\end{center}

\begin{center}
\begin{longtable}{|p{4cm}|p{4cm}|p{4cm}|p{4cm}|}
\hline
\multicolumn{3}{|p{12cm}|}{Course Name in English:} &
\multicolumn{1}{p{4cm}|}{Code:} \\ 
\multicolumn{3}{|p{12cm}|}{\textbf{Internet of Things}} &
\textbf{BCC883}\\ 
\multicolumn{3}{|p{12cm}|}{Course Name in Portuguese:} & \\ 
\multicolumn{3}{|p{12cm}|}{\textbf{Internet das Coisas}} & \\ 
\hline
\multicolumn{3}{|p{12cm}|}{Department Name and Acronym:} & Academic Unit: \\ 
\multicolumn{3}{|p{12cm}|}{Department of Computer Science (DECOM)} & {ICEB} \\ 
\hline
\multicolumn{4}{|p{16cm}|}{Delivery Mode:
[X] in person \hspace{0.6cm}
[ ] blended \hspace{0.6cm}
[ ] distance learning}\\
\hline
\multicolumn{2}{|p{8cm}|}{Semester workload} &
\multicolumn{2}{p{8cm}|}{Weekly workload}\\
\hline
\multicolumn{1}{|p{4cm}|}{Total} &
\multicolumn{1}{p{4cm}|}{Extension} &
\multicolumn{1}{p{4cm}|}{Theoretical} &
\multicolumn{1}{p{4cm}|}{Practical} \\ 
\multicolumn{1}{|p{4cm}|}{60\,hours} &
\multicolumn{1}{p{4cm}|}{0\;hours} &
\multicolumn{1}{p{4cm}|}{4\;class hours} &
\multicolumn{1}{p{4cm}|}{0\;class hours} \\ 
\hline
\multicolumn{4}{|p{16cm}|}{Course Description:}\\
\multicolumn{4}{|p{16cm}|}{}\\
\multicolumn{4}{|p{\dimexpr 16cm + 6\tabcolsep\relax}|}{Introduction to the concepts, architectures, and technologies of the Internet of Things (IoT). Communication protocols, sensors, actuators, security, and intelligent applications in connected environments.}\\
\multicolumn{4}{|p{16cm}|}{}\\
\hline
\multicolumn{4}{|p{16cm}|}{Course Content:}\\*
\multicolumn{4}{|p{\dimexpr 16cm + 6\tabcolsep\relax}|}{\begin{enumerate}[start=1]\item Fundamentals of the Internet of Things: definitions, evolution, and current landscape\end{enumerate}}\\
\multicolumn{4}{|p{\dimexpr 16cm + 6\tabcolsep\relax}|}{\begin{enumerate}[start=2]\item IoT architecture: layers, devices, and embedded systems\end{enumerate}}\\
\multicolumn{4}{|p{\dimexpr 16cm + 6\tabcolsep\relax}|}{\begin{enumerate}[start=3]\item Sensors and actuators: types, operation, and integration\end{enumerate}}\\
\multicolumn{4}{|p{\dimexpr 16cm + 6\tabcolsep\relax}|}{\begin{enumerate}[start=4]\item IoT communication protocols\end{enumerate}}\\
\multicolumn{4}{|p{\dimexpr 16cm + 6\tabcolsep\relax}|}{\begin{enumerate}[start=5]\item Platforms and tools for IoT development\end{enumerate}}\\
\multicolumn{4}{|p{\dimexpr 16cm + 6\tabcolsep\relax}|}{\begin{enumerate}[start=6]\item Security and privacy in IoT environments\end{enumerate}}\\
\multicolumn{4}{|p{\dimexpr 16cm + 6\tabcolsep\relax}|}{\begin{enumerate}[start=7]\item Data collection, transmission, and processing in IoT\end{enumerate}}\\
\multicolumn{4}{|p{\dimexpr 16cm + 6\tabcolsep\relax}|}{\begin{enumerate}[start=8]\item Integration with artificial intelligence and cloud computing\end{enumerate}}\\
\multicolumn{4}{|p{\dimexpr 16cm + 6\tabcolsep\relax}|}{\begin{enumerate}[start=9]\item Use cases: smart cities, precision agriculture, healthcare, and Industry 4.0\end{enumerate}}\\
\multicolumn{4}{|p{\dimexpr 16cm + 6\tabcolsep\relax}|}{\begin{enumerate}[start=10]\item Future trends and challenges of the Internet of Things\end{enumerate}}\\
\multicolumn{4}{|p{\dimexpr 16cm + 6\tabcolsep\relax}|}{}\\
\hline
\multicolumn{4}{|p{16cm}|}{Basic Bibliography:}\\*
\multicolumn{4}{|p{\dimexpr 16cm + 6\tabcolsep\relax}|}{\begin{itemize}\item HANES, David; SALGUEIRO, Gonzalo; GROSSETETE, Patrick; BARTON, Robert; HENRY, Jerome. IoT Fundamentals: Networking Technologies, Protocols, and Use Cases for the Internet of Things. Indianapolis: Cisco Press, 2017.\end{itemize}}\\
\multicolumn{4}{|p{\dimexpr 16cm + 6\tabcolsep\relax}|}{\begin{itemize}\item BUYYA, Rajkumar; DASTJERDI, Amir Vahid (ed.). Internet of Things: Principles and Paradigms. Cambridge, MA: Morgan Kaufmann, 2016.\end{itemize}}\\
\multicolumn{4}{|p{\dimexpr 16cm + 6\tabcolsep\relax}|}{\begin{itemize}\item MINERVA, Roberto et al. Towards a definition of the Internet of Things (IoT). IEEE Internet Initiative, 2015.\end{itemize}}\\
\multicolumn{4}{|p{16cm}|}{}\\
\hline
\multicolumn{4}{|p{16cm}|}{Complementary Bibliography:}\\*
\multicolumn{4}{|p{\dimexpr 16cm + 6\tabcolsep\relax}|}{\begin{itemize}\item GUBBI, Jayavardhana et al. Internet of Things (IoT): A Vision, Architectural Elements, and Future Directions. Future Generation Computer Systems, v. 29, n. 7, p. 1645-1660, 2013.\end{itemize}}\\
\multicolumn{4}{|p{\dimexpr 16cm + 6\tabcolsep\relax}|}{\begin{itemize}\item STANKOVIC, John A. Research Directions for the Internet of Things. IEEE Internet of Things Journal, v. 1, n. 1, p. 3-9, 2014.\end{itemize}}\\
\multicolumn{4}{|p{\dimexpr 16cm + 6\tabcolsep\relax}|}{\begin{itemize}\item CIRANI, Simone; FERRARI, Gianluigi; PICONE, Marco; VELTRI, Luca. Internet of Things: Architectures, Protocols and Standards. Hoboken, NJ: Wiley, 2019.\end{itemize}}\\
\multicolumn{4}{|p{\dimexpr 16cm + 6\tabcolsep\relax}|}{\begin{itemize}\item SALMAN, Tara; JAIN, Raj. A survey of protocols and standards for Internet of Things. arXiv preprint, arXiv:1903.11549, 2019.\end{itemize}}\\
\multicolumn{4}{|p{\dimexpr 16cm + 6\tabcolsep\relax}|}{\begin{itemize}\item VERMESAN, Ovidiu; FRIESS, Peter (ed.). Internet of Things Applications: From Research and Innovation to Market Deployment. River Publishers, 2014.\end{itemize}}\\
\hline
\end{longtable}
\end{center}

\clearpage
\end{document}